%% =========================================================
%% Chapter: Physics examples of the multi-block asymmetric
%% compression theorem, second of three chapters: Kramers--Wannier
%% duality and the Fibonacci and Ising fusion operators. The
%% Fibonacci anomaly file is input right after the Fibonacci file
%% and supplies the closing \subsection of that section.
%% =========================================================

\chapter{Fusion of non-invertible symmetry operators}%
\label{ch:asymmetric_examples_non_invertible}

This chapter continues the physical examples of the multi-block asymmetric
compression theorem (Theorem~\ref{thm:asym_multiblock}) begun in
Chapter~\ref{ch:asymmetric_examples} with symmetry operators that have no
inverse: Kramers--Wannier
duality~\cite{KramersWannier1941,Aasen2016Topological,Seiberg2024Majorana,Lootens2023Dualities},
and the Fibonacci~\cite[Appendix~D.1]{Bultinck2017Anyons} and
Ising~\cite[Appendix~D.2]{Bultinck2017Anyons} fusion categories and their
anyon chains~\cite{Feiguin2007GoldenChain}, together with the review's
Fibonacci operator~\cite[Appendix~A, ``The MPO for the Fibonacci model'']{Cirac2021Matrix}.
Building on the general
definitions of Chapter~\ref{ch:mpo_symmetry_basics} (fusion rules, units,
invertible and non-invertible labels, Perron--Frobenius dimensions,
symmetric matrix product states with their nonnegative integer
representations, and the zipper condition), each section constructs
the explicit local tensors, evaluates their fusion tables, and computes
their periodic compressions and anomaly invariants. The next chapter treats
group symmetries in the same way.

\section{Kramers--Wannier duality}%
\label{sec:asymex_kw}

Kramers--Wannier duality originates in the two-dimensional classical
Ising model, where it relates the high- and low-temperature
regimes~\cite{KramersWannier1941}; in the quantum transverse-field
Ising chain it exchanges the bond and field terms of the Hamiltonian.
On a periodic chain the duality is a noninvertible operator, rather than
a unitary change of variables on the full Hilbert space. Aasen, Mong and
Fendley realize it as a lattice defect~\cite{Aasen2016Topological};
Seiberg and Shao place it in a noninvertible symmetry algebra together
with translation and the spin flip~\cite{Seiberg2024Majorana}; Lootens,
Delcamp, Ortiz and Verstraete write it as a matrix product operator
intertwiner between the two Hamiltonians~\cite{Lootens2023Dualities}.
The review discusses this operator among its matrix product operator
symmetries~\cite[Section~III.A.6]{Cirac2021Matrix}; its relation to fusion
algebras is Remark~\ref{rem:mposym_kramers_wannier}. Here the duality is
built from an explicit two-by-two tensor whose periodic operator is
unnormalized; it reproduces the primal-to-dual kernel
of~\cite{Aasen2016Topological} up to the factor $2^{N/2}$. Its square
and its action on ground states are characterized both as operator identities
and as asymmetric compressions of stacked tensors.

\begin{definition}[The duality kernel and the two translation tensors]
    \label{def:asymex_kw}
    \lean{KWExample.kwTensor, KWExample.shiftTensor, KWExample.flipShiftTensor,
        KWExample.kwSquare, KWExample.kwSquareTargets}
    \leanok
    \uses{def:mpo_tensor, def:mpo_to_mps, def:mpdo_operator_product}
    With the outgoing label $s'$ written first, the duality tensor is
    \begin{align}
        W^{00} & =\begin{pmatrix}1&1\\0&0\end{pmatrix},
               &
        W^{01} & =\begin{pmatrix}0&0\\1&1\end{pmatrix},
        \notag                                           \\
        W^{10} & =\begin{pmatrix}1&-1\\0&0\end{pmatrix},
               &
        W^{11} & =\begin{pmatrix}0&0\\-1&1\end{pmatrix}.
        \notag
    \end{align}
    The translation tensor is $T^{s's}=E_{s,s'}$ and the flipped
    translation tensor is $(\eta T)^{s's}=E_{s,1-s'}$, where $E_{pq}$ is
    the matrix unit with a single entry $1$ in row $p$ and column $q$.
    The same symbols $T$ and $\eta T$ denote their periodic operators.
    The stacked product tensor is
    $B^{s's}=\sum_{m=0}^1W^{s'm}\otimes W^{ms}$, of bond dimension four.
    Its two target tensors are $2T$ and $2\eta T$.
\end{definition}

Close the virtual bonds by the ordinary trace on a ring of $N>0$ sites,
without dividing by the bond dimension. The resulting operator is
$K_N(a,b)=\tr(W^{a_0b_0}\cdots W^{a_{N-1}b_{N-1}})$.
Its normalization relative to the source defect follows from
Theorem~\ref{thm:asymex_kw_kernel} and Corollary~\ref{cor:asymex_kw_square}.
Theorem~\ref{thm:asymex_kw_bimodule_absorption} proves that $W$ is normal,
and Theorem~\ref{thm:asymex_kw_normal} that $T$ and $\eta T$ are;
Corollary~\ref{cor:asymex_kw_exchange} gives the exchange of Ising couplings,
and Lemma~\ref{lem:asymex_kw_states} describes its action on product and
GHZ states. The compression of its square is
Theorem~\ref{thm:asymex_kw_compression}.

Consider $N\geq1$ qubits on a periodic chain, with sites
$I_N=\{0,\ldots,N-1\}$ and all site indices understood modulo $N$.
The Hilbert space $\mathcal H_N=(\C^2)^{\otimes N}$ has computational
basis $\{\ket{b}:b\in\{0,1\}^{I_N}\}$. Write $\bar b$ for the
componentwise complement of $b$, and $b\oplus e_j$ for the configuration
obtained by flipping its bit at site $j$. The Pauli operators act by
\begin{align}
    X_j\ket{b} & =\ket{b\oplus e_j},
    \notag                           \\
    Z_j\ket{b} & =(-1)^{b_j}\ket{b}.
    \label{eq:asymex_kw_pauli}
\end{align}
They are self-adjoint involutions, satisfy $X_jZ_j=-Z_jX_j$, and commute
at distinct sites. Define global spin flip $\eta$ and one-site left
translation $T$ by
\begin{align}
    \eta\ket{b} & =\ket{\bar b},
    \notag                       \\
    T\ket{b}    & =\ket{\tau b},
    \qquad (\tau b)_j=b_{j+1}.
    \label{eq:asymex_kw_translation}
\end{align}
Thus $\eta=\prod_jX_j$, $\eta^2=\Id$, $T^N=\Id$,
$T^\dagger=T^{-1}$, and $T\eta=\eta T$. In this convention,
$TX_jT^{-1}=X_{j-1}$ and $TZ_jT^{-1}=Z_{j-1}$.

For real couplings $J,h$, the periodic transverse-field Ising Hamiltonian is
\begin{align}
    H_N(J,h) & =-J\sum_{j\in I_N}Z_jZ_{j+1}-h\sum_{j\in I_N}X_j.
    \label{eq:asymex_kw_ising}
\end{align}
The sum includes one bond per site. In particular, the bond at $N=1$
is $Z_0^2=\Id$, and at $N=2$ the two bond terms both equal $Z_0Z_1$.
The duality exchanges the bond terms with the field terms.

\subsection{The duality kernel and its square}

The translation tensor is the shift matrix product unitary
of~\cite[Appendix~A, ``The shift MPU'']{Cirac2021Matrix}, taken with
the opposite direction of translation, and the flipped translation
tensor composes it with the global spin flip.

\begin{center}
    % Formula: B^{s's}=\sum_mW^{s'm}\otimes W^{ms}, the stacked product
    % of Definition~\ref{def:asymex_kw}.
    % Source: Definition~\ref{def:asymex_kw}.
    % Ink-to-index map: both rows are the duality tensor $W$; the shared
    % vertical pairleg at each site is the summed index $m$; the drawn
    % upper-row label $r$ is the per-site component of the outgoing
    % index $s'$.
    % Contraction set: the vertical pairlegs (index $m$) at every site
    % and the horizontal virtual bonds within each row; the physical
    % legs $(r,s)$, the per-site components of $(s',s)$, remain open.
    % Boundary signature: (virtual-west=0 | virtual-east=0 |
    % physical-up=3 | physical-down=3) on both sides.
    \tenkzkernel
    \begin{tenkz}[rows={op,op}, cols=4, boundary=periodic]
        \tn[skin=mpo, ports={90:physical:$r_1$, 270:physical}]{W} &
        \tn[skin=mpo, ports={90:physical:$r_2$, 270:physical}]{W} &
        \tn[skin=dots, ports={180:virtual, 0:virtual}]{} &
        \tn[skin=mpo, ports={90:physical:$r_N$, 270:physical}]{W} \\
        \tn[skin=mpo, ports={90:physical, 270:physical:$s_1$}]{W} &
        \tn[skin=mpo, ports={90:physical, 270:physical:$s_2$}]{W} &
        \tn[skin=dots, ports={180:virtual, 0:virtual}]{} &
        \tn[skin=mpo, ports={90:physical, 270:physical:$s_N$}]{W}
    \end{tenkz}
    \quad $=$ \quad
    \begin{tenkz}[rows={op}, cols=4, boundary=periodic]
        \tn[skin=mpo, ports={90:physical:$r_1$, 270:physical:$s_1$}]{B} &
        \tn[skin=mpo, ports={90:physical:$r_2$, 270:physical:$s_2$}]{B} &
        \tn[skin=dots, ports={180:virtual, 0:virtual}]{} &
        \tn[skin=mpo, ports={90:physical:$r_N$, 270:physical:$s_N$}]{B}
    \end{tenkz}
\end{center}

The site labels in the diagrams run from $1$ to $N$, a cyclic relabeling
of $I_N$.

\subsubsection{Operator identities}

Define the unnormalized periodic operator $K_N$ by the ordinary virtual trace,
without dividing by the bond dimension:
\begin{align}
    (K_N)_{ab}
     & =\tr(W^{a_0b_0}\cdots W^{a_{N-1}b_{N-1}}).
    \label{eq:asymex_kw_periodic}
\end{align}
We also write $D=K_N$ for this unnormalized operator, retaining $D$ in the
square identity below. The normalized operators will be denoted
$\mathcal D_N$ and $U_N$, not $D$.

\begin{theorem}[Periodic kernel]
    \label{thm:asymex_kw_kernel}
    \lean{KWExample.kwTensor_apply,
        KWExample.kwTensor_mpo_eq_prod,
        KWExample.kwTensor_mpo_eq_pow_sum}
    \leanok
    \uses{def:asymex_kw}

    For all configurations $s',s$,
    \begin{align}
        \bra{s'}D\ket{s}
         & =\prod_{j\in I_N}(-1)^{s'_j(s_j+s_{j+1})}.
        \label{eq:asymex_kw_kernel}
    \end{align}
\end{theorem}

\begin{proof}
    \leanok
    The local entries are $(W^{ab})_{uv}=\delta_{u,b}(-1)^{a(b+v)}$.
    Expanding the trace over closed virtual configurations $g$ gives
    \begin{align}
        (K_N)_{ab}
         & =\sum_{g\in\{0,1\}^{I_N}}
        \prod_{j\in I_N}\delta_{g_j,b_j}
        (-1)^{a_j(b_j+g_{j+1})}
        =\prod_{j\in I_N}(-1)^{a_j(b_j+b_{j+1})}.
        \notag
    \end{align}
    Only the configuration $g=b$ contributes.
\end{proof}

For subsequent calculations, put $q(a,b)=\sum_j a_j(b_j+b_{j+1})$,
so $(K_N)_{ab}=(-1)^{q(a,b)}$. Exponents may be calculated modulo two.
Unless indicated otherwise, site sums and products run over $I_N$.

\begin{theorem}[Spin-flip absorption and noninvertibility]
    \label{thm:asymex_kw_absorption}
    \lean{KWExample.kwTensor_mpo_flip_input,
        KWExample.kwTensor_mpo_flip_output,
        KWExample.kwTensor_mpo_mul_spinFlip,
        KWExample.spinFlip_mul_kwTensor_mpo,
        KWExample.kwTensor_mpo_mulVec_eq_zero_of_odd,
        KWExample.kwTensor_mpo_not_isUnit}
    \leanok
    \uses{thm:asymex_kw_kernel}
    For every $N>0$, the unnormalized operator satisfies
    $K_N\eta=\eta K_N=K_N$. If $\eta v=-v$, then $K_Nv=0$.
    In particular, $K_N$ is not invertible.
\end{theorem}
\begin{proof}
    \leanok
    Complementing either input or output leaves the kernel unchanged:
    $q(a,\bar b)\equiv q(\bar a,b)\equiv q(a,b)\pmod2$.
    This proves both absorption identities. For an odd vector,
    $K_Nv=K_N\eta v=-K_Nv$, so $K_Nv=0$.
    The nonzero vector $\ket{0}^{\otimes N}-\ket{1}^{\otimes N}$
    has odd spin flip, so $K_N$ has a nontrivial kernel.
\end{proof}

\begin{lemma}[Local duality relations]
    \label{lem:asymex_kw_local}
    \lean{KWExample.kwTensor_mpo_mul_siteX, KWExample.kwTensor_mpo_mul_siteZZ,
        KWExample.siteX, KWExample.siteZ, KWExample.siteZZ}
    \leanok
    \uses{thm:asymex_kw_kernel}
    For every site $j$,
    \begin{align}
        K_NX_j        & =Z_{j-1}Z_jK_N,
        \label{eq:asymex_kw_local_x}    \\
        K_NZ_jZ_{j+1} & =X_jK_N.
        \label{eq:asymex_kw_local_zz}
    \end{align}
\end{lemma}

\begin{proof}
    \leanok
    Flipping one input or output bit gives
    \begin{align}
        q(a,b\oplus e_j) & \equiv q(a,b)+a_{j-1}+a_j\pmod2,
        \notag                                              \\
        q(a\oplus e_j,b) & \equiv q(a,b)+b_j+b_{j+1}\pmod2.
        \notag
    \end{align}
    These congruences prove~(\ref{eq:asymex_kw_local_x})
    and~(\ref{eq:asymex_kw_local_zz}) entrywise.
    They remain valid when the two indicated sites coincide.
\end{proof}

Under the identification of the dual lattice described after
Corollary~\ref{cor:asymex_kw_square}, these are the two local relations
of the duality defect in~\cite{Aasen2016Topological}.

\begin{corollary}[Coupling exchange]
    \label{cor:asymex_kw_exchange}
    \lean{KWExample.isingHamiltonian,
        KWExample.kwTensor_mpo_mul_isingHamiltonian,
        KWExample.kwTensor_mpo_mul_isingHamiltonian_self}
    \leanok
    \uses{lem:asymex_kw_local}
    The Hamiltonians with exchanged couplings satisfy
    \begin{align}
        K_NH_N(J,h) & =H_N(h,J)K_N.
        \label{eq:asymex_kw_exchange}
    \end{align}
    In particular, $K_N$ commutes with $H_N(J,J)$.
\end{corollary}

\begin{proof}
    \leanok
    Summing the local relations of
    Lemma~\ref{lem:asymex_kw_local} and reindexing the bonds yields
    \begin{align}
        K_NH_N(J,h)
         & =\left(-J\sum_jX_j-h\sum_jZ_{j-1}Z_j\right)K_N
        =H_N(h,J)K_N.
        \notag
    \end{align}
\end{proof}

The local identities also exchange strings and endpoint operators.
For consecutive sites $j_r=j_0+r$ and $1\leq m\leq N$,
\begin{align}
    K_N\prod_{r=0}^{m-1}X_{j_r}
     & =Z_{j_0-1}Z_{j_{m-1}}K_N,
    \notag                                        \\
    K_NZ_{j_0}Z_{j_m}
     & =\left(\prod_{r=0}^{m-1}X_{j_r}\right)K_N.
    \label{eq:asymex_kw_strings}
\end{align}
Indeed, iterating~(\ref{eq:asymex_kw_local_x}) and~(\ref{eq:asymex_kw_local_zz})
cancels every interior
factor $Z_j^2$. At $m=N$ the endpoints coincide, and these relations
give $K_N\eta=\eta K_N=K_N$. They are intertwining identities; no
inverse of $K_N$ is involved.

\begin{lemma}[Adjoint and Gram matrices]
    \label{lem:asymex_kw_gram}
    \lean{KWExample.translate_mul_kwTensor_mpo,
        KWExample.kwTensor_mpo_conjTranspose_eq_inv_translate_mul,
        KWExample.kwTensor_mpo_conjTranspose_mul,
        KWExample.kwTensor_mpo_mul_conjTranspose, KWExample.translate,
        KWExample.shiftTensor_mpo_eq_translate, KWExample.translate_inv,
        KWExample.kwTensor_mpo_conjTranspose, KWExample.kwTensor_mpo_transpose,
        KWExample.kwTensor_mpo_transpose_eq_transpose_translate_mul,
        KWExample.kwTensor_mpo_mul_transpose}
    \leanok
    \uses{thm:asymex_kw_kernel, def:asymex_kw}
    The unnormalized operator satisfies
    \begin{align}
        TK_N        & =K_NT,
        \notag                    \\
        K_N^\dagger & =T^{-1}K_N,
        \label{eq:asymex_kw_adjoint}
    \end{align}
    and
    \begin{align}
        K_N^\dagger K_N=K_NK_N^\dagger & =2^N(\Id+\eta).
        \label{eq:asymex_kw_gram}
    \end{align}
\end{lemma}

\begin{proof}
    \leanok
    Cyclic reindexing gives $q(\tau a,\tau b)=q(a,b)$ and
    $q(b,a)=q(\tau a,b)$. Since the kernel is real, these identities
    imply~(\ref{eq:asymex_kw_adjoint}).

    For a binary configuration $c$, the character sum factors as
    \begin{align}
        \sum_a(-1)^{\sum_ja_jc_j}
         & =\prod_j(1+(-1)^{c_j})
        =\begin{cases}
             2^N, & c=\mathbf0,    \\
             0,   & c\neq\mathbf0.
         \end{cases}
        \label{eq:asymex_kw_character}
    \end{align}
    Here $\mathbf0$ and $\mathbf1$ denote the two constant configurations.
    For input configurations $b,c$, put $d=b\oplus c$. Then
    \begin{align}
        (K_N^\dagger K_N)_{bc}
         & =\sum_a(-1)^{\sum_ja_j(d_j+d_{j+1})}
        =2^N(\delta_{b,c}+\delta_{b,\bar c}).
        \notag
    \end{align}
    Equation~(\ref{eq:asymex_kw_character}) imposes $d_j=d_{j+1}$ at
    every site. Such a configuration is constant on the cycle, and
    the two constant configurations are distinct because $N\geq1$.
    For the other Gram matrix, put $d=a\oplus c$ and sum over inputs:
    \begin{align}
        (K_NK_N^\dagger)_{ac}
         & =\sum_b(-1)^{\sum_jb_j(d_j+d_{j-1})}
        =2^N(\delta_{a,c}+\delta_{a,\bar c}).
        \notag
    \end{align}
    This proves~(\ref{eq:asymex_kw_gram}).
\end{proof}

Define the spin-flip sectors $\mathcal H_N^\pm=\{v:\eta v=\pm v\}$ and
their orthogonal projections $P_\pm=(\Id\pm\eta)/2$. Two useful
normalizations of the duality operator are
\begin{align}
    \mathcal D_N & =2^{-N/2}K_N,
    \notag                                               \\
    U_N          & =2^{-(N+1)/2}K_N=\mathcal D_N/\sqrt2.
    \label{eq:asymex_kw_normalizations}
\end{align}

\begin{lemma}[Kernel, image, and partial isometry]
    \label{lem:asymex_kw_sectors}
    \lean{KWExample.kwTensor_mpo_mulVec_eq_zero_iff,
        KWExample.exists_kwTensor_mpo_mulVec_eq_iff,
        KWExample.kwTensor_mpo_rank,
        KWExample.kwDefect_conjTranspose_mul,
        KWExample.kwDefect_mul_conjTranspose,
        KWExample.kwTensor_mpo_normalized_conjTranspose_mul,
        KWExample.kwTensor_mpo_normalized_mul_conjTranspose}
    \leanok
    \uses{thm:asymex_kw_absorption, lem:asymex_kw_gram}
    The kernel and image of the unnormalized duality operator are
    \begin{align}
        \ker K_N  & =\mathcal H_N^-,
        \notag                       \\
        \ran K_N  & =\mathcal H_N^+,
        \notag                       \\
        \rank K_N & =2^{N-1}.
        \label{eq:asymex_kw_sectors}
    \end{align}
    The normalized operators satisfy
    \begin{align}
        \mathcal D_N^\dagger\mathcal D_N
        =\mathcal D_N\mathcal D_N^\dagger & =2P_+,
        \notag                                     \\
        U_N^\dagger U_N=U_NU_N^\dagger    & =P_+.
        \label{eq:asymex_kw_partial}
    \end{align}
    Thus $U_N$ is a partial isometry with initial and final projection
    $P_+$. It vanishes on $\mathcal H_N^-$ and restricts to a unitary
    operator on $\mathcal H_N^+$.
\end{lemma}

\begin{proof}
    \leanok
    If $K_Nv=0$, then~(\ref{eq:asymex_kw_gram}) gives
    $0=K_N^\dagger K_Nv=2^N(v+\eta v)$, so $\eta v=-v$. Conversely,
    Theorem~\ref{thm:asymex_kw_absorption} shows that $K_N$ annihilates
    $\mathcal H_N^-$. The absorption relation $\eta K_N=K_N$ puts the
    image in $\mathcal H_N^+$. Conversely, for $\eta y=y$,
    $K_N(2^{-(N+1)}K_N^\dagger y)=2^{-1}(y+\eta y)=y$, so every vector in
    that sector is in the image. For the rank, $\rank K_N=\rank
        K_N^\dagger K_N=\rank P_+$, and the rank of the orthogonal
    projection $P_+$ is its trace. Since $\eta$ fixes no basis vector,
    $\tr\eta=0$ and $\tr P_+=2^N/2=2^{N-1}$.
    Rescaling~(\ref{eq:asymex_kw_gram}) proves
    (\ref{eq:asymex_kw_partial}).
\end{proof}

Each $H_N(J,h)$ preserves both spin-flip sectors: the field terms
commute with $\eta$, and the two signs acquired by a bond term cancel.
Thus $U_N$ implements a unitary equivalence between $H_N(J,h)$ and
$H_N(h,J)$ on $\mathcal H_N^+$, but not on the full Hilbert space.

\begin{corollary}[Square of the duality]
    \label{cor:asymex_kw_square}
    \lean{KWExample.kwTensor_mpo_mul_self}
    \leanok
    \uses{lem:asymex_kw_gram}
    With $D=K_N$ as above,
    \begin{align}
        D^2 & =2^N(\Id+\eta)T.
        \label{eq:asymex_kw_fusion}
    \end{align}
\end{corollary}

\begin{proof}
    \leanok
    By~(\ref{eq:asymex_kw_adjoint}), $K_N=TK_N^\dagger$ and $T$
    commutes with $K_N$, so $K_N^2=K_NTK_N^\dagger=TK_NK_N^\dagger$.
    Equation~(\ref{eq:asymex_kw_gram}) and $T\eta=\eta T$ give the
    result.
\end{proof}

The corresponding normalized identities are
\begin{align}
    \mathcal D_N^2 & =(\Id+\eta)T,
    \notag                         \\
    U_N^2          & =P_+T,
    \notag                         \\
    U_N^{2N}       & =P_+.
    \label{eq:asymex_kw_normalized_square}
\end{align}
The last identity uses $P_+T=TP_+$ and $T^N=\Id$.
The two summands $T$ and $\eta T$ act on the same Hilbert space.
On $\mathcal H_N^+$ they add, giving $K_N^2=2^{N+1}T$; on
$\mathcal H_N^-$ they cancel. The negative sector is annihilated,
not mapped by flipped translation.

The translation factor depends on identifying the dual lattice with
the original one. In Ref.~\cite{Aasen2016Topological}, identify the
input spin $h_j$ with $b_j$ and the dual spin
$\hat h_{j+1/2}$ with $a_j$. Its primal-to-dual kernel is then
$\mathcal D_N$, since
\begin{align}
    \sum_j(\hat h_{j-1/2}h_j+h_j\hat h_{j+1/2})
     & =\sum_j(a_{j-1}b_j+b_ja_j)=q(a,b).
    \notag
\end{align}
The source returns to the original lattice with the same matrix
elements read with the arguments exchanged, that is, by the transpose
$\mathcal D_N^{\mathsf T}$. Since the kernel is real,
$\mathcal D_N^{\mathsf T}=\mathcal D_N^\dagger$.

\begin{proposition}[Fusion relations of the duality defect]
    \label{prop:asymex_kw_defect_fusion}
    \lean{KWExample.kwDefect, KWExample.spinFlip_mul_self,
        KWExample.kwDefect_mul_spinFlip, KWExample.spinFlip_mul_kwDefect,
        KWExample.kwDefect_transpose_mul, KWExample.kwDefect_mul_transpose,
        KWExample.kwDefect_mul_transpose_mul,
        KWExample.kwTensor_mpo_transpose_mul,
        KWExample.kwTensor_mpo_mul_transpose_mul}
    \leanok
    \uses{lem:asymex_kw_gram, thm:asymex_kw_absorption}
    The normalized duality $\mathcal D_N=2^{-N/2}K_N$ and the spin flip
    satisfy
    \begin{align}
        \eta^2                                           & =\Id,
        \notag                                                             \\
        \mathcal D_N\eta=\eta\mathcal D_N                & =\mathcal D_N,
        \notag                                                             \\
        \mathcal D_N^{\mathsf T}\mathcal D_N
        =\mathcal D_N\mathcal D_N^{\mathsf T}            & =\Id+\eta,
        \notag                                                             \\
        \mathcal D_N\mathcal D_N^{\mathsf T}\mathcal D_N & =2\mathcal D_N.
        \label{eq:asymex_kw_defect_fusion}
    \end{align}
\end{proposition}

\begin{proof}
    \leanok
    The spin flip is an involution, and the second line is
    Theorem~\ref{thm:asymex_kw_absorption} rescaled. The third line is
    (\ref{eq:asymex_kw_gram}) multiplied by $2^{-N}$, because the kernel
    is real. Multiplying $\mathcal D_N\mathcal D_N^{\mathsf T}=\Id+\eta$
    on the right by $\mathcal D_N$ and using $\eta\mathcal D_N=\mathcal
        D_N$ gives the last line.
\end{proof}

With the spin-flip defect equal to $\eta$ and the square of the source
read as $\mathcal D_N^{\mathsf T}\mathcal D_N$, the first three lines
of~(\ref{eq:asymex_kw_defect_fusion}) are the fusion algebra
of~\cite{Aasen2016Topological}, and the last line is its relation
$\mathcal D_\sigma(\mathcal D_\sigma^2-2)=0$. The source computes the
composite on the original lattice only; the identity
$\mathcal D_N\mathcal D_N^{\mathsf T}=\Id+\eta$ on the dual lattice is
not printed there. Repeating the identified endomorphism instead gives
$\mathcal D_N^2=(\Id+\eta)T$.

The partial-isometry convention of Ref.~\cite{Seiberg2024Majorana}
has the scale of $U_N$. Since the kernel is real,
$U_N^\dagger=2^{-(N+1)/2}K_N^{\mathsf T}$.

\begin{proposition}[Relations of the adjoint normalized duality]
    \label{prop:asymex_kw_adjoint_relations}
    \lean{KWExample.kwTranslation, KWExample.kwTranslation_mul_siteX,
        KWExample.kwTranslation_mul_siteZZ, KWExample.kwTranslation_mul_self,
        KWExample.kwTranslation_mul_spinFlip,
        KWExample.spinFlip_mul_kwTranslation,
        KWExample.transpose_translate_mul_kwTranslation,
        KWExample.kwTranslation_conjTranspose,
        KWExample.kwTranslation_mul_conjTranspose,
        KWExample.kwTranslation_conjTranspose_mul,
        KWExample.kwTranslation_pow_two_mul,
        KWExample.transpose_translate_mul_siteX,
        KWExample.transpose_translate_pow_self,
        KWExample.spinFlip_mul_transpose_translate}
    \leanok
    \uses{lem:asymex_kw_local, lem:asymex_kw_gram, thm:asymex_kw_absorption}
    For every site $j$,
    \begin{align}
        U_N^\dagger X_j        & =Z_jZ_{j+1}U_N^\dagger,
        \notag                                           \\
        U_N^\dagger Z_jZ_{j+1} & =X_{j+1}U_N^\dagger.
        \label{eq:asymex_kw_adjoint_local}
    \end{align}
    and
    \begin{align}
        (U_N^\dagger)^2                  & =P_+T^{-1},
        \notag                                                  \\
        U_N^\dagger\eta=\eta U_N^\dagger & =U_N^\dagger,
        \notag                                                  \\
        T^{-1}U_N^\dagger                & =U_N^\dagger T^{-1},
        \notag                                                  \\
        (U_N^\dagger)^\dagger            & =U_N^\dagger T,
        \notag                                                  \\
        U_N^\dagger U_N=U_NU_N^\dagger   & =P_+,
        \notag                                                  \\
        (U_N^\dagger)^{2N}               & =P_+.
        \label{eq:asymex_kw_adjoint_algebra}
    \end{align}
\end{proposition}

\begin{proof}
    \leanok
    Transposing~(\ref{eq:asymex_kw_local_zz}) at
    site $j$ and~(\ref{eq:asymex_kw_local_x}) at site $j+1$
    gives~(\ref{eq:asymex_kw_adjoint_local}), because $X_j$ and
    $Z_jZ_{j+1}$ are real symmetric. Since the kernel is real,
    (\ref{eq:asymex_kw_adjoint}) gives
    $K_N^{\mathsf T}=K_NT^{-1}=T^{-1}K_N$, so
    $(K_N^{\mathsf T})^2=K_N^{\mathsf T}K_NT^{-1}$, and
    (\ref{eq:asymex_kw_gram}) gives the square. The absorption and
    commutation relations follow from
    Theorem~\ref{thm:asymex_kw_absorption} and $TK_N=K_NT$, and the
    adjoint relation from $K_N=K_N^{\mathsf T}T$. The two products
    with the adjoint are~(\ref{eq:asymex_kw_gram}) rescaled. Finally,
    $P_+$ commutes with $T$, $T^N=\Id$ and $P_+^2=P_+$, so
    $(U_N^\dagger)^{2N}=P_+^NT^{-N}=P_+$.
\end{proof}

After exchanging the Pauli operators $X$ and $Z$ at every site, the
relations~(\ref{eq:asymex_kw_adjoint_local})
and~(\ref{eq:asymex_kw_adjoint_algebra}) are those that
Ref.~\cite{Seiberg2024Majorana} prints for its noninvertible
translation, with its translation read as $T^{-1}$ and its projection
$(1+\eta)/2$ read as $P_+$. The proposition shows that $U_N^\dagger$
satisfies these relations. That reference defines its operator by an
explicit product of local factors; Proposition~\ref{prop:asymex_kw_circuit}
below shows that, after the exchange of $X$ and $Z$, this product is
exactly $U_N^\dagger$, with no scalar factor.
The noninvertibility here concerns
the periodic bosonic spin chain, not fermionic Majorana translation.
The translation in~(\ref{eq:asymex_kw_fusion}) must be retained when
comparing this lattice relation with a size-independent fusion algebra.

\begin{definition}[Circuit of the noninvertible translation]
    \label{def:asymex_kw_circuit}
    \lean{KWExample.hadamard, KWExample.ssGateZ, KWExample.ssGateXX,
        KWExample.ssEta, KWExample.ssCircuit}
    \leanok
    Let $H^{\otimes N}$ be the Hadamard gate at every site,
    $\bra{c}H^{\otimes N}\ket{b}=2^{-N/2}(-1)^{\sum_jc_jb_j}$, and let
    $\eta_Z=\prod_jZ_j$. Following~\cite{Seiberg2024Majorana}, with its
    sites $1,\ldots,N$ relabeled $0,\ldots,N-1$, put
    $d^z_j=(\Id+iZ_j)/\sqrt2$, $d^x_j=(\Id+iX_jX_{j+1})/\sqrt2$, and
    \begin{align}
        C_N & =e^{-2\pi iN/8}\,(d^z_0d^x_0)(d^z_1d^x_1)\cdots
        (d^z_{N-2}d^x_{N-2})\,d^z_{N-1}\,\frac{\Id+\eta_Z}{2}.
        \label{eq:asymex_kw_circuit_def}
    \end{align}
    The factors $d^x_j$ run over $0\leq j\leq N-2$; the bond between the
    sites $N-1$ and $0$ does not occur.
\end{definition}

\begin{proposition}[The circuit is the rotated adjoint duality]
    \label{prop:asymex_kw_circuit}
    \lean{KWExample.hadamard_mul_ssCircuit_mul_hadamard,
        KWExample.ssCircuit_mul_siteZ, KWExample.ssCircuit_mul_siteXX,
        KWExample.ssCircuit_mul_ssEta, KWExample.ssEta_mul_ssCircuit,
        KWExample.ssCircuit_mul_conjTranspose,
        KWExample.ssCircuit_conjTranspose_mul, KWExample.ssCircuit_mul_self,
        KWExample.ssCircuit_mulVec_eq_zero_iff,
        KWExample.exists_ssCircuit_mulVec_eq_iff,
        KWExample.transpose_translate_apply,
        KWExample.hadamard_mul_transpose_translate}
    \leanok
    \uses{def:asymex_kw_circuit}
    For every $N\geq1$,
    \begin{align}
        H^{\otimes N}C_NH^{\otimes N} & =U_N^\dagger.
        \label{eq:asymex_kw_circuit}
    \end{align}
    Consequently, for every site $j$, including $j=N-1$, whose bond
    $X_{N-1}X_0$ does not occur in~(\ref{eq:asymex_kw_circuit_def}),
    \begin{align}
        C_NZ_j        & =X_jX_{j+1}C_N,
        \notag                          \\
        C_NX_jX_{j+1} & =Z_{j+1}C_N,
        \notag
    \end{align}
    and
    \begin{align}
        C_N\eta_Z=\eta_ZC_N            & =C_N,
        \notag                                                          \\
        C_NC_N^\dagger=C_N^\dagger C_N & =\frac{\Id+\eta_Z}{2},
        \notag                                                          \\
        C_N^2                          & =\frac{\Id+\eta_Z}{2}\,T^{-1}.
        \notag
    \end{align}
    Here $T^{-1}$ sends $\ket{b_0,\ldots,b_{N-1}}$ to
    $\ket{b_{N-1},b_0,\ldots,b_{N-2}}$, which is the action of the
    product of nearest-neighbor swaps by which~\cite{Seiberg2024Majorana}
    defines its translation, and $T^{-1}$ commutes with $H^{\otimes N}$.
    $C_Nv=0$ exactly when $\eta_Zv=-v$, and $y$ lies in the
    range of $C_N$ exactly when $\eta_Zy=y$: the kernel is the sector
    $\eta_Z=-1$, and on the sector $\eta_Z=+1$ the circuit acts
    unitarily, as stated in~\cite{Seiberg2024Majorana}.
\end{proposition}

\begin{proof}
    \leanok
    \uses{thm:asymex_kw_kernel, prop:asymex_kw_adjoint_relations}
    The Hadamard gate is its own inverse and exchanges $X_j$ and $Z_j$.
    Conjugation therefore maps $d^z_j$ to $(\Id+iX_j)/\sqrt2$, $d^x_j$ to
    $(\Id+iZ_jZ_{j+1})/\sqrt2$, and $\eta_Z$ to $\eta$. In this frame
    the gate $(\Id+iX_k)/\sqrt2$ changes only the bit at site $k$, and
    each site carries exactly one such gate. An entry $\bra{c}G\ket{b}$
    of the product $G$ of the gates, before the projection, is therefore
    a single product along one path of configurations; at the diagonal
    gate of the bond $(k,k+1)$ the two bits are $b_k$ and $c_{k+1}$.
    Writing $[x\neq y]$ for the indicator of $x\neq y$, this gives
    \begin{align}
        \bra{c}G\ket{b}
         & =\frac{i^{[c_{N-1}\neq b_{N-1}]}}{\sqrt2}
        \prod_{k=0}^{N-2}w(c_k,b_k,c_{k+1}),
        \notag                                       \\
        w(x,y,z)
         & =\frac{i^{[x\neq y]}}{\sqrt2}\,
        \frac{1+i(-1)^{y+z}}{\sqrt2}
        =\frac{1+i}{2}\,i^x(-i)^z(-1)^{y(x+z)}.
        \notag
    \end{align}
    The factors $i^{c_k}(-i)^{c_{k+1}}$ telescope to
    $i^{c_0}(-i)^{c_{N-1}}$. The projection adds the entry at $\bar b$,
    whose open-chain signs differ by $(-1)^{c_0+c_{N-1}}$, and
    $e^{-2\pi iN/8}=((1-i)/\sqrt2)^N$. Checking the eight values of
    $(c_0,b_{N-1},c_{N-1})$ shows that the remaining boundary factor is
    $(-1)^{b_{N-1}(c_{N-1}+c_0)}$, the sign of the closing bond. By
    Theorem~\ref{thm:asymex_kw_kernel} the entry is
    $2^{-(N+1)/2}(K_N)_{bc}=\bra{c}U_N^\dagger\ket{b}$, which
    proves~(\ref{eq:asymex_kw_circuit}). The remaining relations are
    those of Proposition~\ref{prop:asymex_kw_adjoint_relations}
    conjugated by the Hadamard gate, which commutes with $T^{-1}$.
    If $C_Nv=0$, then $C_N^\dagger C_Nv=\tfrac12(v+\eta_Zv)=0$;
    conversely, if $\eta_Zv=-v$, then $C_Nv=C_N\eta_Zv=-C_Nv$. The range
    of $C_N$ lies in the sector $\eta_Z=+1$ because $\eta_ZC_N=C_N$, and
    for $\eta_Zy=y$ one has $C_N(2C_N^\dagger y)=y+\eta_Zy=2y$.
\end{proof}

Positive length is essential in these calculations. At $N=1$,
$K_1=\Id+X_0$ and $U_1=P_+$, consistently with the self-bond convention.
The empty virtual word instead has trace $\tr\Id_2=2$, so extending
the unnormalized trace definition to $N=0$ would give $K_0=2$. Its square
would be $4$, whereas the right-hand side of
(\ref{eq:asymex_kw_fusion}) would be $2$.

\subsubsection{Asymmetric compression}

The square identity~(\ref{eq:asymex_kw_fusion}) has a virtual
decomposition, obtained by applying the multi-block compression
theorem to the stacked tensor $B$ of Definition~\ref{def:asymex_kw}.
For the matrix-unit tensors of that definition, expanding the cyclic
trace gives
\begin{align}
    \tr(T^{a_0b_0}\cdots T^{a_{N-1}b_{N-1}})
     & =\prod_j\delta_{a_j,b_{j+1}}=T_{ab},
    \notag                                           \\
    \tr((\eta T)^{a_0b_0}\cdots(\eta T)^{a_{N-1}b_{N-1}})
     & =\prod_j\delta_{a_j,1-b_{j+1}}=(\eta T)_{ab}.
    \notag
\end{align}
The stacked tensor $B$ generates $K_N^2$: multiplication of periodic
operators sums the intermediate physical index at every site.

Normality of a tensor below means that words of some fixed
positive length span its virtual matrix algebra. This differs from
operator normality in~(\ref{eq:asymex_kw_gram}) and does not impose a
normalization on the spectral radius of a transfer map.

\begin{theorem}[Normality of the two translation tensors]
    \label{thm:asymex_kw_normal}
    \lean{KWExample.shiftMPS_isNormal, KWExample.flipShiftMPS_isNormal}
    \leanok
    \uses{def:asymex_kw, def:normal}
    The translation tensor and the flipped translation tensor are normal:
    the four matrices of each are the four matrix units of $\MN{2}$, which
    already span $\MN{2}$ at word length one.
\end{theorem}

\begin{proof}
    \leanok
    Every matrix unit occurs among the four matrices of each tensor, and
    the matrix units span $\MN{2}$. Each tensor is therefore injective
    and hence normal.
\end{proof}

\begin{theorem}[Split compression of the squared duality]
    \label{thm:asymex_kw_compression}
    % Principal declaration: KWExample.kwSquare_dim_eq.
    \lean{KWExample.kwSquare_compression, KWExample.kwSquare_dim_eq,
        KWExample.kwSquare_remainder_eq_zero}
    \leanok
    \uses{def:asymex_kw, thm:asym_multiblock}
    The stacked tensor of Definition~\ref{def:asymex_kw} has the block
    triangular form of Theorem~\ref{thm:asym_multiblock} for the two
    targets $C_1=2T$ and $C_2=2\eta T$, with no zero slots and $4=2+2$.
    There are compression maps $W_s:\C^4\to\C^2$ and
    $V_s:\C^2\to\C^4$, indexed here by $s=1,2$, such that
    $W_sV_t=\delta_{st}\Id_2$ and $W_sB^wV_s=C_s^w$ for every word $w$.
    The remainder vanishes:
    \begin{align}
        R^{ab} & =B^{ab}-\sum_{s=1}^2V_sC_s^{ab}W_s=0.
        \notag
    \end{align}
    Thus the extension splits.
\end{theorem}

\begin{proof}
    \leanok
    \uses{thm:asymex_explicit_gauge}
    The change of bond coordinates is
    \begin{align}
        G & =\frac12\begin{pmatrix}
                        1 & 0 & 1  & 0  \\
                        0 & 1 & 0  & -1 \\
                        1 & 0 & -1 & 0  \\
                        0 & 1 & 0  & 1
                    \end{pmatrix},
        \label{eq:asymex_kw_gauge}
    \end{align}
    whose inverse has integer entries. For each physical pair $(a,b)$,
    direct multiplication gives
    \begin{align}
        GB^{ab}G^{-1} & =(2T^{ab})\oplus(2(\eta T)^{ab}).
        \notag
    \end{align}
    After clearing the denominator, these are equalities of integer
    matrices. Block diagonality gives the triangularity condition and
    the vanishing remainder. The coordinate projections and inclusions
    give the compression pairs by
    Theorem~\ref{thm:asymex_explicit_gauge}; the dimension count is
    clause (vii) of Theorem~\ref{thm:asym_multiblock}.
\end{proof}

\begin{theorem}[Periodic data of the squared duality]
    \label{thm:asymex_kw_trace}
    \lean{KWExample.kwSquare_trace_evalWord}
    \leanok
    \uses{def:asymex_kw}
    For every nonempty word $w$ over the alphabet of index pairs,
    \begin{align}
        \tr(B^w) & =2^{|w|}(\tr(T^w)+\tr((\eta T)^w)).
        \label{eq:asymex_kw_trace}
    \end{align}
    This is~(\ref{eq:asymex_kw_fusion}) at the level of periodic
    coefficients.
\end{theorem}

\begin{proof}
    \leanok
    \uses{thm:asymex_kw_compression, prop:asym_compression_trace}
    Proposition~\ref{prop:asym_compression_trace} applied to the
    compression of Theorem~\ref{thm:asymex_kw_compression} gives
    $\tr(B^w)=\tr((2T)^w)+\tr((2\eta T)^w)$. Pulling the weight $2$ out
    of each word evaluation contributes the factor $2^{|w|}$.
\end{proof}

\begin{theorem}[The sitewise intertwiners of the split case]
    \label{thm:asymex_kw_intertwiners}
    % Principal declaration: KWExample.kwSquare_mul_right0_eq.
    \lean{KWExample.kwSquareLeft0, KWExample.kwSquareLeft1,
        KWExample.kwSquareRight0, KWExample.kwSquareRight1,
        KWExample.kwSquare_mul_right0_eq, KWExample.kwSquare_mul_right1_eq,
        KWExample.left0_mul_kwSquare_eq, KWExample.left1_mul_kwSquare_eq,
        KWExample.kwSquareLeft0_mul_right0, KWExample.kwSquareLeft1_mul_right1}
    \leanok
    \uses{def:asymex_kw, thm:asymex_kw_compression}
    Write
    \begin{align}
        V_1 & =\begin{pmatrix}
                   1 & 0  \\
                   0 & 1  \\
                   1 & 0  \\
                   0 & -1
               \end{pmatrix},
            & V_2
            & =\begin{pmatrix}
                   1  & 0 \\
                   0  & 1 \\
                   -1 & 0 \\
                   0  & 1
               \end{pmatrix},
        \notag                       \\
        W_1 & =\frac12\begin{pmatrix}
                          1 & 0 & 1 & 0  \\
                          0 & 1 & 0 & -1
                      \end{pmatrix},
            & W_2
            & =\frac12\begin{pmatrix}
                          1 & 0 & -1 & 0 \\
                          0 & 1 & 0  & 1
                      \end{pmatrix}.
        \notag
    \end{align}
    Then $W_1V_1=W_2V_2=\Id_2$, and for every index pair,
    $B^{s's}V_1=V_1(2T)^{s's}$, $B^{s's}V_2=V_2(2\eta T)^{s's}$,
    $W_1B^{s's}=(2T)^{s's}W_1$ and $W_2B^{s's}=(2\eta T)^{s's}W_2$.
\end{theorem}

\begin{proof}
    \leanok
    \uses{thm:asymex_explicit_gauge}
    The matrices $W_1,W_2$ are the two row blocks of $G$, and $V_1,V_2$
    are the corresponding column blocks of $G^{-1}$, as in
    Theorem~\ref{thm:asymex_explicit_gauge}. Each intertwining identity
    is verified by multiplying the explicit matrices and clearing the
    denominator. The two normalizations also follow from clause (iv) of
    Theorem~\ref{thm:asym_multiblock}.
\end{proof}

\begin{center}
    % Formula: the four sitewise identities of
    % Theorem~\ref{thm:asymex_kw_intertwiners}: B^{s's}V_1=V_1(2T)^{s's},
    % B^{s's}V_2=V_2(2\eta T)^{s's}, W_1B^{s's}=(2T)^{s's}W_1, and
    % W_2B^{s's}=(2\eta T)^{s's}W_2.
    % Source: Theorem~\ref{thm:asymex_kw_intertwiners}.
    % Ink-to-index map: the ring is the compression map $V_1$, $W_1$,
    % $V_2$ or $W_2$; the square node is the source $B$ or its target
    % $2T$ or $2\eta T$.
    % Contraction set: the single virtual bond joining the ring to the
    % square node in each pair; the physical pair leg remains open.
    % Boundary signature: (virtual-west=1 | virtual-east=1 |
    % physical-up=1 | physical-down=0) for each pictured equality.
    \tenkzkernel
    \begin{tenkz}[rows={wire}, cols=2, west=open, east=open]
        \tn[label pos=s, ports={90:physical:$s's$}]{B} & \tn[skin=ring]{V_1}
    \end{tenkz}
    \quad $=$ \quad
    \begin{tenkz}[rows={wire}, cols=2, west=open, east=open]
        \tn[skin=ring]{V_1} & \tn[label pos=s, ports={90:physical:$s's$}]{2T}
    \end{tenkz}
    \\
    \begin{tenkz}[rows={wire}, cols=2, west=open, east=open]
        \tn[label pos=s, ports={90:physical:$s's$}]{B} & \tn[skin=ring]{V_2}
    \end{tenkz}
    \quad $=$ \quad
    \begin{tenkz}[rows={wire}, cols=2, west=open, east=open]
        \tn[skin=ring]{V_2} &
        \tn[label pos=s, ports={90:physical:$s's$}]{2\eta T}
    \end{tenkz}
    \\
    \begin{tenkz}[rows={wire}, cols=2, west=open, east=open]
        \tn[skin=ring]{W_1} & \tn[label pos=s, ports={90:physical:$s's$}]{B}
    \end{tenkz}
    \quad $=$ \quad
    \begin{tenkz}[rows={wire}, cols=2, west=open, east=open]
        \tn[label pos=s, ports={90:physical:$s's$}]{2T} & \tn[skin=ring]{W_1}
    \end{tenkz}
    \\
    \begin{tenkz}[rows={wire}, cols=2, west=open, east=open]
        \tn[skin=ring]{W_2} & \tn[label pos=s, ports={90:physical:$s's$}]{B}
    \end{tenkz}
    \quad $=$ \quad
    \begin{tenkz}[rows={wire}, cols=2, west=open, east=open]
        \tn[label pos=s, ports={90:physical:$s's$}]{2\eta T} &
        \tn[skin=ring]{W_2}
    \end{tenkz}
\end{center}

\subsection{Action on ground states}%
\label{sec:asymex_kw_action}

The product state and the Greenberger--Horne--Zeilinger (GHZ)
state~\cite{GreenbergerHorneZeilinger1989,Greenberger1990Bell} give
elementary finite-chain examples of the action. Use the unnormalized
vectors
\begin{align}
    p_N & =(\ket{0}+\ket{1})^{\otimes N}=\sum_b\ket{b},
    \notag                                              \\
    g_N & =\ket{\mathbf0}+\ket{\mathbf1}.
    \label{eq:asymex_kw_states}
\end{align}
For $h>0$, $p_N$ is proportional to the unique ground state of
$H_N(0,h)$. For $J>0$ and $N\geq2$, the ground space of $H_N(J,0)$ is
$\spn_{\C}\{\ket{\mathbf0},\ket{\mathbf1}\}$, and $g_N$ spans its
spin-flip-even subspace. These facts follow respectively by minimizing
each field term $-hX_j$ and each diagonal bond term $-JZ_jZ_{j+1}$.
They do not require a thermodynamic limit.

\begin{lemma}[Product-state and GHZ actions]
    \label{lem:asymex_kw_states}
    \lean{KWExample.kwTensor_mpo_mulVec_plus,
        KWExample.kwTensor_mpo_mulVec_single_zero,
        KWExample.kwTensor_mpo_mulVec_single_flipConfig_zero,
        KWExample.kwTensor_mpo_mulVec_ghz,
        KWExample.kwTensor_mpo_normalized_mulVec_plus,
        KWExample.kwTensor_mpo_normalized_mulVec_ghz, KWExample.plusState,
        KWExample.ghzState, KWExample.mpv_kwPlus_eq_smul_mpv_ghz,
        KWExample.mpv_kwGHZ_eq_two_smul_mpv_plusTensor}
    \leanok
    \uses{thm:asymex_kw_kernel, thm:asymex_kw_absorption}
    For every $N\geq1$,
    \begin{align}
        K_Np_N                              & =2^Ng_N,
        \notag                                         \\
        K_N\ket{\mathbf0}=K_N\ket{\mathbf1} & =p_N,
        \notag                                         \\
        K_Ng_N                              & =2p_N.
        \label{eq:asymex_kw_state_actions}
    \end{align}
    Writing $\widehat p_N=2^{-N/2}p_N$ and
    $\widehat g_N=2^{-1/2}g_N$, one has
    $U_N\widehat p_N=\widehat g_N$ and
    $U_N\widehat g_N=\widehat p_N$.
\end{lemma}

\begin{proof}
    \leanok
    Reindexing the exponent and using~(\ref{eq:asymex_kw_character})
    gives
    \begin{align}
        (K_Np_N)_a
         & =\sum_b(-1)^{\sum_jb_j(a_j+a_{j-1})}
        =2^N(\delta_{a,\mathbf0}+\delta_{a,\mathbf1}).
        \notag
    \end{align}
    Both constant inputs have phase one against every output
    configuration, so each maps to $p_N$. Finally,
    $\|p_N\|^2=2^N$ and $\|g_N\|^2=2$ give the normalized statements.
\end{proof}

In particular, the antisymmetric combination
$\ket{\mathbf0}-\ket{\mathbf1}$ is annihilated. These finite-chain
identities alone assert neither spontaneous symmetry breaking nor a
thermodynamic spectral-gap property.

The finite-chain actions in~(\ref{eq:asymex_kw_state_actions}) have
different virtual decompositions. The product-state action splits into
the two GHZ sectors with a local weight $2$. The GHZ action has two
unweighted copies of the same product-state target, together with a
nonzero remainder. The product-state and GHZ tensors are the elementary
matrix product state examples of~\cite[Appendix~A]{Cirac2021Matrix};
here their explicit matrices allow the asymmetric compression and its
sitewise intertwiners to be compared directly.

\begin{definition}[The two action tensors]
    \label{def:asymex_kw_action}
    \lean{KWExample.plusTensor, KWExample.ghz0, KWExample.ghz1,
        KWExample.kwPlus, KWExample.plusTarget, KWExample.ghz,
        KWExample.kwGHZ, KWExample.kwGHZTarget,
        KWExample.mpv_plusTensor_eq_plusState, KWExample.mpv_ghz_eq_ghzState}
    \leanok
    \uses{def:asymex_kw, def:mps_tensor}
    The product-state tensor $A_+$ has bond dimension one and
    $A_+^0=A_+^1=(1)$. Its periodic vector is $p_N$.
    The two ferromagnetic-sector tensors have
    $C_0^i=(\delta_{i,0})$ and $C_1^i=(\delta_{i,1})$.
    Their direct sum is the GHZ tensor
    $A_{\mathrm{GHZ}}^0=\diag(1,0)$,
    $A_{\mathrm{GHZ}}^1=\diag(0,1)$, whose periodic vector is $g_N$.

    For either input tensor $A$, the action tensor is
    $(W\cdot A)^{s'}=\sum_sW^{s's}\otimes A^s$.
    For $A=A_+$ its matrices are
    \begin{align}
        (W\cdot A)^0
         & =\begin{pmatrix}1&1\\1&1\end{pmatrix},
         &
        (W\cdot A)^1
         & =\begin{pmatrix}1&-1\\-1&1\end{pmatrix},
        \notag
    \end{align}
    with targets $2C_0$ and $2C_1$. For $A=A_{\mathrm{GHZ}}$ they are
    \begin{align}
        (W\cdot A)^0
         & =\begin{pmatrix}
                1 & 0 & 1 & 0 \\
                0 & 0 & 0 & 0 \\
                0 & 0 & 0 & 0 \\
                0 & 1 & 0 & 1
            \end{pmatrix},
         &
        (W\cdot A)^1
         & =\begin{pmatrix}
                1 & 0  & -1 & 0 \\
                0 & 0  & 0  & 0 \\
                0 & 0  & 0  & 0 \\
                0 & -1 & 0  & 1
            \end{pmatrix},
        \notag
    \end{align}
    with two unweighted copies of $A_+$ as targets.
\end{definition}

\begin{center}
    % Formula: (W\cdot A)^{s'}=\sum_sW^{s's}\otimes A^s, of
    % Definition~\ref{def:asymex_kw_action}, evaluated at the
    % product-state tensor $A_+$ (left) and at the Greenberger--Horne--
    % Zeilinger tensor $A_{\mathrm{GHZ}}$ (right).
    % Source: Definition~\ref{def:asymex_kw_action}.
    % Ink-to-index map: the upper row is the duality tensor $W$, whose
    % drawn physical legs $r$ are the per-site components of the
    % outgoing index $s'$; the lower row is the acted-on state tensor;
    % the shared vertical pairleg at each site is the summed index $s$.
    % Contraction set: the vertical pairlegs (index $s$) at every site
    % and the horizontal virtual bonds within each row, both closed by
    % the periodic trace; only the upper physical legs of $W$ remain
    % open.
    % Boundary signature: (virtual-west=0 | virtual-east=0 |
    % physical-up=3 | physical-down=0) for both pictures.
    \tenkzkernel
    \begin{tenkz}[rows={op,ket}, cols=4, boundary=periodic]
        \tn[skin=mpo, ports={90:physical:$r_1$}]{W} &
        \tn[skin=mpo, ports={90:physical:$r_2$}]{W} &
        \tn[skin=dots, ports={180:virtual, 0:virtual}]{} &
        \tn[skin=mpo, ports={90:physical:$r_N$}]{W} \\
        \tn[label pos=s, ports={180:virtual, 0:virtual}]{A_+} &
        \tn[label pos=s, ports={180:virtual, 0:virtual}]{A_+} &
        \tn[skin=dots, ports={180:virtual, 0:virtual}]{} &
        \tn[label pos=s, ports={180:virtual, 0:virtual}]{A_+}
    \end{tenkz}
    \qquad
    \begin{tenkz}[rows={op,ket}, cols=4, boundary=periodic]
        \tn[skin=mpo, ports={90:physical:$r_1$}]{W} &
        \tn[skin=mpo, ports={90:physical:$r_2$}]{W} &
        \tn[skin=dots, ports={180:virtual, 0:virtual}]{} &
        \tn[skin=mpo, ports={90:physical:$r_N$}]{W} \\
        \tn[label pos=s, ports={180:virtual, 0:virtual}]{A_{\mathrm{GHZ}}} &
        \tn[label pos=s, ports={180:virtual, 0:virtual}]{A_{\mathrm{GHZ}}} &
        \tn[skin=dots, ports={180:virtual, 0:virtual}]{} &
        \tn[label pos=s, ports={180:virtual, 0:virtual}]{A_{\mathrm{GHZ}}}
    \end{tenkz}
\end{center}

For the product-state action, the matrices are $2Q_+$ and $2Q_-$,
where
\begin{align}
    Q_\pm & =\frac12\begin{pmatrix}1&\pm1\\\pm1&1\end{pmatrix}.
    \notag
\end{align}
They satisfy $Q_\pm^2=Q_\pm$, $Q_+Q_-=Q_-Q_+=0$ and $\tr Q_\pm=1$.
Hence a nonempty word of length $N$ has trace $2^N$ when all letters
agree and trace zero otherwise. This recovers $K_Np_N=2^Ng_N$ and
identifies the factor $2$ as a local tensor weight. By contrast, the
factor $2$ in $K_Ng_N=2p_N$ is the multiplicity of two unweighted
targets. Replacing $p_N$ or $g_N$ by a normalized vector changes the
global scalar as in Lemma~\ref{lem:asymex_kw_states}; it does not
change the weights of these fixed tensors.

In the following two results, $B$ denotes the relevant action tensor,
rather than the stacked square tensor of Definition~\ref{def:asymex_kw}.

\begin{theorem}[The duality on the product state splits]
    \label{thm:asymex_kw_plus}
    % Principal declaration: KWExample.plus_dim_eq.
    \lean{KWExample.plusCompression, KWExample.plus_dim_eq, KWExample.plus_remainder_eq_zero,
        KWExample.kwPlus_trace_evalWord_eq_sum}
    \leanok
    \uses{def:asymex_kw_action, thm:asym_multiblock}
    The first action tensor of Definition~\ref{def:asymex_kw_action} has
    the block triangular form of Theorem~\ref{thm:asym_multiblock} for
    the weighted targets $\widehat C_s=2C_s$, $s=0,1$, with no zero
    slots and $2=1+1$. Its compression pairs satisfy $W_sV_s=\Id_1$
    and $W_sB^wV_s=\widehat C_s^w$ for every word $w$.
    The remainder vanishes identically, and for every nonempty word,
    \begin{align}
        \tr(B^w) & =\tr((2C_0)^w)+\tr((2C_1)^w).
        \notag
    \end{align}
\end{theorem}

\begin{proof}
    \leanok
    \uses{thm:asymex_explicit_gauge, prop:asym_compression_trace}
    Let $F=\begin{pmatrix}1&1\\1&-1\end{pmatrix}$, so
    $F^{-1}=F/2$. Direct multiplication gives
    \begin{align}
        FB^0F^{-1} & =\diag(2,0),
        \notag                    \\
        FB^1F^{-1} & =\diag(0,2).
        \notag
    \end{align}
    Diagonality gives the triangularity condition, identifies the
    matched blocks, and makes the remainder zero. The coordinate
    compression pairs satisfy the stated word identities by
    Theorem~\ref{thm:asymex_explicit_gauge}. The trace identity follows
    from Proposition~\ref{prop:asym_compression_trace}.
\end{proof}

\begin{theorem}[The duality on the ferromagnetic state does not split]
    \label{thm:asymex_kw_ghz}
    % Principal declaration: KWExample.kwGHZ_dim_eq.
    \lean{KWExample.kwGHZCompression, KWExample.kwGHZ_dim_eq,
        KWExample.kwGHZ_evalWord_remainder_eq_zero, KWExample.kwGHZRemainder_eq,
        KWExample.kwGHZ_remainder_mul_remainder, KWExample.kwGHZ_trace_evalWord_eq_two}
    \leanok
    \uses{def:asymex_kw_action, thm:asym_multiblock}
    The second action tensor of Definition~\ref{def:asymex_kw_action}
    has the block triangular form of Theorem~\ref{thm:asym_multiblock}
    for two unweighted copies $C_s=A_+$ of the product-state tensor,
    with $z=2$ zero slots and $4=1+1+2$. The compression pair of each
    copy satisfies $W_sV_s=\Id_1$ and $W_sB^wV_s=C_s^w$ for every word.
    For every nonempty word, $\tr(B^w)=2$.
    Every remainder word of length at least four vanishes; in fact,
    \begin{align}
        R^0 & =\begin{pmatrix}
                   0 & 0 & 1 & 0 \\
                   0 & 0 & 0 & 0 \\
                   0 & 0 & 0 & 0 \\
                   0 & 1 & 0 & 0
               \end{pmatrix},
            & R^1
            & =\begin{pmatrix}
                   0 & 0  & -1 & 0 \\
                   0 & 0  & 0  & 0 \\
                   0 & 0  & 0  & 0 \\
                   0 & -1 & 0  & 0
               \end{pmatrix},
        \notag
    \end{align}
    and $R^iR^j=0$ for all letters $i,j$.
\end{theorem}

\begin{proof}
    \leanok
    \uses{thm:asymex_explicit_gauge, prop:asym_compression_trace}
    The two copies occupy the first and fourth bond coordinates, and
    the zero slots occupy the second and third. Ordering the coordinates
    as $(0,3,1,2)$ makes both letters upper triangular, with diagonal
    $(1,1,0,0)$. Thus the matched blocks are both $A_+$, and the
    coordinate compression pairs give the asserted reductions.
    Every word evaluation of $A_+$ is $1$, so
    Proposition~\ref{prop:asym_compression_trace} gives
    $\tr(B^w)=1+1=2$ for nonempty $w$.

    Subtracting the matched diagonal contributions in the original
    coordinates gives the displayed remainder matrices. Their only
    nonzero entries lie in positions $(0,2)$ and $(3,1)$, so direct
    multiplication gives $R^iR^j=0$. The general bound in clause (vi)
    of Theorem~\ref{thm:asym_multiblock} also gives vanishing at lengths
    at least $|S|+z=4$, where $S$ is the set of two matched slots.
    The obstruction to a sitewise splitting is established in
    Theorem~\ref{thm:asymex_kw_ghz_intertwiner}.
\end{proof}

A word reduction need not provide a pair of sitewise intertwiners.
Here the periodic coefficients have merged into two copies of the same
product state, but the virtual action still has a one-sided obstruction.

\begin{theorem}[One-sided sitewise intertwiners for the merged target]
    \label{thm:asymex_kw_ghz_intertwiner}
    \lean{KWExample.kwGHZ_left_intertwiner_eq_zero,
        KWExample.kwGHZ_right_intertwiner,
        KWExample.kwGHZ_right_intertwiner_ne_zero}
    \leanok
    \uses{def:asymex_kw_action}
    Let $B=W\cdot A_{\mathrm{GHZ}}$. If a row vector $u\in\C^4$
    satisfies $uB^i=u$ for both letters $i$, then $u=0$.
    Thus no nonzero sitewise left intertwiner exists for the
    product-state target. A nonzero sitewise right intertwiner does
    exist: the first coordinate vector satisfies $B^ie_0=e_0$ for both
    letters.
\end{theorem}

\begin{proof}
    \leanok
    Both target matrices are the scalar $1$. In coordinates
    $u=(u_0,u_1,u_2,u_3)$, the left-intertwiner equations are
    \begin{align}
        uB^0 & =(u_0,u_3,u_0,u_3)=u,
        \notag                         \\
        uB^1 & =(u_0,-u_3,-u_0,u_3)=u.
        \notag
    \end{align}
    The second coordinates give $u_3=u_1=0$, and the third give
    $u_0=u_2=0$. For the right intertwiner, the first column of each
    matrix is $e_0$, which is nonzero.
\end{proof}

This obstruction rules out a sitewise split copy of the product-state
target: such a splitting would require a nonzero left intertwiner.
It does not obstruct the word reductions or the periodic trace formula
in Theorem~\ref{thm:asymex_kw_ghz}. Those statements retain the
finite-chain action $K_Ng_N=2p_N$ while distinguishing it from a split
virtual direct sum.

\subsection{The duality as a bimodule over the spin-flip algebra}%
\label{sec:asymex_kw_bimodule}

By~(\ref{eq:asymex_kw_fusion}) (Corollary~\ref{cor:asymex_kw_square}),
$D^2=2^N(\Id+\eta)T$, where the non-trivial translation factor $T$ prevents the
duality $D$ from closing in a fusion algebra with $\Id$ and $\eta$. In the
framework of~\cite{Aasen2016Topological}, it is instead fused with the symmetry
operators on either side and with its own return map.
These relations are recorded below for the unnormalized operator and derived at
the tensor level from the compression theorems, without computing a
gauge.

\begin{lemma}[Spin-flip algebra and the bimodule relations]
    \label{lem:asymex_kw_bimodule}
    \lean{KWExample.flipTensor, KWExample.z2Tensor, KWExample.z2Fusion,
        MPOTensor.mpo_apply_of_bondOne, MPOTensor.mpo_idTensor,
        KWExample.mpo_flipTensor, KWExample.mpo_z2Tensor_zero,
        KWExample.mpo_z2Tensor_one,
        KWExample.isMPOFusionAlgebra_z2Tensor,
        KWExample.kwTensor_mpo_mul_z2Tensor_mpo,
        KWExample.z2Tensor_mpo_mul_kwTensor_mpo, KWExample.kwTransposeTensor,
        KWExample.mpo_kwTransposeTensor,
        KWExample.kwTensor_mpo_mul_kwTransposeTensor_mpo,
        KWExample.kwTransposeTensor_mpo_mul_kwTensor_mpo}
    \leanok
    \uses{def:asymex_kw, thm:asymex_kw_absorption, lem:asymex_kw_gram}
    Let $O_0$ and $O_1$ be the periodic operators of the bond-one tensors
    $O_0^{s's}=\delta_{s's}$ and $O_1^{s's}=\delta_{s',1-s}$, and let
    $W^{\mathsf T}$ be the duality tensor with its physical indices
    exchanged, $(W^{\mathsf T})^{s's}=W^{ss'}$. For every $N>0$,
    $O_0=\Id$, $O_1=\eta$, and $O_aO_b=O_{a+b}$ with $a+b$ taken modulo
    two. The periodic operator of $W^{\mathsf T}$ is $K_N^{\mathsf T}$,
    and
    \begin{align}
        K_NO_a=O_aK_N & =K_N\qquad(a=0,1),
        \notag                             \\
        K_NK_N^{\mathsf T}=K_N^{\mathsf T}K_N
                      & =2^N(O_0+O_1).
        \label{eq:asymex_kw_bimodule}
    \end{align}
\end{lemma}

\begin{proof}
    \leanok
    A periodic contraction of bond-one tensors is the product of its
    one-site entries, which gives $O_0=\Id$ and $O_1=\eta$; the fusion
    rule is $\eta^2=\Id$. Exchanging the physical indices of every tensor
    in the periodic trace transposes the operator. The relations
    (\ref{eq:asymex_kw_bimodule}) are then
    Theorem~\ref{thm:asymex_kw_absorption} and
    (\ref{eq:asymex_kw_gram}), using that the kernel is real.
\end{proof}

\begin{theorem}[Absorption of the spin flip is a gauge equivalence]
    \label{thm:asymex_kw_bimodule_absorption}
    \lean{KWExample.kwTensor_isNormal,
        KWExample.exists_gauge_mulTensor_kwTensor_flipTensor,
        KWExample.exists_gauge_mulTensor_flipTensor_kwTensor}
    \leanok
    \uses{lem:asymex_kw_bimodule, cor:asym_single_block}
    The duality tensor $W$ is normal. The stacked tensors
    $\sum_mW^{s'm}\otimes O_1^{ms}$ and $\sum_mO_1^{s'm}\otimes W^{ms}$,
    of bond dimension two, are each conjugate to $W$: there are
    invertible $2\times2$ matrices $X$ and $Y$ with
    \begin{align}
        \sum_mW^{s'm}\otimes O_1^{ms} & =X^{-1}W^{s's}X,
                                      &
        \sum_mO_1^{s'm}\otimes W^{ms} & =Y^{-1}W^{s's}Y
        \notag
    \end{align}
    for all $s',s$.
\end{theorem}

\begin{proof}
    \leanok
    The four matrices of $W$ give the four matrix units by
    \begin{align}
        E_{00} & =\tfrac12(W^{00}+W^{10}),
               &
        E_{01} & =\tfrac12(W^{00}-W^{10}),
        \notag                             \\
        E_{10} & =\tfrac12(W^{01}-W^{11}),
               &
        E_{11} & =\tfrac12(W^{01}+W^{11}),
        \notag
    \end{align}
    so $W$ is injective and hence normal. The
    periodic operators of both stacked tensors equal $K_N$ for every
    $N>0$ by~(\ref{eq:asymex_kw_bimodule}). Corollary~\ref{cor:asym_single_block}
    gives $V$ and $W'$ with $W'V=\Id_2$ whose remainder vanishes on every
    word of length at least $2-2+1=1$. The remainder therefore vanishes
    letter by letter, and $V$ is square with left inverse $W'$.
\end{proof}

Explicitly, the stacked tensors have the letters
$\sum_mW^{s'm}\otimes O_1^{ms}=W^{s',1-s}$ and
$\sum_mO_1^{s'm}\otimes W^{ms}=W^{1-s',s}$, and a direct check on the
four matrices of Definition~\ref{def:asymex_kw} gives
\begin{align}
    W^{s',1-s} & =XW^{s's}X,
               &
    W^{1-s',s} & =YW^{s's}Y,
               &
    X          & =\begin{pmatrix}0&1\\1&0\end{pmatrix},
               &
    Y          & =\begin{pmatrix}1&0\\0&-1\end{pmatrix},
    \notag
\end{align}
with $X^{-1}=X$ and $Y^{-1}=Y$. The proof above does not use these
matrices.

\begin{theorem}[Fusion of the duality with its adjoint]
    \label{thm:asymex_kw_bimodule_fusion}
    \lean{KWExample.z2Tensor_isNormal,
        KWExample.exists_compression_mulTensor_kwTensor_kwTransposeTensor,
        KWExample.exists_compression_mulTensor_kwTransposeTensor_kwTensor}
    \leanok
    \uses{lem:asymex_kw_bimodule, thm:asym_mpo_product, thm:asym_multiblock}
    The stacked tensor $\sum_mW^{s'm}\otimes(W^{\mathsf T})^{ms}$, of
    bond dimension four, has the block triangular form of
    Theorem~\ref{thm:asym_multiblock} for the two bond-one targets
    $2O_0$ and $2O_1$, with exactly two zero slots: $4=1+1+2$. The same
    holds for $\sum_m(W^{\mathsf T})^{s'm}\otimes W^{ms}$.
\end{theorem}

\begin{proof}
    \leanok
    Both bond-one tensors are normal, since each has a one-site matrix
    equal to $1$. By~(\ref{eq:asymex_kw_bimodule}), the periodic
    operators of the stacked tensor are $\sum_a2^NO_a$ for every $N>0$.
    Theorem~\ref{thm:asym_mpo_product} applies with weights $2$, and the
    dimension count of Theorem~\ref{thm:asym_multiblock} gives the number
    of zero slots.
\end{proof}

The biorthogonal compression pairs of this theorem are fusion tensors
for the channels $\Id$ and $\eta$ in the product of the duality with
its adjoint. The stacked square of Theorem~\ref{thm:asymex_kw_compression}
has no zero slot, whereas the product with the adjoint has two.
%
\section{Fibonacci fusion operators}%
\label{sec:asymex_fibonacci}

The Fibonacci fusion category has two simple objects, the vacuum $1$
and $\tau$, with the fusion rule $\tau\otimes\tau=1\oplus\tau$. Its
string-net construction produces matrix product operators whose fusion
rule, fusion tensors, and $F$-symbols are given
in~\cite[Appendix~D.1]{Bultinck2017Anyons}; the review displays the
string-net operator built from the same Fibonacci data
in~\cite[Appendix~A, ``The MPO for the Fibonacci model'']{Cirac2021Matrix},
and the general string-net construction of such operator symmetries
and their intertwiners is in~\cite{Lootens2021MPOStringNet}. The
physical space used here is the two-label space of the golden
chain~\cite{Feiguin2007GoldenChain}, one label $1$ or $\tau$ per site,
rather than the full triple-index alphabet of the categorical
string-net construction. The entries of
the blocks defined below are obtained from the $F$-symbols
of~\cite[Appendix~D.1]{Bultinck2017Anyons}. From these data one obtains
two operators on the chain whose fusion relation and action on periodic
states reduce to finite matrix computations.

\begin{definition}[The string-net tensors of the source and of the review]
    \label{def:asymex_fib_gsymbol_tensor}
    \lean{FibonacciCompression.fibLoopFactor, FibonacciCompression.fibGSymbol,
        FibonacciCompression.fibBondLetter, FibonacciCompression.fibSiteLetter,
        FibonacciCompression.fibEdgeTauLetter, FibonacciCompression.fibStringNetTensor,
        FibonacciCompression.fibStringNetEdgeTau, FibonacciCompression.fibReviewTensor,
        FibonacciCompression.fibReviewWeight,
        FibonacciCompression.fibReviewWeightedTensor}
    \leanok
    \uses{def:asymex_fib_fsymbols}
    Let $d_1=1$, $d_\tau=\varphi=(1+\sqrt5)/2$ be the quantum dimensions,
    and let $N_{ab}^c$ denote the Fibonacci fusion multiplicities.
    Over $\C$ write $v_a=\sqrt{d_a}$ and $G^{abc}_{def}=v_e^{-1}v_f^{-1}[F^{abc}_d]_e^f$.
    For a label $f$, the bond letters of the block $f$ are the pairs
    $(u,l)$ of labels with $N_{ul}^f>0$: the letters $(1,1)$, $(\tau,\tau)$
    for $f=1$ and $(1,\tau)$, $(\tau,1)$, $(\tau,\tau)$ for $f=\tau$. A
    physical letter of the source's tensor is a pair $(x,e)$ of a plaquette
    label and an edge label. The operator tensor
    of~\cite[Section~6.2 and Appendix~D.1]{Bultinck2017Anyons} is drawn as the
    crossing of the
    operator line $f$ with an edge line $e$, with the plaquettes $b$ upper
    left, $q'$ upper right, $c$ lower left and $q''$ lower right; its
    entries are
    \begin{align}
        \bigl(S_f^{(b,e'),(c,e)}\bigr)_{p,(q',q'')}
         & =\delta_{e'e}\,\delta_{p,(b,c)}\,
        G^{q'bc}_{q''ef}\sqrt{v_{q'}v_bv_cv_{q''}},
        \label{eq:asymex_fib_gsymbol_tensor}
    \end{align}
    and its letters at edge label $\tau$ are
    $S_f^{x'x}=S_f^{(x',\tau),(x,\tau)}$. The periodic operators of $S_f$
    act on $(\C^4)^{\otimes N}$, one plaquette label and one edge label
    per site. The configurations whose edge labels are all $\tau$ span a
    subspace identified with $(\C^2)^{\otimes N}$, and the letters $S_f^{x'x}$ are
    those of $S_f$ on this subspace. The operator tensor
    of~\cite[Appendix~A, ``The MPO for the Fibonacci model'']{Cirac2021Matrix},
    at the vertical label $\tau$, has the entries
    \begin{align}
        \bigl(R_f^{x'x}\bigr)_{p,(q',q'')}
         & =\delta_{p,(x',x)}\,\frac{1}{\sqrt{d_{x'}d_{q''}}}\,
        [F^{f x\tau}_{q'}]_{x'}^{q''},
        \label{eq:asymex_fib_review_tensor}
    \end{align}
    with the $F$-symbols of Definition~\ref{def:asymex_fib_fsymbols} and their
    selection rule, in place of the selection rule printed in the review.
    On configurations $x\in\{1,\tau\}^N$ let
    $W=\operatorname{diag}\prod_k d_{x_k}$, with square roots
    $W^{\pm1/2}=\operatorname{diag}\prod_k d_{x_k}^{\pm1/2}$, and let $\widetilde R_f$ be
    the tensor with letters $d_x R_f^{x'x}$.
\end{definition}

The printed $G$-symbol tensor, including its quantum-dimension factors,
is given in Definition~\ref{def:asymex_fib_gsymbol_tensor}. Its restriction
to edge label $\tau$ is related to the two-label blocks by the diagonal
bond similarity of Theorem~\ref{thm:asymex_fib_gsymbol_similarity}.
The review uses a different physical normalization; its relation to these
blocks is the congruence of Theorem~\ref{thm:asymex_fib_review_congruence}.

For every $N>0$, close the bond indices by the ordinary, unnormalized
trace of the ordered product of local matrices. The first physical
index is the output and the second is the input.

The word-span certificates are in Theorem~\ref{thm:asymex_fib_normal}.
The admissibility projector and its action as the fusion unit are
identified in Lemma~\ref{lem:asymex_fib_vacuum} and
Theorem~\ref{thm:asymex_fib_unit}. The fusion relation is proved for the
two-label blocks in Theorem~\ref{thm:asymex_fib_fusion} and for the
source's blocks on both edge labels in
Theorem~\ref{thm:asymex_fib_stringnet_fusion}. Commutation with the
golden-chain Hamiltonian is established for $N\geq2$ in
Theorem~\ref{thm:asymex_fib_golden_chain}; the regular action on the two
explicit state tensors is in Theorem~\ref{thm:asymex_fib_state_action}.

The physical space at each site is $\C^2$, with computational labels
$0$ for the vacuum and $1$ for $\tau$. At system size $N\geq1$, write
$\mc H_N=(\C^2)^{\otimes N}$ and use the orthonormal basis
$\ket{x}=\ket{x_0,\ldots,x_{N-1}}$. Site indices are cyclic:
$x_N=x_0$. Throughout, the golden-number convention is
$\varphi=(1+\sqrt5)/2$ and
$\sigma=\varphi^{-1/2}=\sqrt{(\sqrt5-1)/2}$.
In particular, $\sigma^2=\varphi^{-1}$ and
$\sigma^4+\sigma^2=1$.

For an MPO tensor $T^{yx}$, the first physical index is the output
index. Its periodic operator is defined by
\begin{align}
    \bra{y}O_N(T)\ket{x}
     & =\tr(T^{y_0x_0}\cdots T^{y_{N-1}x_{N-1}}).
    \label{eq:asymex_fib_periodic}
\end{align}
For the tensors $B_1$ and $B_\tau$ defined below, write
$O_1=O_N(B_1)$ and $O_\tau=O_N(B_\tau)$ when $N$ is fixed.
Their non-invertible fusion relation is
\begin{align}
    O_\tau O_\tau & =O_1+O_\tau,
    \label{eq:asymex_fib_fusion}
\end{align}
as proved below for every positive system size. The label $\tau$ is not
invertible in the sense of Definition~\ref{def:fusion_unit_invertible}. The
vacuum operator $O_1$ is a projection onto admissible configurations, not
the identity on $\mc H_N$, as for the unit of any fusion algebra
(Remark~\ref{rem:mposym_unit_projector}).

Every entry of the tensors below, and of the changes of bond
coordinates in the proofs, lies in a quartic ring generated by
$\sigma$. The identities of this section are first verified in that
ring and then embedded into $\C$. The exceptions are the string-net
tensors of Definition~\ref{def:asymex_fib_gsymbol_tensor} and the bond
similarity relating them to the blocks, whose entries involve
$\varphi^{\pm1/4}$; those identities are proved over $\C$.

\begin{definition}[The ring of golden integers]
    \label{def:asymex_golden_ring}
    \lean{GoldenInt, GoldenInt.sigma, goldenPoly, goldenSigmaReal,
        goldenToComplex, MPSTensor.complexOfGolden, MPSTensor.evalWordGolden,
        MPSTensor.mulGoldenTensor}
    \leanok
    The ring of \emph{golden integers} is $\Z[\sigma]$, where
    $\sigma=\varphi^{-1/2}=\sqrt{(\sqrt5-1)/2}$ is the positive real
    root of $X^4+X^2-1$. It is presented by the four integer
    coordinates of $c_0+c_1\sigma+c_2\sigma^2+c_3\sigma^3$,
    with multiplication reduced using $\sigma^4=1-\sigma^2$.
    It embeds into $\C$ through the quotient
    $\mathbb{Q}[X]/(X^4+X^2-1)$ evaluated at this real root.
    For matrices, the embedding is applied entrywise.
    Word evaluation and the bond-space product of tensors over
    $\Z[\sigma]$ use the same formulas as over $\C$.
\end{definition}

\begin{theorem}[Exact arithmetic in the ring of golden integers]
    \label{thm:asymex_golden_ring}
    % Principal declaration: goldenMulRule.
    \lean{GoldenInt.sigma_quartic, goldenMulRule, goldenSigmaReal_quartic,
        MPSTensor.complexOfGolden_mul, MPSTensor.evalWord_complexOfGolden,
        MPSTensor.mulTensor_complexOfGolden}
    \leanok
    \uses{def:asymex_golden_ring}
    The generator satisfies $\sigma^4+\sigma^2-1=0$ in $\Z[\sigma]$, and so
    does the real number $\sqrt{(\sqrt5-1)/2}$. In any commutative ring
    carrying a root $s$ of $X^4+X^2-1$, the product of two integer
    combinations of $1,s,s^2,s^3$ is the combination prescribed by the
    multiplication of $\Z[\sigma]$. The entrywise embedding of matrices
    over $\Z[\sigma]$ into the complex matrices is multiplicative, and it
    commutes with word evaluation and with the bond-space product of
    tensors.
\end{theorem}

\begin{proof}
    \leanok
    Squaring $\sigma^2=(\sqrt5-1)/2$ gives the quartic for the real root.
    The multiplication rule is the reduction of a product of two
    polynomials of degree at most three modulo $X^4+X^2-1$.
    It is a polynomial identity in their eight integer coefficients.
    The entrywise embedding is induced by a ring homomorphism, so
    expansion of a matrix product proves its multiplicativity.
    Induction on the word proves compatibility with word evaluation.
    Expansion of the bond-space product proves the remaining
    compatibility statement.
\end{proof}

\begin{definition}[The two Fibonacci blocks and the stacked square]
    \label{def:asymex_fib}
    \lean{FibonacciCompression.fibTau, FibonacciCompression.fibOne,
        FibonacciCompression.fibTauMPS, FibonacciCompression.fibOneMPS,
        FibonacciCompression.fibStack, FibonacciCompression.fibTargets}
    \leanok
    \uses{def:asymex_golden_ring, def:mpo_tensor, def:mpdo_operator_product}
    Read a physical letter as the pair $(x',x)$ of outgoing and incoming
    labels, ordered as $(0,0),(0,1),(1,0),(1,1)$. The reduced vertex
    weight is
    \begin{align}
        F & =\begin{pmatrix}
                 0 & \varphi^{-1}   & \varphi^{-1/2} \\
                 1 & 0              & 1              \\
                 1 & \varphi^{-1/2} & -\varphi^{-1}
             \end{pmatrix}.
        \label{eq:asymex_fib_vertex}
    \end{align}
    The $\tau$ tensor $B_\tau$ has bond dimension three and letters
    \begin{align}
        B_\tau^{00} & =0.
        \label{eq:asymex_fib_tau_zero}
    \end{align}
    \begin{align}
        B_\tau^{01} & =\begin{pmatrix}
                           0 & \sigma^2 & \sigma \\0&0&0\\0&0&0
                       \end{pmatrix}.
        \label{eq:asymex_fib_tau_01}
    \end{align}
    \begin{align}
        B_\tau^{10} & =\begin{pmatrix}
                           0 & 0 & 0 \\1&0&1\\0&0&0
                       \end{pmatrix}.
        \label{eq:asymex_fib_tau_10}
    \end{align}
    \begin{align}
        B_\tau^{11} & =\begin{pmatrix}
                           0 & 0 & 0 \\0&0&0\\1&\sigma&-\sigma^2
                       \end{pmatrix}.
        \label{eq:asymex_fib_tau_11}
    \end{align}
    Thus each nonzero letter carries one row of
    (\ref{eq:asymex_fib_vertex}), with the other rows zero.
    The vacuum tensor $B_1$ has bond dimension two and letters
    \begin{align}
        B_1^{00} & =\begin{pmatrix}0&1\\0&0\end{pmatrix}.
        \label{eq:asymex_fib_one_00}
    \end{align}
    \begin{align}
        B_1^{11} & =\begin{pmatrix}0&0\\1&1\end{pmatrix}.
        \label{eq:asymex_fib_one_11}
    \end{align}
    \begin{align}
        B_1^{01} & =B_1^{10}=0.
        \label{eq:asymex_fib_one_off}
    \end{align}
    The stacked square has bond dimension nine:
    \begin{align}
        B^{x'x} & =\sum_{m=0}^1 B_\tau^{x'm}\otimes B_\tau^{mx}.
        \label{eq:asymex_fib_stack}
    \end{align}
    Viewed as MPS tensors over the four-letter pair alphabet, its two
    target blocks are $B_1$ and $B_\tau$, each with multiplicity one
    and weight one.
\end{definition}

The review writes the operator with the prefactor $1/\sqrt{d_Ad_D}$
arising from the normalization of the fusion basis, where $d_A$ and
$d_D$ are the quantum dimensions of the labels on the adjacent
bonds~\cite[Appendix~A, ``The MPO for the Fibonacci model'']{Cirac2021Matrix};
that prefactor is omitted in the blocks defined here, and its effect on the
periodic operators is Theorem~\ref{thm:asymex_fib_review_congruence}.

\begin{theorem}[The one-site matrices of the $\tau$ block]
    \label{thm:asymex_fib_entries}
    \lean{FibonacciCompression.fibTauMPS_apply_one,
        FibonacciCompression.fibTauMPS_apply_three}
    \leanok
    \uses{def:asymex_fib}
    The matrix of $B_\tau$ at the letter $(0,1)$ has the first row of
    (\ref{eq:asymex_fib_vertex}) as its first row and zero elsewhere, and
    its matrix at the letter $(1,1)$ has the third row of
    (\ref{eq:asymex_fib_vertex}) as its third row and zero elsewhere, with
    $\varphi^{-1/2}=\sqrt{(\sqrt5-1)/2}$.
\end{theorem}

\begin{proof}
    \leanok
    \uses{thm:asymex_golden_ring}
    Both matrices are obtained by embedding their coefficients in
    $\Z[\sigma]$ into $\C$ entrywise. Substitution of
    $\sigma=\sqrt{(\sqrt5-1)/2}$ gives
    (\ref{eq:asymex_fib_tau_01}) and
    (\ref{eq:asymex_fib_tau_11}).
\end{proof}

The blocks above are built from the categorical data of the Fibonacci
theory, which~\cite[Appendix~D.1]{Bultinck2017Anyons} and the review
print in the following form. The former gives the operator tensor only
as a diagram and prints no entries of $B_1$ or $B_\tau$; the placement of the $F$-symbols used
here, without the closed-loop factors of the source and without the
prefactor of the review, is recorded in~\cite{gap:bmwshv17_fibonacci_block_entries_provenance}.
Theorem~\ref{thm:asymex_fib_gsymbol_similarity} relates these blocks to the
source's tensor built from the $G$-symbols, and
Theorem~\ref{thm:asymex_fib_stringnet_fusion} carries their fusion rules to
that tensor on both edge labels.

\begin{definition}[Fibonacci fusion data of the review]
    \label{def:asymex_fib_fsymbols}
    \lean{FibonacciCompression.fibQuantumDimGolden,
        FibonacciCompression.fibQuantumDimInvSqrt,
        FibonacciCompression.fibFTauGolden,
        FibonacciCompression.fibFSymbolGolden,
        FibonacciCompression.fibGSymbolGolden}
    \leanok
    \uses{def:asymex_golden_ring, def:asymex_fib_anomaly_nim}
    The Fibonacci theory has the labels $1$ and $\tau$, with fusion
    multiplicities
    $N_{11}^1=N_{\tau1}^\tau=N_{1\tau}^\tau=N_{\tau\tau}^1=N_{\tau\tau}^\tau=1$
    and all others zero, and quantum dimensions $d_1=1$ and
    $d_\tau=\varphi=1+\sigma^2$, the Perron--Frobenius dimensions of
    Example~\ref{ex:mposym_fibonacci_dimension}. Write $v_a=\sqrt{d_a}$, so that
    $v_1^{-1}=1$ and $v_\tau^{-1}=\sigma$. For labels $a,b,c,d,e,f$,
    with $\delta_{ijk}=1$ when $N_{ij}^k>0$ and $\delta_{ijk}=0$
    otherwise, the $F$-symbols are
    \begin{align}
        [F^{abc}_d]_e^f
         & =\delta_{abe}\delta_{cde}\delta_{adf}\delta_{bcf}\,F^{abc}_{def},
        \label{eq:asymex_fib_fsymbol}
    \end{align}
    where $F^{abc}_{def}=1$ unless $a=b=c=d=\tau$, and the block at
    $a=b=c=d=\tau$, indexed by $e,f\in\{1,\tau\}$, is
    \begin{align}
        F^{\tau\tau\tau}_\tau
         & =\begin{pmatrix}
                \sigma^2 & \sigma    \\
                \sigma   & -\sigma^2
            \end{pmatrix}
        =\begin{pmatrix}
             1/\varphi      & 1/\sqrt\varphi \\
             1/\sqrt\varphi & -1/\varphi
         \end{pmatrix}.
        \label{eq:asymex_fib_ftau}
    \end{align}
    The $G$-symbols are
    $G^{abc}_{def}=v_e^{-1}v_f^{-1}[F^{abc}_d]_e^f$. All these
    quantities are taken in $\Z[\sigma]$. This is the data
    of~\cite[Appendix~A, ``The MPO for the Fibonacci model'']{Cirac2021Matrix}
    and~\cite[Appendix~D.1]{Bultinck2017Anyons}.
\end{definition}

\begin{theorem}[Values of the Fibonacci fusion data]
    \label{thm:asymex_fib_fsymbol_values}
    % Principal declaration: FibonacciCompression.complexOfGolden_fibFTauGolden.
    \lean{FibonacciCompression.complexOfGolden_fibFTauGolden,
        FibonacciCompression.fibNim_eq_ite,
        FibonacciCompression.fibQuantumDimGolden_one_sq,
        FibonacciCompression.fibQuantumDimInvSqrt_sq_mul,
        FibonacciCompression.fibFTauGolden_mul_self,
        FibonacciCompression.fibFSymbolGolden_tau,
        FibonacciCompression.goldenToComplex_fibQuantumDimGolden_one,
        goldenSigmaReal_sq, goldenSigmaReal_eq_inv_sqrt_goldenRatio}
    \leanok
    \uses{def:asymex_fib_fsymbols, def:asymex_fib_anomaly_nim,
        def:asymex_golden_ring}
    The multiplicities $N_{ab}^c$ of
    Definition~\ref{def:asymex_fib_anomaly_nim} are the Fibonacci
    fusion multiplicities listed in
    Definition~\ref{def:asymex_fib_fsymbols}. In $\Z[\sigma]$,
    $d_\tau^2=d_\tau+1$, $v_a^{-2}d_a=1$ for both labels, the block
    (\ref{eq:asymex_fib_ftau}) squares to the identity, and it is the
    block of the $F$-symbols (\ref{eq:asymex_fib_fsymbol}) at
    $a=b=c=d=\tau$. Under the embedding into $\C$, $d_\tau$ is the
    golden ratio, $\sigma^2=\varphi^{-1}$, $\sigma=1/\sqrt\varphi$, and
    the block (\ref{eq:asymex_fib_ftau}) has the printed complex entries
    $1/\varphi$, $1/\sqrt\varphi$, $1/\sqrt\varphi$, $-1/\varphi$.
\end{theorem}

\begin{proof}
    \leanok
    \uses{thm:asymex_golden_ring}
    The identities in $\Z[\sigma]$ are finite computations with
    coefficients reduced using $\sigma^4=1-\sigma^2$; the
    multiplicities are compared case by case. Over the reals,
    $\sigma^2=(\sqrt5-1)/2=\varphi^{-1}$, so
    $1+\sigma^2=1+\varphi^{-1}=\varphi$, and $\sigma$ is the
    nonnegative square root of $\varphi^{-1}$.
\end{proof}

\begin{theorem}[The $\tau$ block is the $F$-symbol tensor]
    \label{thm:asymex_fib_tau_fsymbol}
    % Principal declaration: FibonacciCompression.fibTau_eq_fSymbol.
    \lean{FibonacciCompression.fibTau_eq_fSymbol,
        FibonacciCompression.fibTauGolden_eq_fSymbol,
        FibonacciCompression.fibLetter}
    \leanok
    \uses{def:asymex_fib, def:asymex_fib_fsymbols, def:asymex_golden_ring}
    Index the three bond coordinates of $B_\tau$ by the admissible
    letters $(1,\tau)$, $(\tau,1)$, $(\tau,\tau)$, in this order. For
    every physical letter $(x',x)$, every incoming bond letter $p$ and
    every outgoing bond letter $q=(q',q'')$,
    \begin{align}
        (B_\tau^{x'x})_{pq}
         & =\begin{cases}
                [F^{\tau x\tau}_{q'}]_{x'}^{q''} & \text{if } p=(x',x), \\
                0                                & \text{otherwise},
            \end{cases}
        \label{eq:asymex_fib_tau_fsymbol}
    \end{align}
    both in $\Z[\sigma]$ and, after the entrywise embedding, over $\C$.
\end{theorem}

\begin{proof}
    \leanok
    \uses{thm:asymex_golden_ring}
    For each of the four physical letters, both sides are $3\times3$
    matrices over $\Z[\sigma]$ and are compared entrywise using
    (\ref{eq:asymex_fib_fsymbol}) and (\ref{eq:asymex_fib_ftau}). The
    complex statement follows by the entrywise embedding.
\end{proof}



\begin{theorem}[The source's blocks on edge label $\tau$]
    \label{thm:asymex_fib_gsymbol_similarity}
    % Principal declaration: FibonacciCompression.fibStringNetEdgeTau_conj.
    \lean{FibonacciCompression.fibStringNetEdgeTau_conj,
        FibonacciCompression.fibGSymbolGolden_tetrahedral,
        FibonacciCompression.fibOneGolden_eq_fSymbol,
        FibonacciCompression.fibBlock_apply_eq_fSymbol,
        FibonacciCompression.fibStringNetGauge,
        FibonacciCompression.fibStringNetGauge_ne_zero,
        MPOTensor.mpo_eq_of_diagonal_conj,
        FibonacciCompression.mpo_fibStringNetEdgeTau,
        FibonacciCompression.mpo_fibStringNetTensor_edgeTau,
        MPOTensor.mpo_apply_comp}
    \leanok
    \uses{def:asymex_fib_gsymbol_tensor, thm:asymex_fib_tau_fsymbol}
    The $G$-symbols satisfy $G^{abc}_{def}=G^{fce}_{abd}$, and both blocks
    are $F$-symbol tensors:
    $(B_f^{x'x})_{p,(q',q'')}=\delta_{p,(x',x)}[F^{fx\tau}_{q'}]_{x'}^{q''}$
    for $f=1,\tau$. Let $h$ be the diagonal
    matrix on the bond letters with $h(u,l)=(v_u/v_l)^{1/2}$. Then
    \begin{align}
        h\,S_f^{x'x}\,h^{-1} & =B_f^{x'x}
        \label{eq:asymex_fib_gsymbol_conj}
    \end{align}
    for both labels $f$ and all $x',x$. Consequently, on the
    configurations whose edge labels are all $\tau$, the periodic operators
    of $S_1$ and $S_\tau$ equal those of $B_1$ and $B_\tau$ at every length.
\end{theorem}

\begin{proof}
    \leanok
    \uses{thm:asymex_golden_ring}
    The symmetry of the $G$-symbols and the $F$-symbol form of $B_1$ are
    finite computations in $\Z[\sigma]$. By the symmetry,
    $G^{q'x'x}_{q''\tau f}=v_{x'}^{-1}v_{q''}^{-1}[F^{fx\tau}_{q'}]_{x'}^{q''}$,
    so the nonzero entries of $S_f^{x'x}$ are those of $B_f^{x'x}$ multiplied
    by
    $\sqrt{v_{q'}v_{x'}v_xv_{q''}}/(v_{x'}v_{q''})=h(q',q'')/h(x',x)$, which
    is (\ref{eq:asymex_fib_gsymbol_conj}). The similarity $h$ cancels in the
    trace defining the periodic operator, so the operators agree.
\end{proof}

\begin{theorem}[Fusion rules of the source's blocks on both edge labels]
    \label{thm:asymex_fib_stringnet_fusion}
    % Principal declaration: FibonacciCompression.isMPOFusionAlgebra_fibStringNetTensor.
    \lean{FibonacciCompression.isMPOFusionAlgebra_fibStringNetTensor,
        FibonacciCompression.fibEdgeOneLetter,
        FibonacciCompression.exists_fibEdgeOneLetter_or_fibEdgeTauLetter,
        FibonacciCompression.fibEdgeGraded, FibonacciCompression.fibEdgeOneGolden,
        FibonacciCompression.fibOneFull, FibonacciCompression.fibTauFull,
        FibonacciCompression.fibBlockFull,
        FibonacciCompression.fibGSymbolGolden_edgeOne,
        FibonacciCompression.fibStringNetTensor_edgeOne_conj,
        FibonacciCompression.fibStringNetTensor_conj,
        FibonacciCompression.mpo_fibStringNetTensor,
        FibonacciCompression.fibConjFull, FibonacciCompression.fibConjFull_mul_eq,
        FibonacciCompression.fibOneOneFull_compression,
        FibonacciCompression.fibOneTauFull_compression,
        FibonacciCompression.fibTauOneFull_compression,
        FibonacciCompression.fibTauTauFull_compression,
        FibonacciCompression.isMPOFusionAlgebra_fibBlockFull,
        FibonacciCompression.fibSiteLetter_fibEdgeOneLetter,
        FibonacciCompression.fibEdgeGraded_edgeTau,
        FibonacciCompression.fibEdgeGraded_edgeOne,
        FibonacciCompression.fibEdgeGraded_edgeOne_edgeTau,
        FibonacciCompression.fibEdgeGraded_edgeTau_edgeOne,
        FibonacciCompression.goldenToComplex_fibEdgeOneGolden,
        FibonacciCompression.fibBlockFull_edgeTau,
        FibonacciCompression.fibBlockFull_edgeOne_apply,
        FibonacciCompression.fibBlockFull_edgeOne_edgeTau,
        FibonacciCompression.fibBlockFull_edgeTau_edgeOne,
        FibonacciCompression.fibBlockFull_toMPSTensor_eq,
        FibonacciCompression.fibOneOneFullStack_mul_gaugeInv,
        FibonacciCompression.fibOneTauFullStack_mul_gaugeInv,
        FibonacciCompression.fibTauOneFullStack_mul_gaugeInv,
        FibonacciCompression.fibTauTauFullStack_mul_gaugeInv,
        FibonacciCompression.fibOneOneEdgeOneStackGolden,
        FibonacciCompression.fibOneTauEdgeOneStackGolden,
        FibonacciCompression.fibTauOneEdgeOneStackGolden,
        FibonacciCompression.fibTauTauEdgeOneStackGolden,
        FibonacciCompression.fibOneOneFullStack_eq, FibonacciCompression.fibOneTauFullStack_eq,
        FibonacciCompression.fibTauOneFullStack_eq, FibonacciCompression.fibTauTauFullStack_eq,
        MPOTensor.IsMPOFusionAlgebra.of_mpo_eq_mul_mul}
    \leanok
    \uses{thm:asymex_fib_gsymbol_similarity, thm:asymex_fib_anomaly_fusion_algebra,
        def:mpo_fusion_algebra}
    Extend $B_1$ and $B_\tau$ to the four physical letters (plaquette, edge
    label) by keeping $B_f^{x'x}$ at edge label $\tau$, putting the diagonal
    matrix unit $E_{pp}$ at the bond letter $p=(x',x)$ at edge label $1$
    (zero when $(x',x)$ is not a bond letter of the block $f$), and zero when
    the edge label changes. Then $h\,S_f^{ij}\,h^{-1}=B_f^{ij}$ for both labels
    $f$ and all physical letters $i,j$, and the periodic operators of the
    source's blocks $S_1$, $S_\tau$ on $(\C^4)^{\otimes N}$ satisfy the
    Fibonacci fusion rules $O_1^2=O_1$, $O_1O_\tau=O_\tau O_1=O_\tau$ and
    $O_\tau^2=O_1+O_\tau$ at every positive length. This is the fusion
    statement of~\cite[Appendix~D.1]{Bultinck2017Anyons}.
\end{theorem}

\begin{proof}
    \leanok
    \uses{thm:asymex_golden_ring, thm:asymex_fib_anomaly_fusion_algebra}
    At edge label $1$ the selection rule of $G^{abc}_{d1f}$ forces $a=b$ and
    $c=d$, and $G^{bbc}_{c1f}v_bv_c=1$ whenever $(b,c)$ is a bond letter of
    the block $f$. So the letter of $S_f$ at $\bigl((b,1),(c,1)\bigr)$ is the
    diagonal matrix unit $E_{pp}$ at $p=(b,c)$, which the diagonal matrix $h$
    leaves unchanged; together with
    Theorem~\ref{thm:asymex_fib_gsymbol_similarity} this gives the similarity
    on all letters, and the periodic operators of $S_f$ and of the extended
    $B_f$ agree. For each product of two extended blocks, the gauge of the
    corresponding compression of
    Theorem~\ref{thm:asymex_fib_anomaly_fusion_algebra} makes every one of the
    sixteen letters of the stacked tensor block upper triangular: at edge
    label $\tau$ the conjugated letters are the block-diagonal ones already
    recorded, and at edge label $1$ they carry the edge-$1$ letters of the
    target blocks on the diagonal and further entries only in the columns of
    the zero slots. These are finite computations in $\Z[\sigma]$. The trace
    identity of a compression holds for a nonzero nilpotent remainder, which
    gives the four fusion rules.
\end{proof}

\begin{theorem}[The prefactor of the review]
    \label{thm:asymex_fib_review_congruence}
    % Principal declaration: FibonacciCompression.mpo_fibReviewTensor.
    \lean{FibonacciCompression.mpo_fibReviewTensor,
        FibonacciCompression.fibReviewTensor_conj,
        FibonacciCompression.not_isMPOFusionAlgebra_fibReviewTensor,
        FibonacciCompression.isMPOFusionAlgebra_fibReviewWeightedTensor,
        MPOTensor.mpo_letterwise_smul}
    \leanok
    \uses{def:asymex_fib_gsymbol_tensor, thm:asymex_fib_anomaly_fusion_algebra,
        def:mpo_fusion_algebra}
    For both labels $f$ and every length $N$, the periodic operator of the
    tensor (\ref{eq:asymex_fib_review_tensor}) is the congruence
    \begin{align}
        O_N(R_f) & =W^{-1/2}\,O_N(B_f)\,W^{-1/2}.
        \label{eq:asymex_fib_review_congruence}
    \end{align}
    These operators do not satisfy the Fibonacci fusion rules: at $N=1$,
    $O_N(R_1)=\operatorname{diag}(0,1/\varphi)$, whose square has the entry
    $1/\varphi^2\neq1/\varphi$. The operators
    $O_N(\widetilde R_f)=O_N(R_f)W=W^{-1/2}O_N(B_f)W^{1/2}$
    satisfy them at every positive length.
\end{theorem}

\begin{proof}
    \leanok
    \uses{thm:asymex_fib_anomaly_fusion_algebra}
    The diagonal bond matrix $g(u,l)=d_l^{-1/2}$ carries $R_f^{x'x}$ to
    $(d_{x'}d_x)^{-1/2}B_f^{x'x}$, a multiple of each letter by its own
    scalar. The periodic operator of such a tensor is the operator of $B_f$
    with its entry at $(x,y)$ multiplied by $\prod_k(d_{x_k}d_{y_k})^{-1/2}$,
    which is (\ref{eq:asymex_fib_review_congruence}). At $N=1$,
    $O_N(B_1)=\operatorname{diag}(0,1)$ gives the stated entry. Finally,
    $O_N(\widetilde R_f)=W^{-1/2}O_N(B_f)W^{1/2}$ is similar to
    $O_N(B_f)$ by one matrix for all labels, so products and linear
    combinations are carried along and the fusion rules of
    Theorem~\ref{thm:asymex_fib_anomaly_fusion_algebra} transfer.
\end{proof}

While the review states $O_aO_b=\sum_cN_{ab}^cO_c$ for the prefactored
review tensor $R_f$~\cite[Section~3]{Cirac2021Matrix}, that operator already
contains the prefactor $1/\sqrt{d_{x'}d_{q''}}$ and fails the fusion rules
in the orthonormal plaquette basis.
Theorem~\ref{thm:asymex_fib_review_congruence} shows that the fusion rules
hold after composing it with $W$;
see~\cite{gap:bmwshv17_fibonacci_block_entries_provenance}.

\begin{theorem}[Normality of the two Fibonacci blocks]
    \label{thm:asymex_fib_normal}
    \lean{FibonacciCompression.fibOneMPS_isNormal,
        FibonacciCompression.fibTauMPS_isNormal}
    \leanok
    \uses{def:asymex_fib, def:normal}
    Both blocks $B_1$ and $B_\tau$ are normal: their length-three words
    satisfy $\spn_{\C}\{B_1^w : |w|=3\} = \MN{2}$ and
    $\spn_{\C}\{B_\tau^w : |w|=3\} = \MN{3}$.
\end{theorem}

\begin{proof}
    \leanok
    \uses{thm:asymex_golden_ring}
    For the $\tau$ block, let $e_i$ be the $i$-th coordinate column
    vector and $f_i$ the $i$-th row of
    (\ref{eq:asymex_fib_vertex}). In this proof, superscripts
    $0,1,2,3$ enumerate the ordered pair alphabet. Thus
    $B_\tau^{i+1}=e_i f_i$ for $i\in\{0,1,2\}$, and
    \begin{align}
        B_\tau^{i+1}B_\tau^3B_\tau^{k+1}
         & =F_{i2}F_{2k}\,e_i f_k.
    \end{align}
    All factors $F_{i2}F_{2k}$ are nonzero, and the rows $f_k$
    form a basis. Each of the nine matrix units is an explicit
    combination over $\Z[\sigma]$ of these nine words.
    For the vacuum block, four length-three words suffice.
    Both families of coefficients are verified in $\Z[\sigma]$
    and mapped to complex matrices by
    Theorem~\ref{thm:asymex_golden_ring}.
\end{proof}

\subsection{Admissible configurations and the vacuum operator}

Define the adjacency matrix $K=\begin{pmatrix}0&1\\1&1\end{pmatrix}$
and, for a cyclic word $x\in\{0,1\}^N$, set
\begin{align}
    p_N(x) & =\prod_{v=0}^{N-1}K_{x_vx_{v+1}}.
    \label{eq:asymex_fib_admissibility}
\end{align}
A cyclic word $x$ is admissible if $p_N(x)=1$, equivalently if it contains no
cyclic neighboring pair $00$.

\begin{lemma}[Vacuum projection and its kernel]
    \label{lem:asymex_fib_vacuum}
    \lean{FibonacciCompression.mpo_fibOne_apply,
        FibonacciCompression.mpo_fibOne_isStarProjection,
        FibonacciCompression.range_mpo_fibOne,
        FibonacciCompression.ker_mpo_fibOne,
        FibonacciCompression.mpo_fibOne_ne_one,
        FibonacciCompression.fibAdjacency,
        FibonacciCompression.fibAdmissibility,
        FibonacciCompression.fibAdmissibility_eq_one_iff}
    \leanok
    \uses{def:asymex_fib, def:mpo_family}
    For $N\geq1$, the vacuum operator $P_N=O_N(B_1)$ has matrix elements
    \begin{align}
        \bra{y}P_N\ket{x} & =\delta_{yx}p_N(x).
        \label{eq:asymex_fib_vacuum_kernel}
    \end{align}
    Thus $P_N\ket{x}=p_N(x)\ket{x}$. It is the orthogonal projection
    onto the admissible subspace:
    \begin{align}
        \ran P_N & =\spn_{\C}\{\ket{x}:p_N(x)=1\}.
        \label{eq:asymex_fib_vacuum_range}
    \end{align}
    \begin{align}
        \ker P_N & =\spn_{\C}\{\ket{x}:p_N(x)=0\}.
        \label{eq:asymex_fib_vacuum_null}
    \end{align}
    In particular, $P_N^2=P_N=P_N^\dagger$, but $P_N\neq\Id_{\mc H_N}$.
\end{lemma}

\begin{proof}
    \leanok
    The vacuum letters satisfy
    $(B_1^{yx})_{ab}=\delta_{yx}\delta_{ax}K_{xb}$.
    Expanding the periodic trace over bond labels
    $g=(g_0,\ldots,g_{N-1})$, with $g_N=g_0$, gives
    \begin{align}
        \bra{y}P_N\ket{x}
         & =\sum_g\prod_{v=0}^{N-1}
        \delta_{y_vx_v}\delta_{g_vx_v}K_{x_vg_{v+1}}
        =\delta_{yx}\prod_{v=0}^{N-1}K_{x_vx_{v+1}}.
    \end{align}
    Only $g=x$ contributes. Since $p_N(x)\in\{0,1\}$,
    the operator is diagonal with real idempotent entries.
    This proves the projection, range, and kernel statements.
    Since $\ket{0,\ldots,0}\in\ker P_N$, $P_N\neq\Id_{\mc H_N}$,
    with $P_1=\diag(0,1)$ at one site.
\end{proof}

\begin{theorem}[Vacuum unit laws]
    \label{thm:asymex_fib_unit}
    \lean{FibonacciCompression.fibOne_mul_fibOne,
        FibonacciCompression.fibOne_mul_fibTau,
        FibonacciCompression.fibTau_mul_fibOne}
    \leanok
    \uses{def:asymex_fib, def:mpo_family}
    For every $N\geq1$, the periodic operators satisfy
    \begin{align}
        O_N(B_1)^2 & =O_N(B_1).
        \label{eq:asymex_fib_unit_square}
    \end{align}
    \begin{align}
        O_N(B_1)O_N(B_\tau) & =O_N(B_\tau).
        \label{eq:asymex_fib_unit_left}
    \end{align}
    \begin{align}
        O_N(B_\tau)O_N(B_1) & =O_N(B_\tau).
        \label{eq:asymex_fib_unit_right}
    \end{align}
\end{theorem}

\begin{proof}
    \leanok
    For MPO tensors $T$ and $U$, denote their stacked product by
    $(TU)^{yx}=\sum_mT^{ym}\otimes U^{mx}$.
    Write $0_r$ for the tensor whose letters are zero $r\times r$
    matrices, and $\sim$ for simultaneous similarity of all letters.
    The three stacked products admit the decompositions
    \begin{align}
        B_1B_1 & \sim B_1\oplus0_2.
        \label{eq:asymex_fib_unit_blocks}
    \end{align}
    \begin{align}
        B_1B_\tau & \sim B_\tau\oplus0_3.
    \end{align}
    \begin{align}
        B_\tau B_1 & \sim B_\tau\oplus0_3.
    \end{align}
    The changes of bond coordinates and their inverses have entries
    in $\Z[\sigma]$. The decompositions, of dimensions $4=2+2$ and
    $6=3+3$, are verified by matrix multiplication with coefficients
    reduced using $\sigma^4=1-\sigma^2$.
    The periodic product formula for a multi-block compression gives
    the trace of the single nonzero target block. Since $N>0$, every
    zero block contributes zero, proving the three identities.
\end{proof}

By~(\ref{eq:asymex_fib_unit_left}), the range of $O_N(B_\tau)$ lies in
the admissible subspace. By~(\ref{eq:asymex_fib_unit_right}), it vanishes
on the orthogonal complement of that subspace. Thus the vacuum acts
as the identity on the admissible subspace and is the unit of the
algebra generated there by $O_N(B_\tau)$.

This vacuum block is distinct from the weighted string-net projector
of~\cite[Appendix~D.1]{Bultinck2017Anyons}, which combines the vacuum and
$\tau$ blocks with weights $w_1=1/(1+\varphi^2)$ and
$w_\tau=\varphi/(1+\varphi^2)$.

\begin{definition}[The weighted Fibonacci projector]
    \label{def:asymex_fib_projector}
    \lean{FibonacciCompression.fibWeight, FibonacciCompression.fibPMPO,
        FibonacciCompression.fibDelta, FibonacciCompression.fibProjector}
    \leanok
    \uses{def:asymex_fib, def:mpo_direct_sum, def:mpo_family}
    The weights of the two labels are
    $w_1=1/(1+\varphi^2)$ and $w_\tau=\varphi/(1+\varphi^2)$, the quantum
    dimensions divided by the square of the total quantum dimension.
    The bond-dimension-five tensor is the direct sum $B_1\oplus B_\tau$,
    and its boundary matrix is
    $\Delta=w_1\Id_2\oplus w_\tau\Id_3$. At every system size $N$, the
    weighted projector is
    \begin{align}
        \Pi_N & =w_1O_N(B_1)+w_\tau O_N(B_\tau)
        =\frac{O_N(B_1)+\varphi\,O_N(B_\tau)}{1+\varphi^2},
        \label{eq:asymex_fib_projector}
    \end{align}
    following~\cite[Sections~3.1 and~3.2 and Appendix~D.1]{Bultinck2017Anyons}.
\end{definition}

\begin{theorem}[The weight relation]
    \label{thm:asymex_fib_projector_weights}
    \lean{FibonacciCompression.fibWeight_fusion}
    \leanok
    \uses{def:asymex_fib_projector, def:asymex_fib_anomaly_nim}
    For each label $c\in\{1,\tau\}$,
    $\sum_{a,b\in\{1,\tau\}}N_{ab}^cw_aw_b=w_c$.
\end{theorem}

\begin{proof}
    \leanok
    For $c=1$ the relation reads $w_1=w_1^2+w_\tau^2$, and both
    sides equal $1/(1+\varphi^2)$. For $c=\tau$ it reads
    $w_\tau=2w_1w_\tau+w_\tau^2$; after multiplication by
    $(1+\varphi^2)^2/\varphi$ it becomes $1+\varphi^2=2+\varphi$, a
    consequence of $\varphi^2=\varphi+1$.
\end{proof}

\begin{theorem}[The projector from the bond-dimension-five tensor]
    \label{thm:asymex_fib_projector_trace}
    % Principal declaration: FibonacciCompression.fibProjector_eq_trace_fibDelta.
    \lean{FibonacciCompression.fibProjector_eq_trace_fibDelta,
        FibonacciCompression.mpo_fibPMPO}
    \leanok
    \uses{def:asymex_fib_projector}
    At every system size $N$, write $B=B_1\oplus B_\tau$. Then
    \begin{align}
        \bra{y}\Pi_N\ket{x}
         & =\tr(\Delta B^{y_0x_0}\cdots B^{y_{N-1}x_{N-1}}),
        \label{eq:asymex_fib_projector_trace}
    \end{align}
    and without the boundary matrix the periodic operator is
    $O_N(B)=O_N(B_1)+O_N(B_\tau)$.
\end{theorem}

\begin{proof}
    \leanok
    \uses{lem:mpo_direct_sum_eval_word, thm:mpo_direct_sum_additive}
    Word evaluation of the direct sum is block diagonal by
    Lemma~\ref{lem:mpo_direct_sum_eval_word}. Multiplication by the
    block-diagonal matrix $\Delta$ scales the two blocks by $w_1$ and
    $w_\tau$, and the trace of a block-diagonal matrix is the sum of the
    traces of its blocks. The second statement is
    Theorem~\ref{thm:mpo_direct_sum_additive}.
\end{proof}

\begin{theorem}[The weighted Fibonacci projector is idempotent]
    \label{thm:asymex_fib_projector_idempotent}
    \lean{FibonacciCompression.fibProjector_mul_self}
    \leanok
    \uses{def:asymex_fib_projector}
    For every system size $N\geq1$,
    $(\Pi_N)^2=\Pi_N$.
\end{theorem}

\begin{proof}
    \leanok
    \uses{thm:asymex_fib_unit, thm:asymex_fib_fusion,
        thm:asymex_fib_projector_weights}
    Expanding the square and using the unit laws of
    Theorem~\ref{thm:asymex_fib_unit} and the fusion rule of
    Theorem~\ref{thm:asymex_fib_fusion} gives
    $(w_1^2+w_\tau^2)O_N(B_1)+(2w_1w_\tau+w_\tau^2)O_N(B_\tau)$.
    Theorem~\ref{thm:asymex_fib_projector_weights} identifies the two
    coefficients with $w_1$ and $w_\tau$.
\end{proof}

At $N=0$ the empty traces give $O_0(B_1)=2$ and $O_0(B_\tau)=3$, so
$\Pi_0=(2+3\varphi)/(1+\varphi^2)$, a scalar different from $0$ and
$1$, hence $\Pi_0^2\neq\Pi_0$;
\cite[Section~3.2]{Bultinck2017Anyons} asks for idempotence at every
length, and the empty chain is excluded as a degenerate reading;
see~\cite{gap:bmwshv17_fibonacci_projector_positive_length}.

\subsection{The golden chain}

The golden chain of Feiguin et al.~\cite{Feiguin2007GoldenChain} acts on
the labels $x_0,\ldots,x_{N-1}\in\{1,\tau\}$ of a periodic chain. Its
three-site term projects the two anyons at links $i$ and $i+1$ onto their
trivial fusion channel, with the amplitudes
\begin{align}
    (p_i)^{x_i'}_{x_i}
     & =\left(F^{x_{i-1}\tau\tau}_{x_{i+1}}\right)_{x_i'}^{1}
    \left(F^{x_{i-1}\tau\tau}_{x_{i+1}}\right)_{x_i}^{1},
    \label{eq:asymex_fib_golden_term}
\end{align}
where $p_i$ changes only the label $x_i$.
The Hamiltonian is $H_N=-\sum_ip_i$. The review states that such anyonic
chains commute with the whole algebra of periodic
operators~\cite{Cirac2021Matrix}.

\begin{definition}[The golden chain]%
    \label{def:asymex_fib_golden_chain}
    \lean{FibonacciCompression.goldenChainTermGolden,
        FibonacciCompression.goldenChainTerm,
        FibonacciCompression.goldenChainHamiltonian,
        FibonacciCompression.fibYGolden,
        FibonacciCompression.fibOneKernelGolden}
    \leanok
    \uses{def:asymex_fib_fsymbols}
    The golden chain is $H_N=-\sum_ip_i$ with the local terms
    of~(\ref{eq:asymex_fib_golden_term}). The topological symmetry has the
    factors $\left(F^{\tau x_n\tau}_{x_{n+1}'}\right)_{x_n'}^{x_{n+1}}$, and the
    vacuum factor is $\delta_{x_n'x_n}\delta_{x_{n+1}'x_{n+1}}$ unless
    $x_n=x_{n+1}=1$, where it vanishes.
\end{definition}

\begin{theorem}[Values of the golden-chain projector]%
    \label{thm:asymex_fib_golden_term_values}
    \lean{FibonacciCompression.goldenChainTermGolden_tau_tau,
        FibonacciCompression.goldenChainTermGolden_of_ne_tau_tau}
    \leanok
    \uses{def:asymex_fib_golden_chain}
    The local term $p_i$ is the projector onto $\ket{1\tau1}$ between two
    trivial neighbours, it vanishes between a trivial and a $\tau$ neighbour,
    and between two $\tau$ neighbours it is
    \begin{align}
        \begin{pmatrix}
            \varphi^{-2}   & \varphi^{-3/2} \\
            \varphi^{-3/2} & \varphi^{-1}
        \end{pmatrix}
        \notag
    \end{align}
    in the basis $1,\tau$ of~$x_i$.
\end{theorem}

\begin{proof}
    \leanok
    \uses{thm:asymex_fib_fsymbol_values}
    The entries are products of two F-symbols and are checked entrywise in
    $\Z[\sigma]$, where $\varphi^{-1}=\sigma^2$.
\end{proof}

\begin{theorem}[Site operators and nearest-neighbour kernels]%
    \label{thm:mpo_site_operator_neighbour_kernel}
    \lean{MPOTensor.neighbourKernel, MPOTensor.siteOperator,
        MPOTensor.sum_ite_agree, MPOTensor.siteOperator_mul_neighbourKernel}
    \leanok
    On a periodic chain of $N\geq2$ sites, let $K$ have the entries
    $\bra{s}K\ket{t}=\prod_nf(s_n,s_{n+1},t_n,t_{n+1})$, and let $p_i$ change
    only the state at site $i$, with amplitude $h(t_{i-1},t_{i+1},s_i,t_i)$.
    If
    \begin{align}
        \sum_yh(a',c',b',y)f(a',y,a,b)f(y,c',b,c)
         & =\sum_yf(a',b',a,y)f(b',c',y,c)h(a,c,y,b)
        \notag
    \end{align}
    for all labels, then $p_iK=Kp_i$.
\end{theorem}

\begin{proof}
    \leanok
    Write $R(s,t)=\prod_{n\neq i-1,i}f(s_n,s_{n+1},t_n,t_{n+1})$ for the
    factors of $K$ away from the two bonds at site $i$; it does not depend
    on $s_i$ or $t_i$. Summing over the label $y$ at site $i$ of the
    intermediate configuration,
    \begin{align}
        \bra{s}p_iK\ket{t}
         & =R(s,t)\sum_yh(s_{i-1},s_{i+1},s_i,y)
        f(s_{i-1},y,t_{i-1},t_i)f(y,s_{i+1},t_i,t_{i+1}),
        \notag                                   \\
        \bra{s}Kp_i\ket{t}
         & =R(s,t)\sum_yf(s_{i-1},s_i,t_{i-1},y)
        f(s_i,s_{i+1},y,t_{i+1})h(t_{i-1},t_{i+1},y,t_i).
        \notag
    \end{align}
    The hypothesis with $(a',b',c')=(s_{i-1},s_i,s_{i+1})$ and
    $(a,b,c)=(t_{i-1},t_i,t_{i+1})$ equates the two sums.
\end{proof}

\begin{theorem}[The golden chain commutes with the Fibonacci operators]%
    \label{thm:asymex_fib_golden_chain}
    \lean{FibonacciCompression.mpo_fibTau_eq_neighbourKernel,
        FibonacciCompression.goldenChainTerm_local_tau,
        FibonacciCompression.mpo_fibOne_eq_neighbourKernel,
        FibonacciCompression.goldenChainTerm_local_one,
        FibonacciCompression.goldenChainHamiltonian_commute_mpo_fibTau,
        FibonacciCompression.goldenChainHamiltonian_commute_mpo_fibOne,
        FibonacciCompression.goldenChainHamiltonian_commute_fibProjector,
        FibonacciCompression.goldenChainHamiltonian_commute_mpo_fibStringNetEdgeTau}
    \leanok
    \uses{def:asymex_fib, def:asymex_fib_golden_chain, def:asymex_fib_projector,
        def:asymex_fib_gsymbol_tensor}
    For every $N\geq2$, the golden chain $H_N$ commutes with $O_N(B_\tau)$,
    with $O_N(B_1)$ and with the projector $\Pi_N$, and, for each label $f$,
    with the periodic operator of the operator tensor $S_f$ restricted to
    configurations with edge label $\tau$.
\end{theorem}

\begin{proof}
    \leanok
    \uses{thm:mpo_site_operator_neighbour_kernel,
        thm:asymex_fib_gsymbol_similarity}
    In the $\tau$ block the bond index at site $n$ is determined by the
    pair $(x_n',x_n)$, so the only closed bond configuration that survives
    gives
    $\bra{x'}O_N(B_\tau)\ket{x}=\prod_n
        \left(F^{\tau x_n\tau}_{x_{n+1}'}\right)_{x_n'}^{x_{n+1}}$,
    the topological symmetry of~\cite{Feiguin2007GoldenChain}; similarly
    $O_N(B_1)$ is diagonal with entry $1$ exactly on the configurations
    without two neighbouring trivial labels. In both cases the entry is a
    product of factors $f(x_n',x_{n+1}',x_n,x_{n+1})$ around the ring. The
    product $p_iO_N(B_\tau)$ and the product $O_N(B_\tau)p_i$ share the
    factors away from the two bonds at site $i$, and the remaining sums
    over the intermediate label at site $i$ agree by a three-site identity,
    checked on the sixty-four choices of the six outer and inner labels in
    exact golden arithmetic. The same identity holds for the kernel of
    $O_N(B_1)$. Summing over $i$ gives the first two statements, and the
    third follows since $\Pi_N=w_1O_N(B_1)+w_\tau O_N(B_\tau)$. On edge label
    $\tau$ the periodic operator of $S_f$ equals $O_N(B_f)$ by
    Theorem~\ref{thm:asymex_fib_gsymbol_similarity}, which gives the last
    statement.
\end{proof}

\subsection{Two periodic states and the Fibonacci action}

For an MPS tensor $D$, write
$\ket{V^{(N)}(D)}=\sum_x\tr(D^{x_0}\cdots D^{x_{N-1}})\ket{x}$.
These vectors are not normalized unless this is stated explicitly.

\begin{definition}[The all-$\tau$ and chain tensors]
    \label{def:asymex_fib_states}
    \lean{FibonacciCompression.fibAllTau,
        FibonacciCompression.fibChain}
    \leanok
    The all-$\tau$ tensor $A$ has bond dimension one, with
    $A^0=(0)$ and
    \begin{align}
        A^1 & =(1).
        \label{eq:asymex_fib_product_tensor}
    \end{align}
    The chain tensor $C$ has bond dimension two:
    \begin{align}
        C^0 & =\begin{pmatrix}0&\sigma\\0&0\end{pmatrix}.
        \label{eq:asymex_fib_chain_zero}
    \end{align}
    \begin{align}
        C^1 & =\begin{pmatrix}0&0\\1&-\sigma^2\end{pmatrix}.
        \label{eq:asymex_fib_chain_one}
    \end{align}
\end{definition}

\begin{lemma}[Amplitudes and norms]
    \label{lem:asymex_fib_amplitudes}
    \lean{FibonacciCompression.mpvState_fibAllTau,
        FibonacciCompression.norm_mpvState_fibAllTau,
        FibonacciCompression.mpv_fibChain,
        FibonacciCompression.norm_mpvState_fibChain_sq,
        FibonacciCompression.norm_mpvState_fibChain_pos}
    \leanok
    \uses{def:asymex_fib_states}
    Let $N\geq1$. The product tensor gives
    $\ket{V^{(N)}(A)}=\ket{1,\ldots,1}$, of norm one.
    For a word $x$, let $n_0(x)$ count its zeros and let $n_{11}(x)$
    count its cyclic neighboring pairs $11$. Then
    \begin{align}
        V^{(N)}(C)_x
         & =p_N(x)\sigma^{n_0(x)}(-\sigma^2)^{n_{11}(x)}.
        \label{eq:asymex_fib_chain_amplitude}
    \end{align}
    Both vectors are fixed by $P_N$. The chain norm is
    \begin{align}
        \bigl\|\ket{V^{(N)}(C)}\bigr\|^2
         & =\sum_{x\in\{0,1\}^N}
        p_N(x)\sigma^{2n_0(x)+4n_{11}(x)}.
        \label{eq:asymex_fib_chain_norm}
    \end{align}
    This norm is positive.
\end{lemma}

\begin{proof}
    \leanok
    \uses{thm:asymex_fib_state_action, thm:asymex_fib_unit,
        thm:asymex_fib_anomaly_nim_rep}
    The bond-one trace is $\prod_v\delta_{x_v1}$.
    For the chain, put
    $H=\begin{pmatrix}0&\sigma\\1&-\sigma^2\end{pmatrix}$.
    Since $(C^x)_{ab}=\delta_{ax}H_{xb}$, the cyclic contraction gives
    \begin{align}
        V^{(N)}(C)_x
         & =\sum_g\prod_v\delta_{g_vx_v}H_{x_vg_{v+1}}
        =\prod_vH_{x_vx_{v+1}}.
    \end{align}
    A pair $00$ makes the product zero. In an admissible word, every
    zero is followed by a one, giving $n_0(x)$ factors of $\sigma$.
    Each pair $10$ contributes $1$, and each pair $11$ contributes
    $-\sigma^2$. This proves~(\ref{eq:asymex_fib_chain_amplitude}).
    The operator $P_N=O_N(B_1)$ fixes the product vector by
    Theorem~\ref{thm:asymex_fib_state_action}. The chain vector is
    $O_N(B_\tau)\ket{V^{(N)}(A)}$ by the same theorem, and the unit law
    $O_N(B_1)O_N(B_\tau)=O_N(B_\tau)$ of
    Theorem~\ref{thm:asymex_fib_unit} shows that $P_N$ fixes it
    (Theorem~\ref{thm:asymex_fib_anomaly_nim_rep}).
    Summing squared amplitudes proves~(\ref{eq:asymex_fib_chain_norm}).
    The all-one amplitude is $(-\sigma^2)^N\neq0$.
\end{proof}

The normalized chain vector is obtained by dividing by the positive
square root of~(\ref{eq:asymex_fib_chain_norm}). For example,
$\ket{V^{(1)}(C)}=-\sigma^2\ket{1}$ and its squared norm is $\sigma^4$.
In particular, the two periodic vectors are not independent at every
length.

\begin{theorem}[Normality of the two state tensors]
    \label{thm:asymex_fib_states_normal}
    \lean{FibonacciCompression.fibAllTau_isNormal,
        FibonacciCompression.fibChain_isNormal}
    \leanok
    \uses{def:asymex_fib_states, def:normal}
    The length-one words of $A$ span $\MN{1}$, and the length-three
    words of $C$ span $\MN{2}$. Thus both tensors are normal.
\end{theorem}

\begin{proof}
    \leanok
    The letter $A^1$ spans $\MN{1}$. For $C$, write $E_{ab}$ for
    the matrix units of $\MN{2}$. Direct multiplication gives
    \begin{align}
        C^0C^1C^0 & =\sigma^2E_{01}.
    \end{align}
    \begin{align}
        C^0C^1C^1 & =-\sigma^3E_{00}+\sigma^5E_{01}.
    \end{align}
    \begin{align}
        C^1C^0C^1 & =\sigma C^1.
    \end{align}
    \begin{align}
        C^1C^1C^0 & =-\sigma^3E_{11}.
    \end{align}
    These four words span all four matrix units, since $\sigma\neq0$.
    The coefficients expressing the units in their span can all be
    chosen in $\Z[\sigma]$, using
    $\sigma^{-1}=\sigma+\sigma^3$.
    Embedding these identities into $\C$ proves that the words span the
    matrix algebra.
\end{proof}

\begin{theorem}[Action on the two periodic states]
    \label{thm:asymex_fib_state_action}
    \lean{FibonacciCompression.mpo_fibOne_mulVec_allTau,
        FibonacciCompression.mpo_fibTau_mulVec_allTau,
        FibonacciCompression.mpo_fibTau_mulVec_chain}
    \leanok
    \uses{def:asymex_fib, def:asymex_fib_states, def:mpo_family}
    For every $N\geq1$,
    \begin{align}
        O_N(B_1)\ket{V^{(N)}(A)}
         & =\ket{V^{(N)}(A)}.
        \label{eq:asymex_fib_action_unit}
    \end{align}
    \begin{align}
        O_N(B_\tau)\ket{V^{(N)}(A)}
         & =\ket{V^{(N)}(C)}.
        \label{eq:asymex_fib_action_product}
    \end{align}
    \begin{align}
        O_N(B_\tau)\ket{V^{(N)}(C)}
         & =\ket{V^{(N)}(A)}+\ket{V^{(N)}(C)}.
        \label{eq:asymex_fib_action_chain}
    \end{align}
\end{theorem}

\begin{proof}
    \leanok
    For an MPO tensor $T$ and MPS tensor $D$, their action tensor
    has letters $(T\cdot D)^y=\sum_xT^{yx}\otimes D^x$.
    As in the proof of Theorem~\ref{thm:asymex_fib_unit}, let $\sim$
    denote simultaneous similarity of letters and let $0_r$ denote
    a zero tensor of bond dimension $r$.
    The three action tensors have block decompositions
    \begin{align}
        B_1\cdot A & \sim A\oplus0_1.
        \label{eq:asymex_fib_action_blocks_unit}
    \end{align}
    \begin{align}
        B_\tau\cdot A & \sim C\oplus0_1.
        \label{eq:asymex_fib_action_blocks_product}
    \end{align}
    \begin{align}
        B_\tau\cdot C & \sim A\oplus C\oplus0_3.
        \label{eq:asymex_fib_action_blocks_chain}
    \end{align}
    In lexicographic order on the product bond indices, matrices whose
    columns give these adapted bases are respectively
    \begin{align}
        Q_2 & =\begin{pmatrix}0&1\\1&-1\end{pmatrix}.
    \end{align}
    \begin{align}
        Q_3 & =\begin{pmatrix}1&0&1\\0&0&-\sigma\\0&1&\sigma^2\end{pmatrix}.
    \end{align}
    \begin{align}
        Q_6 & =\begin{pmatrix}
                   0               & 0 & 0         & 0          & 1        & 1         \\
                   0               & 1 & 0         & 0          & 0        & \sigma^2  \\
                   1               & 0 & \sigma    & 1          & -\sigma  & -\sigma^3 \\
                   0               & 0 & 0         & 1+\sigma^2 & 0        & 0         \\
                   0               & 0 & 0         & 0          & \sigma^2 & 0         \\
                   \sigma+\sigma^3 & 0 & -\sigma^2 & 0          & 0        & -\sigma^2
               \end{pmatrix}.
    \end{align}
    They are invertible, with inverses over $\Z[\sigma]$.
    For each action tensor, let $D_{\mathrm{blk}}$ denote the
    corresponding block tensor in
    (\ref{eq:asymex_fib_action_blocks_unit}),
    (\ref{eq:asymex_fib_action_blocks_product}), or
    (\ref{eq:asymex_fib_action_blocks_chain}).
    Matrix multiplication verifies
    $(T\cdot D)^yQ=QD_{\mathrm{blk}}^y$.
    All off-diagonal blocks vanish, giving the three multi-block
    compressions. The periodic action formula for each compression
    sums its target vectors. The zero blocks contribute nothing
    because $N>0$. The first two sums have one target each, and the
    third has the targets $A$ and $C$.
\end{proof}

The first nontrivial action also has a direct contraction proof.
The all-one input selects the letters $B_\tau^{01}$ and
$B_\tau^{11}$. Both have zero row $1$, so a cyclic bond configuration
using that coordinate contributes zero. Restricting to the ordered
coordinates $(0,2)$ gives exactly $C^0$ and $C^1$. Hence
\begin{align}
    \bra{x}O_N(B_\tau)\ket{1,\ldots,1}
     & =\tr(B_\tau^{x_0 1}\cdots B_\tau^{x_{N-1}1})
    =\tr(C^{x_0}\cdots C^{x_{N-1}}).
\end{align}
There is no additional normalization factor.

Relative to the ordered pair of these unnormalized state families,
the action coefficients are the Fibonacci matrix
$N_\tau=\begin{pmatrix}0&1\\1&1\end{pmatrix}$, realizing the nonnegative-integer
action discussed in~\cite{GarreRubio2023MPOSym}.
The positive-length hypothesis is essential: at $N=0$ the periodic trace is the
bond dimension, yielding $O_0(B_1)=2$, $O_0(B_\tau)=3$, and empty amplitudes
$1$ and $2$ for $A$ and $C$.

\subsection{Asymmetric compression of the stacked square}

The stacked square in~(\ref{eq:asymex_fib_stack}) contracts the physical
index between two copies of $B_\tau$ while retaining both virtual
indices. Its periodic operator is therefore the square of
$O_N(B_\tau)$. The compression maps onto its two target blocks also
have entries in $\Z[\sigma]$, so the fusion rule below is again an
exact computation in the ring of golden integers.

\begin{center}
    % Formula: B^{x'x}=\sum_m(B_\tau)^{x'm}\otimes(B_\tau)^{mx}, the
    % stacked product of Definition~\ref{def:asymex_fib}, realizing
    % $\tau\otimes\tau$; Theorem~\ref{thm:asymex_fib_fusion} identifies
    % its periodic trace with that of $B_1\oplus B_\tau$, the
    % realization $1\oplus\tau$.
    % Source: Definition~\ref{def:asymex_fib}.
    % Ink-to-index map: both rows are the $\tau$ tensor $B_\tau$; the
    % shared vertical pairleg at each site is the summed intermediate
    % label $m$; the drawn upper-row label $y$ is the per-site
    % component of the outgoing index $x'$.
    % Contraction set: the vertical pairlegs (index $m$) at every site
    % and the horizontal virtual bonds within each row; the physical
    % legs $(y,x)$, the per-site components of $(x',x)$, remain open.
    % Boundary signature: (virtual-west=0 | virtual-east=0 |
    % physical-up=3 | physical-down=3) on both sides.
    \tenkzkernel
    \begin{tenkz}[rows={op,op}, cols=4, boundary=periodic]
        \tn[skin=mpo, ports={90:physical:$y_1$, 270:physical}]{B_\tau} &
        \tn[skin=mpo, ports={90:physical:$y_2$, 270:physical}]{B_\tau} &
        \tn[skin=dots, ports={180:virtual, 0:virtual}]{} &
        \tn[skin=mpo, ports={90:physical:$y_N$, 270:physical}]{B_\tau} \\
        \tn[skin=mpo, ports={90:physical, 270:physical:$x_1$}]{B_\tau} &
        \tn[skin=mpo, ports={90:physical, 270:physical:$x_2$}]{B_\tau} &
        \tn[skin=dots, ports={180:virtual, 0:virtual}]{} &
        \tn[skin=mpo, ports={90:physical, 270:physical:$x_N$}]{B_\tau}
    \end{tenkz}
    \quad $=$ \quad
    \begin{tenkz}[rows={op}, cols=4, boundary=periodic]
        \tn[skin=mpo, ports={90:physical:$y_1$, 270:physical:$x_1$}]{B} &
        \tn[skin=mpo, ports={90:physical:$y_2$, 270:physical:$x_2$}]{B} &
        \tn[skin=dots, ports={180:virtual, 0:virtual}]{} &
        \tn[skin=mpo, ports={90:physical:$y_N$, 270:physical:$x_N$}]{B}
    \end{tenkz}
\end{center}

\begin{theorem}[Split compression of the Fibonacci square]
    \label{thm:asymex_fib_compression}
    % Principal declaration: FibonacciCompression.fibonacci_dim_eq.
    \lean{FibonacciCompression.fibonacci_compression, FibonacciCompression.fibonacci_dim_eq,
        FibonacciCompression.fibonacci_evalWord_remainder_eq_zero,
        FibonacciCompression.fibonacci_remainder_eq_zero}
    \leanok
    \uses{def:asymex_fib, thm:asym_multiblock}
    The source of Definition~\ref{def:asymex_fib} carries the block
    triangular form of Theorem~\ref{thm:asym_multiblock}, with target
    blocks $C_1=B_1$ and $C_\tau=B_\tau$, and $z=4$ zero slots.
    The dimension count is $9=2+3+4$. Here the indexed tensors $C_s$
    denote the MPO target blocks, not the chain tensor $C$.
    There are compression maps
    $V_s:\C^{D_s}\to\C^9$ and $W_s:\C^9\to\C^{D_s}$,
    where $D_1=2$ and $D_\tau=3$, such that
    $W_sV_s=\Id_{D_s}$ and $W_sV_t=0$ for $s\neq t$.
    For every word $w$ over the pair alphabet,
    \begin{align}
        W_sB^wV_s & =C_s^w.
        \label{eq:asymex_fib_reduction}
    \end{align}
    Every remainder word of length at least six vanishes.
    In fact, each remainder letter is zero, so the extension splits.
\end{theorem}

\begin{proof}
    \leanok
    \uses{thm:asymex_explicit_gauge, thm:asymex_golden_ring}
    The adapted basis has as its first columns the two fusion tensors,
    followed by a basis of the four-dimensional subspace annihilated
    by every source letter. Writing $Q$ for this basis matrix,
    the finite matrix calculation over $\Z[\sigma]$ gives
    \begin{align}
        Q^{-1}B^{x'x}Q
         & =B_1^{x'x}\oplus B_\tau^{x'x}\oplus0_4.
        \label{eq:asymex_fib_split_letters}
    \end{align}
    The corresponding column blocks of $Q$ and row blocks of $Q^{-1}$
    give $V_s$ and $W_s$. Their products give biorthogonality, and
    multiplication of~(\ref{eq:asymex_fib_split_letters}) along a word
    gives~(\ref{eq:asymex_fib_reduction}).
    Block diagonality gives block triangularity, the target diagonal
    blocks, and the vanishing remainder. The generic nilpotency
    bound of clause (vi) of Theorem~\ref{thm:asym_multiblock} is
    $|S|+z=6$, where $S$ is the two-element set of target blocks;
    the dimension count is clause (vii).
    The four zero directions are virtual bond directions, not
    physical configurations of the chain.
\end{proof}

\begin{theorem}[The fusion rule]
    \label{thm:asymex_fib_fusion}
    % Principal declaration: FibonacciCompression.fibonacci_fusion_rule.
    \lean{FibonacciCompression.fibStack_trace_evalWord,
        FibonacciCompression.fibonacci_fusion_rule,
        FibonacciCompression.fibStack_mpv}
    \leanok
    \uses{def:asymex_fib, def:mpo_family}
    For every nonempty word $w$ over the alphabet of label pairs,
    \begin{align}
        \tr(B^w) & =\tr(B_1^w)+\tr(B_\tau^w).
        \label{eq:asymex_fib_trace}
    \end{align}
    Thus at every positive system size $N$ the periodic operators
    satisfy~(\ref{eq:asymex_fib_fusion}), and the periodic vector
    of the stacked tensor is the sum of the periodic vectors of
    the two target blocks.
\end{theorem}

\begin{proof}
    \leanok
    \uses{thm:asymex_fib_compression, prop:asym_compression_trace}
    Apply Proposition~\ref{prop:asym_compression_trace} to
    Theorem~\ref{thm:asymex_fib_compression}. This gives
    (\ref{eq:asymex_fib_trace}).
    For output and input words $y,x$, the periodic product contracts
    as
    \begin{align}
        \bra{y}O_N(B_\tau)^2\ket{x}
         & =\sum_{m\in\{0,1\}^N}
        \tr(B_\tau^{y_0m_0}\cdots B_\tau^{y_{N-1}m_{N-1}})
        \tr(B_\tau^{m_0x_0}\cdots B_\tau^{m_{N-1}x_{N-1}})
        \nonumber                                     \\
         & =\tr(B^{y_0x_0}\cdots B^{y_{N-1}x_{N-1}}).
    \end{align}
    Substitution of~(\ref{eq:asymex_fib_trace}) proves the operator
    identity entry by entry. The periodic-vector statement is the
    same trace identity read over the pair alphabet.
\end{proof}

At length zero, the traces of the stacked tensor and the two targets
are $9$, $2$, and $3$, respectively. The omitted four-dimensional
zero block has empty trace four, although all its positive-length
traces vanish. This explains the positive-length restriction in the
fusion rule without changing the trace convention. The source asks for
the fusion rules at every length. At the empty chain the periodic
operators are the scalars $O_0(B_1)=2$ and $O_0(B_\tau)=3$, so
$O_0(B_1)^2=4\neq2$ and $O_0(B_1)O_0(B_\tau)=O_0(B_\tau)O_0(B_1)=6\neq3$:
the unit laws of Theorem~\ref{thm:asymex_fib_unit} fail there as well, and
the empty chain is excluded as a degenerate reading;
see~\cite{gap:bmwshv17_fibonacci_projector_positive_length}.

\begin{theorem}[The fusion tensors]
    \label{thm:asymex_fib_fusion_tensors}
    % Principal declaration: FibonacciCompression.fibStack_mul_fusionRight.
    \lean{FibonacciCompression.fibFusionRight, FibonacciCompression.fibFusionLeft,
        FibonacciCompression.fibonacci_right_eq, FibonacciCompression.fibonacci_left_eq,
        FibonacciCompression.fibStack_mul_fusionRight,
        FibonacciCompression.fusionLeft_mul_fibStack}
    \leanok
    \uses{def:asymex_fib, thm:asymex_fib_compression}
    The compression maps of the two target blocks are the fusion tensors
    of the Fibonacci algebra. The map $V_s$ is the column block of the
    adapted basis carrying $s$, and $W_s$ is the matching row block of
    its inverse. For every physical letter and $s\in\{1,\tau\}$,
    \begin{align}
        B^{x'x}V_s & =V_sC_s^{x'x}.
        \label{eq:asymex_fib_right_intertwiner}
    \end{align}
    \begin{align}
        W_sB^{x'x} & =C_s^{x'x}W_s.
        \label{eq:asymex_fib_left_intertwiner}
    \end{align}
\end{theorem}

\begin{proof}
    \leanok
    \uses{thm:asymex_explicit_gauge}
    Theorem~\ref{thm:asymex_explicit_gauge} identifies the maps with
    the row and column blocks of the adapted basis.
    Each intertwining relation is verified by multiplication over
    $\Z[\sigma]$ and mapped to complex matrices by
    Theorem~\ref{thm:asymex_golden_ring}.
    Equivalently, taking the corresponding column or row block
    of~(\ref{eq:asymex_fib_split_letters}) gives
    (\ref{eq:asymex_fib_right_intertwiner}) or
    (\ref{eq:asymex_fib_left_intertwiner}).
\end{proof}

\begin{center}
    % Formula: B^{x'x}V_s=V_sC_s^{x'x} and W_sB^{x'x}=C_s^{x'x}W_s of
    % Theorem~\ref{thm:asymex_fib_fusion_tensors}, for $s\in\{1,\tau\}$.
    % Source: Theorem~\ref{thm:asymex_fib_fusion_tensors}.
    % Ink-to-index map: the ring is the fusion tensor $V_s$ or $W_s$;
    % the square node is the stacked source $B$ or the target block
    % $C_s\in\{B_1,B_\tau\}$.
    % Contraction set: the single virtual bond joining the ring to the
    % square node in each pair; the physical pair leg $(x',x)$ remains
    % open.
    % Boundary signature: (virtual-west=1 | virtual-east=1 |
    % physical-up=1 | physical-down=0) for each pictured equality.
    \tenkzkernel
    \begin{tenkz}[rows={wire}, cols=2, west=open, east=open]
        \tn[label pos=s, ports={90:physical:$x'x$}]{B} & \tn[skin=ring]{V_s}
    \end{tenkz}
    \quad $=$ \quad
    \begin{tenkz}[rows={wire}, cols=2, west=open, east=open]
        \tn[skin=ring]{V_s} & \tn[label pos=s, ports={90:physical:$x'x$}]{C_s}
    \end{tenkz}
    \\
    \begin{tenkz}[rows={wire}, cols=2, west=open, east=open]
        \tn[skin=ring]{W_s} & \tn[label pos=s, ports={90:physical:$x'x$}]{B}
    \end{tenkz}
    \quad $=$ \quad
    \begin{tenkz}[rows={wire}, cols=2, west=open, east=open]
        \tn[label pos=s, ports={90:physical:$x'x$}]{C_s} & \tn[skin=ring]{W_s}
    \end{tenkz}
\end{center}
%
%% =========================================================
%% Subsection: the categorical anomaly of the Fibonacci symmetry as a
%% rank-one no-go. This file continues the Fibonacci section
%% (chapter/ch25_asymmetric_examples_fibonacci.tex,
%% Section~\ref{sec:asymex_fibonacci}) and is input right after it.
%% =========================================================

\subsection{The categorical anomaly}%
\label{sec:asymex_fib_anomaly}

The Fibonacci fusion ring $\Z[\tau]/(\tau^2-\tau-1)$ has no unital ring
homomorphism to $\Z$: the two possible images of $\tau$ over $\C$,
$\varphi=(1+\sqrt5)/2$ and $-1/\varphi$, are both irrational. This
arithmetic fact is the form in which the anomaly of the Fibonacci
fusion rules~\cite[Section~III.A.6]{Cirac2021Matrix} enters the
argument below. Chapter~\ref{ch:mpo_symmetry_basics} defines fusion algebras with a
unit, symmetric families of matrix product states, and their
nonnegative integer representations; by
Theorem~\ref{thm:mpo_symmetric_nimrep}, the multiplicities with which the
normal blocks of an invariant state reappear under the symmetry
operators form such a representation. In the
Fibonacci case the same source exhibits an invariant pair of states
carrying the regular action: two blocks $x_1,x_\tau$ with
$\tau\cdot x_1=x_\tau$ and $\tau\cdot x_\tau=x_1+x_\tau$.

Here the two blocks are concrete state tensors on the two-label
physical space of this section: the all-$\tau$ product tensor and a
bond-two tensor generating its image under $O_\tau$. The explicit
tensors are defined on the two-label physical alphabet rather than the
three-index alphabet of the source. Their unnormalized periodic
vectors carry the regular action. The operator $O_1$ is the
admissibility projector introduced earlier in this section, not the
identity on the ambient physical space. Its action as the identity on
these states is therefore part of the symmetry statement.

The no-go proved below concerns a single word-normal tensor $A$ of positive
bond dimension whose periodic vectors satisfy
$O_1\ket{V^{(N)}(A)}=\ket{V^{(N)}(A)}$ and
$O_\tau\ket{V^{(N)}(A)}=m\ket{V^{(N)}(A)}$ for a length-independent scalar
$m$ at every $N>0$. Then $m\in\N$ must satisfy $m^2=m+1$, which has no integer
solutions. In the Fibonacci case of
Corollary~\ref{cor:mpo_symmetric_no_go}, which specializes the statement
of~\cite{GarreRubio2023MPOSym} to one normal tensor with a
length-independent scalar action (without a proof for this tensor in the
source), this contradiction provides a non-invertible analogue of the
three-cocycle obstruction to an invariant single-block state, without
classifying multi-block families or excluding direct sums of the regular
representation.

Theorem~\ref{thm:asymex_fib_anomaly_nim_rep} records the explicit regular
action on two blocks, and Theorem~\ref{thm:asymex_fib_anomaly_no_go}
establishes the single-block obstruction under the hypotheses just stated.
Neither result classifies all invariant families with several blocks.

\begin{definition}[The Fibonacci fusion matrices]%
    \label{def:asymex_fib_anomaly_nim}
    \lean{FibonacciCompression.fibFusionMatrix, FibonacciCompression.fibNim,
        FibonacciCompression.fibBlock}
    \leanok
    \uses{def:asymex_fib}
    Index the two simple objects of Definition~\ref{def:asymex_fib} by
    $0$ and $1$, corresponding respectively to the trivial object $1$
    and the object $\tau$. The fusion matrix
    \begin{align}
        N_\tau & =\begin{pmatrix}0&1\\1&1\end{pmatrix}
        \label{eq:asymex_fib_anomaly_matrix}
    \end{align}
    represents multiplication by $\tau$ in the basis $1,\tau$ of the
    fusion ring. Together with $N_1=\Id_2$, it gives the regular
    representation by matrices with nonnegative integer entries. The
    structure constants are $N_{ab}^c=(N_a)_{bc}$ for
    $a,b,c\in\{1,\tau\}$, so that $N_1$ and $N_\tau$ are the fusion
    matrices of Definition~\ref{def:mposym_pf_dimension}. The corresponding
    matrix product operator tensors are the admissibility block $B_1$ and
    the block $B_\tau$ of Definition~\ref{def:asymex_fib}.
\end{definition}

\begin{theorem}[The Fibonacci relation for the fusion matrix]%
    \label{thm:asymex_fib_anomaly_matrix_sq}
    \lean{FibonacciCompression.fibFusionMatrix_sq}
    \leanok
    \uses{def:asymex_fib_anomaly_nim}
    $N_\tau^2=N_\tau+\Id_2$, as matrices with nonnegative integer entries.
\end{theorem}

\begin{proof}
    \leanok
    Direct multiplication gives
    $N_\tau^2=\begin{pmatrix}1&1\\1&2\end{pmatrix}=N_\tau+\Id_2$.
\end{proof}

\begin{theorem}[The fusion ring on the periodic operators]%
    \label{thm:asymex_fib_anomaly_fusion_algebra}
    \lean{FibonacciCompression.isMPOFusionAlgebra_fibBlock}
    \leanok
    \uses{def:asymex_fib_anomaly_nim}
    For every pair of labels $a,b\in\{1,\tau\}$ and every positive system
    size $N$, the periodic operators $O_a=O_a(N)=O_N(B_a)$ satisfy
    \begin{align}
        O_aO_b & =\sum_{c\in\{1,\tau\}}N_{ab}^cO_c,
        \label{eq:asymex_fib_anomaly_fusion_algebra}
    \end{align}
    that is, $O_1O_1=O_1$, $O_1O_\tau=O_\tau O_1=O_\tau$, and
    $O_\tau O_\tau=O_1+O_\tau$.
    \begin{center}
        % Formula/source: the (tau,tau) and (1,tau) cases of
        % (\ref{eq:asymex_fib_anomaly_fusion_algebra}), i.e.
        % O_tau O_tau = O_1 + O_tau (already (\ref{eq:asymex_fib_fusion}))
        % and O_1 O_tau = O_tau.
        % Ink-to-index map: each row of L tensors is the periodic operator
        % O_a(L), built from L copies of the block B_a; the top row's
        % physical legs i_1,...,i_L are the free outgoing labels, the
        % bottom row's physical legs j_1,...,j_L are the free incoming
        % labels, and the unlabelled physical legs shared between the two
        % rows are the L contracted intermediate labels of the
        % composition O_a(L)*O_b(L).
        % Contraction set: the L intermediate physical legs between the
        % two rows, and, independently on each row, the periodic virtual
        % bond of that row's block.
        % Boundary signature: (virtual-west=0 | virtual-east=0 |
        % physical-up=L | physical-down=L); no virtual legs remain open,
        % both rows being closed by their own periodic trace.
        \tenkzkernel
        \begin{tenkz}[rows={op,op}, cols=4, boundary=periodic]
            \tn[skin=mpo, ports={90:physical:$i_1$, 270:physical}]{B_\tau} &
            \tn[skin=mpo, ports={90:physical:$i_2$, 270:physical}]{B_\tau} &
            \tn[skin=dots, ports={180:virtual, 0:virtual}]{} &
            \tn[skin=mpo, ports={90:physical:$i_L$, 270:physical}]{B_\tau} \\
            \tn[skin=mpo, ports={90:physical, 270:physical:$j_1$}]{B_\tau} &
            \tn[skin=mpo, ports={90:physical, 270:physical:$j_2$}]{B_\tau} &
            \tn[skin=dots, ports={180:virtual, 0:virtual}]{} &
            \tn[skin=mpo, ports={90:physical, 270:physical:$j_L$}]{B_\tau}
        \end{tenkz}
        \qquad\text{and}\qquad
        \begin{tenkz}[rows={op,op}, cols=4, boundary=periodic]
            \tn[skin=mpo, ports={90:physical:$i_1$, 270:physical}]{B_1} &
            \tn[skin=mpo, ports={90:physical:$i_2$, 270:physical}]{B_1} &
            \tn[skin=dots, ports={180:virtual, 0:virtual}]{} &
            \tn[skin=mpo, ports={90:physical:$i_L$, 270:physical}]{B_1} \\
            \tn[skin=mpo, ports={90:physical, 270:physical:$j_1$}]{B_\tau} &
            \tn[skin=mpo, ports={90:physical, 270:physical:$j_2$}]{B_\tau} &
            \tn[skin=dots, ports={180:virtual, 0:virtual}]{} &
            \tn[skin=mpo, ports={90:physical, 270:physical:$j_L$}]{B_\tau}
        \end{tenkz}
    \end{center}
    In the diagrams the system size is denoted by $L$: the first
    contraction is $O_\tau(L)O_\tau(L)$ and the second is
    $O_1(L)O_\tau(L)$.
\end{theorem}

\begin{proof}
    \leanok
    \uses{thm:asymex_fib_fusion}
    The case $a=b=1$ is the admissibility projector composed with itself,
    the cases $\{a,b\}=\{1,\tau\}$ are the two unit laws obtained by
    block-triangular compression over $\Z[\sigma]$ and embedding into the
    complex matrices via Theorem~\ref{thm:asymex_golden_ring}, and
    Theorem~\ref{thm:asymex_fib_fusion} settles the case $a=b=\tau$.
    Together, these four identities determine all coefficients in
    Eq.~\ref{eq:asymex_fib_anomaly_fusion_algebra}.
\end{proof}

The pair of states is the all-$\tau$ product tensor and the chain
tensor of Definition~\ref{def:asymex_fib_states}, and the identities
$O_1\Psi_1=\Psi_1$, $O_\tau\Psi_1=\Psi_\tau$, and
$O_\tau\Psi_\tau=\Psi_1+\Psi_\tau$ are those of
Theorem~\ref{thm:asymex_fib_state_action}. The content beyond that
theorem is the remaining identity $O_1\Psi_\tau=\Psi_\tau$ and the
packaging of all four as the regular representation of the fusion
matrices.

\begin{theorem}[The regular representation on a pair of normal states]%
    \label{thm:asymex_fib_anomaly_nim_rep}
    \lean{FibonacciCompression.isMPOSymmetricFamily_fibNimTargets,
        FibonacciCompression.mpo_fibOne_mulVec_chain}
    \leanok
    \uses{def:asymex_fib_anomaly_nim, cor:asym_action_tensors}
    Set $s=\varphi^{-1/2}$, the scalar denoted by $\sigma$ in
    Definition~\ref{def:asymex_golden_ring}. On the physical alphabet
    $\{0,1\}$, with $0$ representing the trivial label and $1$
    representing $\tau$, define the bond-one tensor $A$ and bond-two
    tensor $C$ by
    \begin{align}
        A^0 & =(0),                                  & A^1                                       & =(1), \nonumber \\
        C^0 & =\begin{pmatrix}0&s\\0&0\end{pmatrix},
            & C^1                                    & =\begin{pmatrix}0&0\\1&-s^2\end{pmatrix}.
        \label{eq:asymex_fib_anomaly_states}
    \end{align}
    Both tensors are word-normal: their words of a fixed positive
    length span their full bond matrix algebra. For $A$ length one
    suffices, and for $C$ length three suffices.

    Define their unnormalized periodic vectors
    \begin{align}
        \Psi_1(N)    & =\ket{V^{(N)}(A)}
        =\ket{1,\ldots,1}, \nonumber     \\
        \Psi_\tau(N) & =\ket{V^{(N)}(C)}
        =\sum_{x\in\{0,1\}^N}
        \tr(C^{x_0}\cdots C^{x_{N-1}})\ket{x}.
        \label{eq:asymex_fib_anomaly_vectors}
    \end{align}
    At every positive system size $N$ and for all
    $a,r\in\{1,\tau\}$, these states satisfy
    \begin{align}
        O_a\Psi_r & =\sum_{t\in\{1,\tau\}}N_{ar}^t\Psi_t,
        \label{eq:asymex_fib_anomaly_nim_rep}
    \end{align}
    with dependence on $N$ suppressed. Explicitly,
    $O_\tau\Psi_1=\Psi_\tau$,
    $O_\tau\Psi_\tau=\Psi_1+\Psi_\tau$, and $O_1$ fixes both states.
    The multiplicities are the rows of the fusion matrix in
    Eq.~\ref{eq:asymex_fib_anomaly_matrix}.

    These identities use the tensor normalization in
    Eq.~\ref{eq:asymex_fib_anomaly_states}; at length $N=1$,
    $\Psi_1(1)=\ket{1}$ and $\Psi_\tau(1)=-s^2\ket{1}$.

    The following diagrams recall the operator fusion intertwiners of
    Theorem~\ref{thm:asymex_fib_fusion_tensors}. Here $B$ is the stacked
    operator square and $C_s$ denotes an operator target block, not the
    state tensor $C$ of Eq.~\ref{eq:asymex_fib_anomaly_states}.
    The state action uses the separate action tensors described in the
    proof.
    \begin{center}
        % Formula/source: the sitewise fusion tensors underlying
        % (\ref{eq:asymex_fib_anomaly_nim_rep}), specializing
        % Theorem~thm:asymex_fib_fusion_tensors to the slots s=1,tau; the
        % remainder of the compression of Theorem~thm:asymex_fib_compression
        % vanishes identically, so the intertwining identities below hold
        % exactly, with no residual term.
        % Ink-to-index map: B is the stacked source tensor of
        % Definition~def:asymex_fib, with physical leg i; V_s is the
        % column block of the adapted basis carrying the slot s, and W_s
        % the matching row block of its inverse; C_s is the block of the
        % target of label s.
        % Contraction set: only the wire joining B (or C_s) to V_s (or
        % W_s) is summed; no other index is contracted.
        % Boundary signature: (virtual-west=1 | virtual-east=1 |
        % physical-up=1 | physical-down=0) on every wire.
        % Clarification: these retained diagrams show operator fusion,
        % not the distinct compressions of B_1 acting on A, B_tau acting
        % on A, and B_tau acting on C used for the state action.
        \tenkzkernel
        \begin{tenkz}[rows={wire}, cols=2, west=open, east=open]
            \tn[label pos=s, ports={90:physical:$i$}]{B} & \tn[skin=ring]{V_s}
        \end{tenkz}
        \quad $=$ \quad
        \begin{tenkz}[rows={wire}, cols=2, west=open, east=open]
            \tn[skin=ring]{V_s} & \tn[label pos=s, ports={90:physical:$i$}]{C_s}
        \end{tenkz}
        \\
        \begin{tenkz}[rows={wire}, cols=2, west=open, east=open]
            \tn[skin=ring]{W_s} & \tn[label pos=s, ports={90:physical:$i$}]{B}
        \end{tenkz}
        \quad $=$ \quad
        \begin{tenkz}[rows={wire}, cols=2, west=open, east=open]
            \tn[label pos=s, ports={90:physical:$i$}]{C_s} & \tn[skin=ring]{W_s}
        \end{tenkz}
    \end{center}
\end{theorem}

\begin{proof}
    \leanok
    \uses{thm:asymex_fib_fusion_tensors}
    The state tensors are $A$ and $C$ in
    Eq.~\ref{eq:asymex_fib_anomaly_states}. The nonzero letter of $A$
    spans its one-dimensional matrix algebra. For $C$, the four words
    $010$, $011$, $101$, and $110$ span $\MN{2}$: each matrix unit is
    a linear combination of these words with coefficients in $\Z[s]$.

    For an operator tensor $T$ and a state tensor $X$, write
    $(T\cdot X)^i=\sum_jT^{ij}\otimes X^j$ for the action tensor of
    Corollary~\ref{cor:asym_action_tensors}. Simultaneous changes of bond
    coordinates over $\Z[s]$ give
    \begin{align}
        B_1\cdot A    & \sim A\oplus 0_1, \nonumber                   \\
        B_\tau\cdot A & \sim C\oplus 0_1, \nonumber                   \\
        B_\tau\cdot C & \sim A\oplus C\oplus 0_1\oplus 0_1\oplus 0_1.
        \label{eq:asymex_fib_anomaly_action_blocks}
    \end{align}
    Here $\sim$ denotes simultaneous similarity of all letters, and
    $0_1$ is the bond-one tensor whose letters vanish. The source bond
    dimensions are respectively two, three, and six. All three
    decompositions are block diagonal, with no remainder.

    Taking traces of nonempty words in
    Eq.~\ref{eq:asymex_fib_anomaly_action_blocks} gives
    $O_1\Psi_1=\Psi_1$, $O_\tau\Psi_1=\Psi_\tau$, and
    $O_\tau\Psi_\tau=\Psi_1+\Psi_\tau$. The remaining case follows from
    the unit law:
    \begin{align}
        O_1\Psi_\tau
         & =O_1O_\tau\Psi_1=O_\tau\Psi_1=\Psi_\tau.
    \end{align}
    This proves Eq.~\ref{eq:asymex_fib_anomaly_nim_rep}.
    The operator fusion intertwiners of
    Theorem~\ref{thm:asymex_fib_fusion_tensors} give the compatible
    operator identity $O_\tau^2=O_1+O_\tau$; the state identities above
    follow from the action decompositions in
    Eq.~\ref{eq:asymex_fib_anomaly_action_blocks}.
\end{proof}

The action above uses the admissibility projector and chain amplitudes on the
two-label physical space from earlier in this section, with unnormalized
periodic contractions; as noted after Theorem~\ref{thm:asymex_fib_state_action},
the identities do not extend to length zero.

By Theorem~\ref{thm:asymex_fib_anomaly_integral_action} of
Chapter~\ref{ch:mpo_symmetry_basics}, any length-independent scalar action on a
normal tensor takes nonnegative integer values, which drives the single-block
no-go below.

\begin{theorem}[A nonzero length-independent eigenvalue is impossible]%
    \label{thm:asymex_fib_anomaly_eigenvalue_zero}
    \lean{FibonacciCompression.fibTau_eigenvalue_eq_zero}
    \leanok
    \uses{def:asymex_fib, def:normal}
    Let $A$ be a word-normal tensor of physical dimension two and
    positive bond dimension, and set
    $\Psi_A(N)=\ket{V^{(N)}(A)}$. If a single scalar $c\in\C$ satisfies
    $O_\tau(N)\Psi_A(N)=c\Psi_A(N)$ for every $N\geq1$, then $c=0$.
    No unit-fixing hypothesis is required for this conclusion.
\end{theorem}

\begin{proof}
    \leanok
    \uses{thm:asymex_fib_anomaly_fusion_algebra, thm:mpo_symmetric_character,
        ex:mpo_symmetry_fibonacci}
    Suppose $c\neq0$ and suppress $N\geq1$ in the notation. The unit law
    $O_1O_\tau=O_\tau$ of
    Theorem~\ref{thm:asymex_fib_anomaly_fusion_algebra} gives
    \begin{align}
        c\,O_1\Psi_A
         & =O_1O_\tau\Psi_A=O_\tau\Psi_A=c\Psi_A.
    \end{align}
    Cancelling $c$ yields $O_1\Psi_A=\Psi_A$ for every $N\geq1$. Hence $A$
    is symmetric under the Fibonacci fusion algebra with eigenvalues
    $(c_1,c_\tau)=(1,c)$, and
    Theorem~\ref{thm:mpo_symmetric_character} makes $c\in\N$ with
    \begin{align}
        c^2 & =c_\tau c_\tau=c_1+c_\tau=1+c.
    \end{align}
    No nonnegative integer satisfies this equation
    (Example~\ref{ex:mpo_symmetry_fibonacci}), so $c=0$.
\end{proof}

\begin{theorem}[No normal matrix product state is Fibonacci symmetric]%
    \label{thm:asymex_fib_anomaly_no_go}
    \lean{FibonacciCompression.not_exists_normal_fibonacci_symmetric}
    \leanok
    \uses{def:asymex_fib, def:normal}
    There is no positive bond dimension $D$, word-normal tensor $A$ of
    physical dimension two and bond dimension $D$, and single scalar
    $c\in\C$, such that for every $N\geq1$,
    \begin{align}
        O_1(N)\Psi_A(N)    & =\Psi_A(N), \nonumber \\
        O_\tau(N)\Psi_A(N) & =c\Psi_A(N),
        \label{eq:asymex_fib_anomaly_symmetry}
    \end{align}
    where $\Psi_A(N)=\ket{V^{(N)}(A)}$.
    Thus no single normal tensor of positive bond dimension realizes
    a unital rank-one action of these periodic Fibonacci operators
    with one length-independent scalar. The obstruction is the
    absence of a nonnegative integer $m$ satisfying $m^2=m+1$.

    The explicit pair $A,C$ of
    Theorem~\ref{thm:asymex_fib_anomaly_nim_rep} shows that two normal
    blocks suffice. Among nonempty finite families of word-normal
    tensors of positive bond dimension, fixed by $O_1$ and carrying
    length-independent nonnegative integer action coefficients at
    every positive length, the minimal number of blocks is therefore
    two. This counts tensor blocks, not the dimension of their
    periodic-vector span at each short length.
\end{theorem}

\begin{proof}
    \leanok
    \uses{cor:mpo_symmetric_no_go, thm:asymex_fib_anomaly_fusion_algebra,
        ex:mpo_symmetry_fibonacci}
    Suppose $A$ and $c$ satisfy
    Eq.~\ref{eq:asymex_fib_anomaly_symmetry}. Then $A$ is symmetric under
    the Fibonacci fusion algebra of
    Theorem~\ref{thm:asymex_fib_anomaly_fusion_algebra} with eigenvalues
    $(c_1,c_\tau)=(1,c)$. A fusion character $m$ with values in $\N$ would
    satisfy $m_1=1$ and
    \begin{align}
        m_\tau^2 & =m_1+m_\tau=1+m_\tau,
    \end{align}
    which no nonnegative integer does
    (Example~\ref{ex:mpo_symmetry_fibonacci}).
    Corollary~\ref{cor:mpo_symmetric_no_go} therefore excludes $A$.
\end{proof}

The scalar hypothesis applies to the unnormalized periodic vectors of a single
word-normal tensor across all positive lengths; it does not exclude
eigenvectors at a fixed chain length, length-dependent eigenvalues, or
non-normal state tensors.
%
%% =========================================================
%% Physics example: the weighted Ising bond-object twist as an
%% instance of the multi-block asymmetric compression theorem.
%% Covers Notes/OpenProblemsTN/checks/asym_ising_action_data.md,
%% verified by checks/asym_ising_action_verify.py, and the P6
%% renormalization-fixed-point twist of
%% Notes/OpenProblemsTN/problems/p6_rfp_structure_constant_l_dependence.tex.
%% The structure coefficients built from the fusion rules were
%% moved here from a deleted machinery file.
%% =========================================================

\section{The Ising anyon chain}%
\label{sec:asymex_ising_tensors}

The Ising fusion category has three simple objects, the vacuum $1$, the
fermion $\psi$, and the non-Abelian anyon $\sigma$; its F-symbols define
the string-net operators of Bultinck et
al.~\cite[Appendix~D.2]{Bultinck2017Anyons}. Chains of interacting anyons were
introduced for Fibonacci anyons by Feiguin et
al.~\cite{Feiguin2007GoldenChain}, and the
algebra of topological defects of the Ising chain is described by Aasen,
Mong, and Fendley~\cite{Aasen2016Topological}. Matrix product operators
realize the corresponding topological symmetries of anyon chains, as
reviewed in Ref.~\cite[Section~III.A.6]{Cirac2021Matrix}. The matrix product
operator tensors for each simple object are taken in the F-symbol convention of
Ref.~\cite[Appendix~D.2]{Bultinck2017Anyons}; the sigma tensor is scaled
by $\sqrt2$ so that every entry lies in $\Z[\sqrt2]$ and all local
computations are exact. Weighted direct sums of these tensors and
the weighted Verlinde action on a fusion-tree product state address
the question of whether the structure coefficients of a renormalization fixed
point can depend on the chain length~\cite[Theorem~4.14]{Cirac2016MPDO_arXiv}.

\begin{definition}[The Ising string-net tensors of the source]
    \label{def:asymex_ising_gsymbol_tensor}
    \lean{IsingTwist.isingLoopFactor, IsingTwist.isingFSymbolZ,
        IsingTwist.isingFSymbol, IsingTwist.isingGSymbol,
        IsingTwist.isingBondLetter, IsingTwist.isingStringNetTensor}
    \leanok
    \uses{def:asymex_ising_tensors, def:mpo_tensor}
    The quantum dimensions are $d_1=d_\psi=1$ and $d_\sigma=\sqrt2$;
    write $N_{ab}^c$ for the Ising fusion multiplicities.
    Write $v_a=\sqrt{d_a}$ for the Ising labels, so $v_1=v_\psi=1$ and
    $v_\sigma=2^{1/4}$, and let $F^{abc}_{def}$ be the $F$-symbols
    of~\cite[Appendix~D.2]{Bultinck2017Anyons}, nonzero only when
    $N_{ab}^e$, $N_{cd}^e$, $N_{ad}^f$ and $N_{bc}^f$ are positive, with the
    values $\pm1/\sqrt2$ on $F^{\sigma\sigma\sigma}_{\sigma ef}$, $-1$ on
    $F^{\psi\sigma\psi}_{\sigma\sigma\sigma}$ and
    $F^{\sigma\psi\sigma}_{\psi\sigma\sigma}$, and $1$ otherwise. The
    $G$-symbols are $G^{abc}_{def}=v_e^{-1}v_f^{-1}F^{abc}_{def}$. For a
    label $f$, a basis element $(x,y)$ of $P_f$ in
    Definition~\ref{def:asymex_ising_tensors} is read as the pair
    $(u,l)=(y,x)$ of plaquette labels above and below the operator line; so
    $P_\psi$ has the bond pairs $[(\psi,1),(1,\psi),(\sigma,\sigma)]$ and
    $P_\sigma$ has $[(\sigma,1),(\sigma,\psi),(1,\sigma),(\psi,\sigma)]$. The operator tensor
    of~\cite[Section~6.2 and Appendix~D.1]{Bultinck2017Anyons} is the
    crossing of the operator line $f$ with an edge line $e$, with the
    plaquettes $b$ upper left, $a$ upper right, $c$ lower left and $d$
    lower right. On the fusion-tree labels $k=(b,e,a)$ and $h=(c,e',d)$ its
    entries are
    \begin{align}
        \bigl(S_f^{kh}\bigr)_{p,q}
         & =\delta_{ee'}\,\delta_{p,(b,c)}\,\delta_{q,(a,d)}\,
        G^{abc}_{def}\sqrt{v_av_bv_cv_d}.
        \label{eq:asymex_ising_gsymbol_tensor}
    \end{align}
\end{definition}

Definition~\ref{def:asymex_ising_gsymbol_tensor} gives the printed
$G$-symbol tensor with all quantum-dimension factors. The diagonal bond
similarity of Theorem~\ref{thm:asymex_ising_gsymbol_similarity} relates it
to the $F$-symbol blocks used here. The additional factor $\sqrt2$ on the
$\sigma$ block multiplies its periodic operator by $2^{N/2}$.

For every $N>0$, close the bond indices by the ordinary, unnormalized
trace of the ordered product of local matrices. The first physical
index is the output and the second is the input.

The vacuum and fermion kernels are identified in
Theorem~\ref{thm:asymex_ising_path_operators}; the fermion is a partial
isometry whose initial space is the admissible-path subspace
(Corollary~\ref{cor:asymex_ising_fermion_partial_isometry}). The fusion
identities with the $\sigma$ block are in
Theorems~\ref{thm:asymex_ising_sigma_fusion}
and~\ref{thm:asymex_ising_sigma_square}, with the scaling factors retained.
The word-span certificates are in Theorem~\ref{thm:asymex_ising_normal}.
For weighted direct sums, the compression and periodic trace formula
are given in Theorems~\ref{thm:asymex_ising_compression}
and~\ref{thm:asymex_ising_trace}; the weighted action on the state
family is in Theorem~\ref{thm:asymex_ising_action_trace}.

The quantum dimensions of $1$, $\psi$, and $\sigma$ are $1,1,\sqrt2$,
respectively; they are the Perron--Frobenius dimensions of the fusion rules
below (Example~\ref{ex:mposym_ising_dimension}). For every simple object
$a$, the vacuum satisfies $1\otimes a=a\otimes1=a$. The remaining fusion
rules are
\begin{align}
    \psi\otimes\psi     & =1,                        \\
    \sigma\otimes\psi   & =\psi\otimes\sigma=\sigma, \\
    \sigma\otimes\sigma & =1\oplus\psi.
\end{align}
With rows and columns ordered as $(1,\psi)$, the nontrivial F-symbols are
\begin{align}
    F^{\sigma\sigma\sigma}_{\sigma}
                                                    & =\frac{1}{\sqrt2}
    \begin{pmatrix}1&1\\1&-1\end{pmatrix},                              \\
    [F^{\psi\sigma\psi}_{\sigma}]_{\sigma}^{\sigma} & =-1,              \\
    [F^{\sigma\psi\sigma}_{\psi}]_{\sigma}^{\sigma} & =-1.
\end{align}
All other allowed components are one, and forbidden components vanish.

The physical site is $\C^{10}$, with orthonormal basis $\ket{h}$,
$h=0,\ldots,9$. A label $h$ is a locally admissible fusion vertex
$(\ell(h),\rho(h),r(h))$, meaning that $r(h)$ occurs in
$\ell(h)\otimes\rho(h)$. On $N$ sites the ambient Hilbert space is
$\mathcal H_N=(\C^{10})^{\otimes N}$.
For an MPO tensor $T$, the first physical index is the output index,
and its periodic contraction is
\begin{align}
    \bra{k}O_N(T)\ket{h}
     & =\tr\!\left(T^{k_0h_0}\cdots T^{k_{N-1}h_{N-1}}\right).
    \label{eq:asymex_ising_periodic}
\end{align}
Here $h$ and $k$ are words of length $N$, and $\ket{h}$ denotes the
corresponding tensor-product basis vector. Unless stated otherwise,
$N\geq1$, and site indices are cyclic.

The entries of the three tensors lie in the ring $\Z[\sqrt2]$, so their
local products and changes of auxiliary coordinates can be calculated in
this ring before embedding into $\C$. We fix the ring and its embedding
before writing down the tensors.

\begin{definition}[The ring of \texorpdfstring{$\sqrt2$}{sqrt(2)}-integers and its complex embedding]%
    \label{def:asymex_zsqrt2_ring}
    \lean{zsqrt2ToComplex, MPSTensor.complexOfZsqrt2, MPSTensor.mulZsqrt2Tensor,
        MPSTensor.actZsqrt2Tensor}
    \leanok
    The ring $\Z[\sqrt2]$ consists of the numbers $a+b\sqrt2$ with
    $a,b\in\Z$, with multiplication determined by $(\sqrt2)^2=2$.
    It embeds into $\C$ by sending $\sqrt2$ to the positive real square
    root of two. Applying this embedding entrywise gives a complex matrix
    from a matrix over $\Z[\sqrt2]$. For matrix product operator tensors
    $M$ and $N$ over $\Z[\sqrt2]$, their local product is
    $(M\mathbin{\cdot}N)^{ik}=\sum_jM^{ij}\otimes N^{jk}$.
    For a matrix product operator tensor $M$ and a state tensor $A$,
    the local action is $(M\cdot A)^i=\sum_jM^{ij}\otimes A^j$.
\end{definition}

\begin{theorem}[Exact arithmetic in \texorpdfstring{$\Z[\sqrt2]$}{Z[sqrt(2)]}]%
    \label{thm:asymex_zsqrt2_ring}
    % Principal declaration: MPSTensor.complexOfZsqrt2_mul.
    \lean{zsqrt2ToComplex_injective, zsqrt2ToComplex_ne_zero,
        MPSTensor.complexOfZsqrt2_mul, MPSTensor.mulTensor_complexOfRing,
        MPSTensor.actTensor_complexOfRing,
        MPSTensor.isNormal_of_single_eq_smul_zsqrt2}
    \leanok
    \uses{def:asymex_zsqrt2_ring}
    The embedding of $\Z[\sqrt2]$ into $\C$ is injective and sends every
    nonzero element to a nonzero complex number. The entrywise embedding
    of matrices is multiplicative and commutes with the local product
    and local action of tensors over $\Z[\sqrt2]$.
    If the one-site matrices of a tensor over $\Z[\sqrt2]$ include a
    nonzero $\Z[\sqrt2]$-multiple of every matrix unit, its complex
    embedding is injective at one site.
\end{theorem}

\begin{proof}
    \leanok
    If $a+b\sqrt2=a'+b'\sqrt2$ as real numbers, then
    $(b-b')\sqrt2=a'-a$. Irrationality of $\sqrt2$ gives $b=b'$,
    whence $a=a'$. Multiplicativity of the entrywise embedding follows
    because it is induced by a ring homomorphism. Expanding the finite
    sums and tensor products proves compatibility with local products
    and actions. Finally, dividing each designated one-site matrix by
    its nonzero scalar factor exhibits every matrix unit in the complex
    span, as in Theorem~\ref{thm:asymex_normal_units}.
\end{proof}

\begin{definition}[The three tensors \texorpdfstring{$A_1$, $A_\psi$, $A_\sigma$}{A1, Apsi, Asigma}]%
    \label{def:asymex_ising_tensors}
    \lean{IsingTwist.isingRho, IsingTwist.isingOneZ, IsingTwist.isingPsiZ,
        IsingTwist.isingSigmaZ, IsingTwist.isingOne, IsingTwist.isingPsi,
        IsingTwist.isingSigma}
    \leanok
    \uses{def:asymex_zsqrt2_ring, def:mpo_tensor}
    The ten fusion-tree labels and the functions $f$ and $\epsilon$
    used in the tensor formulas are
    \begin{center}
        \begin{tabular}{c@{\quad}ccc@{\quad}cc}
            $h$ & $\ell(h)$ & $\rho(h)$ & $r(h)$   & $f(h)$ & $\epsilon(h)$ \\[2pt]
            $0$ & $1$       & $1$       & $1$      & $3$    & $1$           \\
            $1$ & $1$       & $\psi$    & $\psi$   & $4$    & $1$           \\
            $2$ & $1$       & $\sigma$  & $\sigma$ & $5$    & $1$           \\
            $3$ & $\psi$    & $1$       & $\psi$   & $0$    & $1$           \\
            $4$ & $\psi$    & $\psi$    & $1$      & $1$    & $1$           \\
            $5$ & $\psi$    & $\sigma$  & $\sigma$ & $2$    & $1$           \\
            $6$ & $\sigma$  & $1$       & $\sigma$ & $6$    & $1$           \\
            $7$ & $\sigma$  & $\psi$    & $\sigma$ & $7$    & $-1$          \\
            $8$ & $\sigma$  & $\sigma$  & $1$      & $9$    & $1$           \\
            $9$ & $\sigma$  & $\sigma$  & $\psi$   & $8$    & $1$
        \end{tabular}
    \end{center}
    The auxiliary spaces $P_1$, $P_\psi$, and $P_\sigma$ have ordered bases
    $[(1,1),(\psi,\psi),(\sigma,\sigma)]$,
    $[(1,\psi),(\psi,1),(\sigma,\sigma)]$, and
    $[(1,\sigma),(\psi,\sigma),(\sigma,1),(\sigma,\psi)]$,
    respectively. Write $E_{ab}$ for the matrix units in these ordered
    bases, with indices starting at zero. When a matrix-unit subscript
    is an anyon label, its numerical index is given by the order
    $(1,\psi,\sigma)$.

    The complete tensors $A_1$ and $A_\psi$ are
    \begin{align}
        A_1^{kh}    & =\delta_{kh}E_{\ell(h),r(h)},
        \label{eq:asymex_ising_one_letters}                        \\
        A_\psi^{kh} & =\delta_{k,f(h)}\epsilon(h)E_{\ell(h),r(h)}.
        \label{eq:asymex_ising_psi_letters}
    \end{align}
    Thus $f=(0\,3)(1\,4)(2\,5)(8\,9)$ fixes $6$ and $7$, and
    $\epsilon(h)=-1$ only at $h=7$. In particular,
    $A_\psi^{7,7}=-E_{22}$ encodes
    $[F^{\psi\sigma\psi}_{\sigma}]_{\sigma}^{\sigma}=-1$.

    Put $\widehat A_\sigma=\sqrt2A_\sigma$. Its complete list of
    nonzero matrices is
    \begin{center}
        \begin{tabular}{c@{\quad}c@{\qquad}c@{\quad}c}
            $(k,h)$ & $\widehat A_\sigma^{kh}$ &
            $(k,h)$ & $\widehat A_\sigma^{kh}$                             \\[2pt]
            $(0,6)$ & $\sqrt2E_{22}$           & $(1,7)$ & $\sqrt2E_{23}$  \\
            $(2,8)$ & $E_{20}$                 & $(2,9)$ & $E_{21}$        \\
            $(3,6)$ & $\sqrt2E_{33}$           & $(4,7)$ & $\sqrt2E_{32}$  \\
            $(5,8)$ & $E_{30}$                 & $(5,9)$ & $-E_{31}$       \\
            $(6,0)$ & $\sqrt2E_{00}$           & $(6,3)$ & $\sqrt2E_{11}$  \\
            $(7,1)$ & $\sqrt2E_{01}$           & $(7,4)$ & $\sqrt2E_{10}$  \\
            $(8,2)$ & $\sqrt2E_{02}$           & $(8,5)$ & $\sqrt2E_{12}$  \\
            $(9,2)$ & $\sqrt2E_{03}$           & $(9,5)$ & $-\sqrt2E_{13}$
        \end{tabular}
    \end{center}
    All unlisted matrices vanish. The factors $\pm1$ in this table
    are the rescaled $\sigma\sigma\sigma$ F-symbols $\pm1/\sqrt2$.
    The entries of $A_1$, $A_\psi$, and $\widehat A_\sigma$ belong
    to the subring $\Z[\sqrt2]=\{a+b\sqrt2:a,b\in\Z\}$ of $\C$.
\end{definition}

\begin{theorem}[Sector vanishing of the three tensors]%
    \label{thm:asymex_ising_sector_vanishing}
    \lean{IsingTwist.isingOneZ_eq_zero_of_rho_ne,
        IsingTwist.isingPsiZ_eq_zero_of_rho_ne,
        IsingTwist.isingSigmaZ_eq_zero_of_rho_ne}
    \leanok
    \uses{def:asymex_ising_tensors}
    For each tensor $A \in \{A_1, A_\psi, \widehat A_\sigma\}$ and each pair
    of fusion-tree labels $(k,h)$, if $\rho(k) \ne \rho(h)$ then
    $A(k,h) = 0$.
\end{theorem}

\begin{proof}
    \leanok
    Every nonzero matrix in Definition~\ref{def:asymex_ising_tensors}
    has $\rho(k)=\rho(h)$, as inspection of the tables shows.
\end{proof}

\subsection{The periodic operators of the three tensors}

The vacuum and fermion tensors are evaluated on closed paths first,
where their periodic operators become projections, and the rescaled
sigma tensor is evaluated on a product state second.

\subsubsection{Closed paths and the vacuum and fermion operators}

Local admissibility does not impose compatibility between neighboring
vertices. A word $h$ is a closed path when
$r(h_v)=\ell(h_{v+1})$ at every site, including the last-to-first bond.
Define its closure indicator and sign by
\begin{align}
    q_N(h)    & =\prod_{v=0}^{N-1}\delta_{r(h_v),\ell(h_{v+1})},
    \label{eq:asymex_ising_closure}                              \\
    \eta_N(h) & =\prod_{v=0}^{N-1}\epsilon(h_v)
    =(-1)^{\#\{v:h_v=7\}}.
    \label{eq:asymex_ising_path_sign}
\end{align}
The map $f$ acts on words site by site.
The closed-path subspace is
$\mathcal K_N=\spn_{\C}\{\ket{h}:q_N(h)=1\}$.

\begin{lemma}[Cyclic trace of matrix units]%
    \label{lem:asymex_ising_unit_loop}
    \lean{Matrix.trace_ofFn_prod_smul_single}
    \leanok
    For $N\geq1$, scalars $c_v$, and matrix-unit indices $a_v,b_v$,
    \begin{align}
        \tr\!\left(\prod_{v=0}^{N-1}c_vE_{a_vb_v}\right)
         & =\left(\prod_{v=0}^{N-1}c_v\right)
        \prod_{v=0}^{N-1}\delta_{b_v,a_{v+1}},
        \label{eq:asymex_ising_unit_loop}
    \end{align}
    where $a_N=a_0$.
\end{lemma}

\begin{proof}
    \leanok
    Successive multiplication uses $E_{ab}E_{cd}=\delta_{bc}E_{ad}$.
    Taking the trace gives the final factor
    $\tr(E_{a_0b_{N-1}})=\delta_{b_{N-1},a_0}$.
\end{proof}

\begin{theorem}[Vacuum projection and fermion action on closed paths]%
    \label{thm:asymex_ising_path_operators}
    % Principal declarations: IsingTwist.mpo_isingOne_apply and
    % IsingTwist.mpo_isingPsi_apply (the two kernels).
    \lean{IsingTwist.mpo_isingOne_apply, IsingTwist.mpo_isingPsi_apply,
        IsingTwist.mpo_isingOne_conjTranspose,
        IsingTwist.mpo_isingPsi_conjTranspose,
        IsingTwist.isingOne_mul_isingOne, IsingTwist.isingPsi_mul_isingPsi,
        IsingTwist.isingOne_mul_isingPsi, IsingTwist.isingPsi_mul_isingOne,
        IsingTwist.isingLeft, IsingTwist.isingRight,
        IsingTwist.IsAdmissiblePath, IsingTwist.isingPsiFuse,
        IsingTwist.isingOne_eq_smul_single,
        IsingTwist.isingPsi_eq_smul_single,
        MPOTensor.mpo_apply_of_eq_smul_single,
        IsingTwist.isingGaugeOneOneZ, IsingTwist.isingGaugePsiPsiZ}
    \leanok
    \uses{def:asymex_ising_tensors, lem:asymex_ising_unit_loop}
    Set $Q_N=O_N(A_1)$ and $S_N=O_N(A_\psi)$. For $N\geq1$,
    \begin{align}
        \bra{k}Q_N\ket{h} & =\delta_{kh}q_N(h),
        \label{eq:asymex_ising_vacuum_kernel}                \\
        \bra{k}S_N\ket{h} & =\delta_{k,f(h)}\eta_N(h)q_N(h).
        \label{eq:asymex_ising_fermion_kernel}
    \end{align}
    The operator $Q_N$ is the orthogonal projection onto $\mathcal K_N$,
    and
    \begin{align}
        Q_N^2       & =Q_N=Q_N^\dagger, \\
        S_N^2       & =Q_N,             \\
        Q_NS_N      & =S_NQ_N=S_N,      \\
        S_N^\dagger & =S_N.
    \end{align}
\end{theorem}

\begin{proof}
    \leanok
    Substituting~(\ref{eq:asymex_ising_one_letters})
    and~(\ref{eq:asymex_ising_psi_letters})
    into~(\ref{eq:asymex_ising_unit_loop}) gives the two kernels.
    Since $q_N(h)\in\{0,1\}$, the vacuum kernel is the orthogonal
    projection onto the stated basis subspace.

    Let $j$ exchange $1$ and $\psi$ and fix $\sigma$. The label table gives
    $\ell(f(h))=j(\ell(h))$, $r(f(h))=j(r(h))$,
    $\rho(f(h))=\rho(h)$, and $\epsilon(f(h))=\epsilon(h)$.
    Since $j$ is bijective, closure is preserved:
    $q_N(f(h))=q_N(h)$. Also,
    $\eta_N(f(h))=\eta_N(h)$ and $f^2(h)=h$. Therefore
    \begin{align}
        S_N^2\ket{h}
         & =q_N(h)q_N(f(h))\eta_N(h)\eta_N(f(h))\ket{f^2(h)}
        =q_N(h)\ket{h}.
    \end{align}
    Closure invariance also gives $Q_NS_N=S_NQ_N=S_N$.
    The real signs and the involution imply
    \begin{align}
        \bra{k}S_N^\dagger\ket{h}
         & =\overline{\delta_{h,f(k)}\eta_N(k)q_N(k)}
        =\delta_{k,f(h)}\eta_N(h)q_N(h).
    \end{align}
    Thus $S_N^\dagger=S_N$.
\end{proof}

The kernels show
$\range S_N=\range Q_N=\mathcal K_N$ and
$\ker S_N=\ker Q_N=\mathcal K_N^\perp$. On $\mathcal K_N$, the operator
$S_N$ is a self-adjoint signed permutation involution.

\begin{corollary}[The fermion operator is a partial isometry]%
    \label{cor:asymex_ising_fermion_partial_isometry}
    \lean{IsingTwist.isingPsi_conjTranspose_mul_isingPsi,
        IsingTwist.mpo_isingOne_ne_one,
        IsingTwist.mpo_isingPsi_notMem_unitaryGroup}
    \leanok
    \uses{thm:asymex_ising_path_operators}
    For $N\geq1$, $S_N^\dagger S_N=Q_N$ and $Q_N\neq\Id$. Thus $S_N$ is
    a partial isometry with initial projection $Q_N$, and it is not
    unitary on $\mathcal H_N$.
\end{corollary}

\begin{proof}
    \leanok
    By Theorem~\ref{thm:asymex_ising_path_operators},
    $S_N^\dagger S_N=S_N^2=Q_N$. The constant word with label
    $1=(1,\psi,\psi)$ is not closed, since $r(1)=\psi$ and
    $\ell(1)=1$, so the vacuum kernel vanishes on its diagonal entry
    and $Q_N\neq\Id$. A unitary $S_N$ would give $Q_N=S_N^\dagger S_N=\Id$.
\end{proof}
At one site, the only closed labels are $0,3,6,7$. Using physical
matrix units in $\MN{10}$,
\begin{align}
    Q_1 & =E_{00}+E_{33}+E_{66}+E_{77}, \\
    S_1 & =E_{30}+E_{03}+E_{66}-E_{77}.
\end{align}
Thus the vacuum label acts as the identity on closed paths, not on
the entire physical space. These identities follow from the concrete
tensor letters rather than from the abstract fusion rules alone.

\subsubsection{The sigma normalization and a product-state action}

The rescaling $\widehat A_\sigma=\sqrt2A_\sigma$ occurs at every site.
By multilinearity of the periodic contraction,
\begin{align}
    O_N(\widehat A_\sigma) & =(\sqrt2)^N O_N(A_\sigma).
    \label{eq:asymex_ising_sigma_scale}
\end{align}

\begin{theorem}[Sigma action on the constant closed path]%
    \label{thm:asymex_ising_sigma_product}
    \lean{IsingTwist.mpo_isingSigma_mulVec_const_six,
        IsingTwist.mpo_isingSigma_normalized_mulVec_const_six}
    \leanok
    \uses{def:asymex_ising_tensors}
    For $N\geq1$,
    \begin{align}
        O_N(\widehat A_\sigma)\ket{6}^{\otimes N}
         & =(\sqrt2)^N
        \left(\ket{0}^{\otimes N}+\ket{3}^{\otimes N}\right),
        \label{eq:asymex_ising_sigma_product}        \\
        O_N(A_\sigma)\ket{6}^{\otimes N}
         & =\ket{0}^{\otimes N}+\ket{3}^{\otimes N}.
    \end{align}
\end{theorem}

\begin{proof}
    \leanok
    \uses{def:asymex_ising_tensors}
    At input label $6=(\sigma,1,\sigma)$ the only nonzero matrices are
    $\widehat A_\sigma^{0,6}=\sqrt2E_{22}$ and
    $\widehat A_\sigma^{3,6}=\sqrt2E_{33}$.
    Products containing both matrix units vanish because
    $E_{22}E_{33}=E_{33}E_{22}=0$.
    Each constant output word has trace $(\sqrt2)^N$.
    This proves~(\ref{eq:asymex_ising_sigma_product});
    dividing by the factor in~(\ref{eq:asymex_ising_sigma_scale})
    gives the unscaled formula.
\end{proof}

All three constant paths in this identity lie in $\mathcal K_N$.
The input is a normalized product state, while the unscaled output
has squared norm two. No sigma-square operator identity is needed
for this calculation.

Positive length matters here. At length zero, the empty matrix
product has trace equal to its auxiliary dimension. Thus
$O_0(A_1)=O_0(A_\psi)=3$ and
$O_0(\widehat A_\sigma)=4$, as operators on a one-dimensional
physical space. The projection and product-state identities above
do not extend to that case.

\subsubsection{The fusion algebra of the three periodic operators}

The operator blocks of the Fibonacci string-net satisfy the Fibonacci
fusion rules~\cite[Appendix~D.1]{Bultinck2017Anyons}, and the Ising
operators are built in the same way from the Ising
F-symbols~\cite[Appendix~D.2]{Bultinck2017Anyons}. For the Ising category
the source prints only the fusion rules themselves. The theorems below
prove the corresponding operator identities for the tensors of
Definition~\ref{def:asymex_ising_tensors}. These tensors use the F-symbols
directly, without the square roots of quantum dimensions that the source
divides them by; by Theorem~\ref{thm:asymex_ising_gsymbol_similarity} this
choice does not change the periodic
operators~\cite{gap:bmwshv17_ising_boundary_tensor_normalization}. The
same abstract algebra is realized by the topological defects of the
Ising lattice model~\cite{Aasen2016Topological}.

Each product is computed by one change of auxiliary coordinates of the
product tensor.

\begin{lemma}[Operator fusion from a change of auxiliary coordinates]%
    \label{lem:asymex_ising_fusion_similarity}
    \lean{MPSTensor.mpo_mul_eq_of_zsqrt2_conj, MPOTensor.mpo_eq_of_conj,
        MPOTensor.evalWord_eq_mul_mul_of_conj, MPOTensor.mpo_zero_of_pos,
        MPSTensor.padZsqrt2, MPSTensor.complexOfZsqrt2_padZsqrt2,
        MPSTensor.padZsqrt2_smul_single}
    \leanok
    \uses{def:asymex_zsqrt2_ring, def:mpdo_operator_product}
    Let $X$, $Y$, and $T$ be MPO tensors on the same physical space, with
    entries in $\Z[\sqrt2]$ and auxiliary dimensions $D_X$, $D_Y$, and
    $D_T$, and let $z\geq0$. Suppose that a nonzero $c\in\Z[\sqrt2]$ and
    matrices $G$ of size $(D_T+z)\times D_XD_Y$ and $H$ of size
    $D_XD_Y\times(D_T+z)$ with entries in $\Z[\sqrt2]$ satisfy
    $GH=c\Id$, $HG=c\Id$, and
    \begin{align}
        G\,(X\mathbin{\cdot}Y)^{kh}\,H
         & =c\,\bigl(T^{kh}\oplus0_z\bigr)
        \label{eq:asymex_ising_fusion_similarity}
    \end{align}
    for all physical labels $k,h$, where $0_z$ is the zero matrix of
    size $z$. Then $O_N(X)\,O_N(Y)=O_N(T)$ for every $N\geq1$.
\end{lemma}

\begin{proof}
    \leanok
    \uses{thm:mpdo_operator_product_mpo, thm:mpo_direct_sum_additive}
    The periodic operator of the product tensor $X\mathbin{\cdot}Y$ is
    $O_N(X)\,O_N(Y)$. Over $\C$, the matrices $G$ and $c^{-1}H$ are
    mutually inverse, so conjugating every matrix of $X\mathbin{\cdot}Y$
    by them leaves the traces of all words unchanged. By
    (\ref{eq:asymex_ising_fusion_similarity}) the conjugated tensor is
    the direct sum of $T$ and the zero tensor of auxiliary dimension $z$.
    Its periodic operator is $O_N(T)+O_N(0)$, and $O_N(0)=0$ for
    $N\geq1$ because every word of positive length contains a zero
    matrix.
\end{proof}

\begin{theorem}[Fusion with the sigma operator]%
    \label{thm:asymex_ising_sigma_fusion}
    % Principal declarations: the four products; the gauges and the
    % letter lemmas are the verification data of the proof.
    \lean{IsingTwist.isingOne_mul_isingSigma,
        IsingTwist.isingSigma_mul_isingOne,
        IsingTwist.isingPsi_mul_isingSigma,
        IsingTwist.isingSigma_mul_isingPsi,
        IsingTwist.isingGaugeOneSigmaZ, IsingTwist.isingGaugeSigmaOneZ,
        IsingTwist.isingGaugePsiSigmaZ, IsingTwist.isingGaugeSigmaPsiZ,
        IsingTwist.isingSignPsiSigma, IsingTwist.isingSignSigmaPsi,
        IsingTwist.isingSigmaZ_eq_smul_single,
        IsingTwist.isingConj_of_rho_eq,
        MPSTensor.mul_mulTensorR_mul_transpose_eq_list}
    \leanok
    \uses{def:asymex_ising_tensors, lem:asymex_ising_fusion_similarity}
    For $N\geq1$,
    \begin{align}
        O_N(A_1)\,O_N(\widehat A_\sigma)
         & =O_N(\widehat A_\sigma)\,O_N(A_1)
        =O_N(\widehat A_\sigma),                \\
        O_N(A_\psi)\,O_N(\widehat A_\sigma)
         & =O_N(\widehat A_\sigma)\,O_N(A_\psi)
        =O_N(\widehat A_\sigma).
    \end{align}
\end{theorem}

\begin{proof}
    \leanok
    \uses{thm:asymex_ising_sector_vanishing}
    Each of the four products is an instance of
    Lemma~\ref{lem:asymex_ising_fusion_similarity} with $c=1$, $z=8$,
    a signed permutation matrix $G$ of size $12$, and $H=G^{\mathsf T}$. The
    matrices of all three tensors are multiples of single matrix units,
    so each conjugated matrix of the product tensor is a short explicit
    sum. By Theorem~\ref{thm:asymex_ising_sector_vanishing}, both sides
    of~(\ref{eq:asymex_ising_fusion_similarity}) vanish unless the two
    physical labels have the same middle label $\rho$, and the remaining
    identities are checked in exact arithmetic over $\Z[\sqrt2]$.
\end{proof}

\begin{corollary}[Fusion with the unscaled sigma operator]%
    \label{cor:asymex_ising_sigma_fusion_unscaled}
    \lean{IsingTwist.isingOne_mul_isingSigma_normalized,
        IsingTwist.isingSigma_normalized_mul_isingOne,
        IsingTwist.isingPsi_mul_isingSigma_normalized,
        IsingTwist.isingSigma_normalized_mul_isingPsi}
    \leanok
    \uses{thm:asymex_ising_sigma_fusion, thm:mpdo_mpo_scalar_scaling}
    For $N\geq1$,
    \begin{align}
        O_N(A_1)\,O_N(A_\sigma)
         & =O_N(A_\sigma)\,O_N(A_1)
        =O_N(A_\sigma),                \\
        O_N(A_\psi)\,O_N(A_\sigma)
         & =O_N(A_\sigma)\,O_N(A_\psi)
        =O_N(A_\sigma).
    \end{align}
    For the tensors of Definition~\ref{def:asymex_ising_tensors}, these
    relations realize the fusion rules $1\times\sigma=\sigma\times1=\sigma$
    and $\psi\times\sigma=\sigma\times\psi=\sigma$
    of~\cite[Appendix~D.2]{Bultinck2017Anyons} as periodic operator
    identities.
\end{corollary}

\begin{proof}
    \leanok
    Both sides of each identity in
    Theorem~\ref{thm:asymex_ising_sigma_fusion} are linear in
    $O_N(\widehat A_\sigma)$. Dividing by the factor $(\sqrt2)^N$
    of~(\ref{eq:asymex_ising_sigma_scale}) gives the four identities.
\end{proof}

\begin{theorem}[The square of the sigma operator]%
    \label{thm:asymex_ising_sigma_square}
    % Principal declarations: IsingTwist.isingSigma_mul_isingSigma and
    % IsingTwist.isingSigma_normalized_mul_self.
    \lean{IsingTwist.isingSigma_mul_isingSigma,
        IsingTwist.isingSigma_normalized_mul_self,
        IsingTwist.isingGaugeSigmaSigmaZ, IsingTwist.isingOnePsiZ,
        IsingTwist.isingSigmaSigma_conj,
        IsingTwist.isingGaugeSigmaSigmaZ_mul_transpose,
        IsingTwist.isingGaugeSigmaSigmaZ_transpose_mul,
        IsingTwist.complexOfZsqrt2_isingOnePsiZ,
        MPSTensor.twoTermMatrix, MPSTensor.twoTermMatrix_mul_mul_transpose}
    \leanok
    \uses{def:asymex_ising_tensors, lem:asymex_ising_fusion_similarity}
    For $N\geq1$,
    \begin{align}
        O_N(\widehat A_\sigma)^2
         & =2^N\bigl(O_N(A_1)+O_N(A_\psi)\bigr),
        \label{eq:asymex_ising_sigma_square_scaled} \\
        O_N(A_\sigma)^2
         & =O_N(A_1)+O_N(A_\psi).
        \label{eq:asymex_ising_sigma_square}
    \end{align}
\end{theorem}

\begin{proof}
    \leanok
    \uses{thm:asymex_ising_sector_vanishing, thm:mpo_direct_sum_additive,
        thm:mpdo_mpo_scalar_scaling}
    Lemma~\ref{lem:asymex_ising_fusion_similarity} applies with
    $X=Y=\widehat A_\sigma$, $c=2$, $z=10$, and $H=G^{\mathsf T}$, where the
    matrix $G$ of size $16$ has fourteen rows of the form $\sqrt2e_i$ and
    the two rows $e_2\pm e_7$. These two rows carry the F-move of
    $F^{\sigma\sigma\sigma}_{\sigma}$, and $GG^{\mathsf T}=G^{\mathsf T}G=2\Id$. The target
    tensor $T$ is the direct sum of $2A_1$ and $2A_\psi$. As in the proof
    of Theorem~\ref{thm:asymex_ising_sigma_fusion}, the identity
    (\ref{eq:asymex_ising_fusion_similarity}) is checked in exact
    arithmetic within each sector of $\rho$. The periodic operator of
    $T$ is $2^N\bigl(O_N(A_1)+O_N(A_\psi)\bigr)$ by additivity and the
    scaling law, which gives~(\ref{eq:asymex_ising_sigma_square_scaled}).
    Dividing by the square of the factor
    in~(\ref{eq:asymex_ising_sigma_scale}) gives~(\ref{eq:asymex_ising_sigma_square}).
\end{proof}

\subsubsection{Bond similarity for the source tensors}



\begin{theorem}[The source's tensors are bond-similar to the formal tensors]
    \label{thm:asymex_ising_gsymbol_similarity}
    % Principal declaration: IsingTwist.mpo_isingStringNetTensor.
    \lean{IsingTwist.mpo_isingStringNetTensor,
        IsingTwist.isingGSymbol_tetrahedral,
        IsingTwist.isingBlock_apply_eq_fSymbol,
        IsingTwist.isingStringNetGauge,
        IsingTwist.isingStringNetGauge_ne_zero,
        IsingTwist.isingStringNetTensor_conj,
        IsingTwist.isMPOFusionAlgebra_isingStringNetTensor,
        IsingTwist.mpo_isingStringNetTensor_sigma}
    \leanok
    \uses{def:asymex_ising_gsymbol_tensor, ex:mpo_symmetry_ising,
        def:mpo_fusion_algebra}
    The $G$-symbols satisfy $G^{abc}_{def}=G^{fce}_{abd}$, and the tensors
    $A_1$, $A_\psi$ and $A_\sigma$ are $F$-symbol tensors:
    $(A_f^{kh})_{p,q}=\delta_{ee'}\delta_{p,(b,c)}\delta_{q,(a,d)}F^{fce}_{abd}$.
    Let $g_f$ be the diagonal matrix on $P_f$ with
    $g_f(u,l)=(v_u/v_l)^{1/2}$. Then
    \begin{align}
        g_f\,S_f^{kh}\,g_f^{-1} & =A_f^{kh}
        \label{eq:asymex_ising_gsymbol_conj}
    \end{align}
    for $f=1,\psi,\sigma$ and all labels $k,h$. Consequently
    $O_N(S_f)=O_N(A_f)$ at every length, in particular
    $O_N(S_\sigma)=2^{-N/2}O_N(\widehat A_\sigma)$, and the operators
    $O_N(S_1)$, $O_N(S_\psi)$, $O_N(S_\sigma)$ satisfy the Ising fusion rules
    at every positive length.
\end{theorem}

\begin{proof}
    \leanok
    \uses{thm:asymex_zsqrt2_ring, lem:asymex_ising_fusion_similarity,
        thm:mpdo_mpo_scalar_scaling}
    Write $v_a=t^{n_a}$ with $t=2^{1/4}$, $n_1=n_\psi=0$ and $n_\sigma=1$.
    Multiplied by $\sqrt2\,t^{n_b+n_d+n_e+n_f}$, the symmetry of the
    $G$-symbols becomes an identity between entries of $\Z[\sqrt2]$ and
    powers of $t^2=\sqrt2$, checked in exact arithmetic over all labels; so
    is the $F$-symbol form of the three tensors. By the symmetry,
    $G^{abc}_{def}=v_b^{-1}v_d^{-1}F^{fce}_{abd}$, so the nonzero entries of
    $S_f^{kh}$ are those of $A_f^{kh}$ multiplied by
    $\sqrt{v_av_bv_cv_d}/(v_bv_d)=g_f(a,d)/g_f(b,c)$, which
    is~(\ref{eq:asymex_ising_gsymbol_conj}). The similarity cancels in the
    trace defining the periodic operator, so the operators agree, the
    relation to $\widehat A_\sigma$ is the scaling law, and the fusion rules
    are those of Example~\ref{ex:mpo_symmetry_ising}.
\end{proof}

\subsection{One-site injectivity by exact arithmetic}%
\label{sec:asymex_ising_ring}

The embedding of $\Z[\sqrt2]$ into $\C$ from
Theorem~\ref{thm:asymex_zsqrt2_ring} preserves the matrix-unit spanning
criterion of Theorem~\ref{thm:asymex_normal_units}.
Here that criterion establishes one-site injectivity, not the
spectral-radius-one normalization required of a normal tensor (NT)
in the CPSV convention.

\begin{theorem}[One-site injectivity of the three tensors]%
    \label{thm:asymex_ising_normal}
    \lean{IsingTwist.isingOne_isNormal, IsingTwist.isingPsi_isNormal,
        IsingTwist.isingSigma_isNormal}
    \leanok
    \uses{def:asymex_ising_tensors, def:normal}
    Each of the three tensors is injective at one site: the ten nonzero
    local matrices of $A_1$ span $\MN{3}$, those of $A_\psi$ span
    $\MN{3}$, and the sixteen nonzero local matrices of $\sqrt2A_\sigma$
    span $\MN{4}$.
\end{theorem}

\begin{proof}
    \leanok
    \uses{thm:asymex_zsqrt2_ring}
    Every matrix unit of the relevant auxiliary matrix algebra occurs,
    up to a nonzero factor in $\Z[\sqrt2]$, among the local matrices.
    This follows directly from the tables.
    Theorem~\ref{thm:asymex_zsqrt2_ring} gives the spanning statement
    over $\C$.
\end{proof}

\subsection{The weighted direct sum and its local product}%
\label{sec:asymex_ising_bond_object}

We now apply asymmetric compression to products of these operators.
Consider
$\Theta=\lambda_1A_1\oplus\lambda_\psi A_\psi\oplus\lambda_\sigma A_\sigma$.
First retain the two summands with labels $1$ and $\psi$ and rescale
$(\lambda_1,\lambda_\psi)=(1/2,3/10)$ to the integer weights $5,3$.
Write $\Theta_2=5A_1\oplus3A_\psi$.

The local product $\Theta_2\star(\sqrt2A_\sigma)$ represents a product
of operators, as in Section~\ref{sec:asym_mpo_product}. Here $\star$
denotes the local MPO product of
Definition~\ref{def:asymex_zsqrt2_ring}, distinguished from the action
of an MPO on an MPS. Theorem~\ref{thm:asym_mpo_product} applies to
this product. Section~\ref{sec:asymex_ising_action} instead applies the
full weighted direct sum to a product state using
Corollary~\ref{cor:asym_action_tensors}.
Neither calculation alone establishes a positive vertical decomposition
of an MPDO renormalization fixed point.

\begin{definition}[The weighted direct sum and the product tensor]%
    \label{def:asymex_ising_bond_object}
    \lean{IsingTwist.thetaTwoZ, IsingTwist.thetaTwo, IsingTwist.isingBondObjectZ,
        IsingTwist.isingBondObject}
    \leanok
    \uses{def:asymex_ising_tensors, def:mpdo_operator_product}
    Set $\Theta_2=5A_1\oplus3A_\psi$, with auxiliary space
    $\C^6=P_1\oplus P_\psi$. The local product tensor
    $B=\Theta_2\star(\sqrt2A_\sigma)$ has a hundred-letter pair alphabet
    and auxiliary dimension $24=6\cdot4$. If $u=0,\ldots,5$ indexes
    $P_1\oplus P_\psi$ and $v=0,\ldots,3$ indexes $P_\sigma$, their
    tensor-product coordinate is $4u+v$. Explicitly,
    \begin{align}
        B^{(h,h')} & =\sum_{g=0}^{9}
        \Theta_2^{hg}\otimes(\sqrt2A_\sigma)^{gh'}.
        \label{eq:asymex_ising_product_letters}
    \end{align}
\end{definition}

In the periodic diagrams below, $L$ denotes the number of sites.

\begin{center}
    % Formula: O_L(Theta_2) O_L(sqrt2 A_sigma) = O_L(B), the stacked product
    %   of Definition~\ref{def:asymex_ising_bond_object} closed into full
    %   periodic operators (data file, section 1.2, and Theorem~\ref{thm:asymex_ising_trace}
    %   below).
    % Ink-to-index map: the upper row of operator nodes is Theta_2, the
    %   lower row is sqrt2 A_sigma; the physical legs h_1,...,h_L of the top
    %   row and g_1,...,g_L of the bottom row are open, and the shared
    %   virtual wire between the two rows at every site is the summed middle
    %   label of the stacked product. On the right, the single row of B
    %   carries both physical labels at each site as one pair letter.
    % Contraction set: the L vertical wires joining the two rows on the
    %   left, one per site; the horizontal virtual wires within each row,
    %   closed periodically.
    % Boundary signature (virtual-west, virtual-east, physical-up,
    %   physical-down): left, (0,0,L,L) with each three-site representative
    %   emitting (0,0,3,3); right, (0,0,L,0) with the representative emitting
    %   (0,0,3,0), the pair label absorbing both physical legs into one node.
    % The vertical contractions are physical input-output contractions.
    % With both physical ports shown, the right-hand picture has signature
    % (0,0,3,3); flattening each pair gives the equivalent (0,0,3,0) signature.
    \tenkzkernel
    \begin{tenkz}[rows={op,op}, cols=4, boundary=periodic]
        \tn[skin=mpo, ports={90:physical:$h_1$, 270:physical}]{\Theta_2} &
        \tn[skin=mpo, ports={90:physical:$h_2$, 270:physical}]{\Theta_2} &
        \tn[skin=dots, ports={180:virtual, 0:virtual}]{} &
        \tn[skin=mpo, ports={90:physical:$h_L$, 270:physical}]{\Theta_2} \\
        \tn[skin=mpo, ports={90:physical, 270:physical:$g_1$}]{\sqrt2A_\sigma} &
        \tn[skin=mpo, ports={90:physical, 270:physical:$g_2$}]{\sqrt2A_\sigma} &
        \tn[skin=dots, ports={180:virtual, 0:virtual}]{} &
        \tn[skin=mpo, ports={90:physical, 270:physical:$g_L$}]{\sqrt2A_\sigma}
    \end{tenkz}
    \quad $=$ \quad
    \begin{tenkz}[rows={op}, cols=4, boundary=periodic]
        \tn[skin=mpo, ports={90:physical:$h_1$, 270:physical:$g_1$}]{B} &
        \tn[skin=mpo, ports={90:physical:$h_2$, 270:physical:$g_2$}]{B} &
        \tn[skin=dots, ports={180:virtual, 0:virtual}]{} &
        \tn[skin=mpo, ports={90:physical:$h_L$, 270:physical:$g_L$}]{B}
    \end{tenkz}
\end{center}

\subsection{The gauge and the block decomposition}%
\label{sec:asymex_ising_gauge}

The change of auxiliary coordinates is a signed permutation of the
twenty-four coordinates of $B$. Its first four rows are the F-move
isometry for $1\otimes\sigma=\sigma$, the next four are the F-move
isometry for $\psi\otimes\sigma=\sigma$, and the remaining sixteen
form a coordinate complement. The single negative sign is the F-symbol
$[F^{\psi\sigma\psi}_\sigma]=-1$.

\begin{definition}[Signed permutation matrices and the Ising gauge]%
    \label{def:asymex_ising_gauge}
    \lean{MPSTensor.signedPermMatrix, IsingTwist.isingPerm, IsingTwist.isingSign,
        IsingTwist.isingGaugeZ, IsingTwist.isingGauge, IsingTwist.isingSlots,
        IsingTwist.isingWeights, IsingTwist.isingTargets, IsingTwist.isingBlockZ}
    \leanok
    \uses{def:asymex_ising_bond_object, def:asymex_explicit_gauge}
    For a map $\pi$ on an index set and a sign vector $s$, form the
    matrix with entry $s(x)$ in row $x$ and column $\pi(x)$ and zero
    elsewhere. When $\pi$ is a permutation, this is a signed permutation
    matrix. The Ising gauge is the signed permutation matrix given by the
    F-move rows described above, in the sense of
    Definition~\ref{def:asymex_explicit_gauge}.
    There are two nonzero blocks: one from $1\otimes\sigma=\sigma$,
    with weight $5$, and one from $\psi\otimes\sigma=\sigma$, with
    weight $3$. Both carry $\sqrt2A_\sigma$, of auxiliary dimension
    four. The zero block has dimension $16=24-2\cdot4$.
\end{definition}

\begin{theorem}[Orthogonality of the Ising gauge]%
    \label{thm:asymex_ising_gauge_unitary}
    \lean{MPSTensor.signedPermMatrix_mul_mul_transpose,
        IsingTwist.isingGaugeZ_mul_transpose, IsingTwist.isingGaugeZ_transpose_mul}
    \leanok
    \uses{def:asymex_ising_gauge}
    For a permutation $\pi$ and sign vector $s$, let
    $G_{x,\pi(x)}=s(x)$ and let all other entries of $G$ vanish.
    For every matrix $M$,
    \begin{align}
        (GMG^{\mathsf T})_{xy} & =s(x)\,M_{\pi(x)\pi(y)}\,s(y).
        \label{eq:asymex_ising_gauge_conj}
    \end{align}
    The Ising gauge $G$ satisfies $GG^{\mathsf T}=G^{\mathsf T}G=\Id$,
    so its transpose is its inverse.
\end{theorem}

\begin{proof}
    \leanok
    Identity~(\ref{eq:asymex_ising_gauge_conj}) follows by expanding the
    matrix product against the two Kronecker deltas of the signed
    permutation matrices. For the Ising gauge, both products with its
    transpose are the identity, by finite verification over the
    twenty-four rows.
\end{proof}

\begin{theorem}[The local block identity for the Ising product tensor]%
    \label{thm:asymex_ising_letter_identity}
    % Principal declaration: IsingTwist.isingGaugeZ_conj.
    \lean{IsingTwist.isingGaugeZ_conj, IsingTwist.isingBondObject_conjMatrix}
    \leanok
    \uses{def:asymex_ising_gauge}
    In the block coordinates of Definition~\ref{def:asymex_ising_gauge},
    conjugation by the Ising gauge gives, for every pair $(h,h')$,
    \begin{align}
        GB^{(h,h')}G^{-1}
         & =5(\sqrt2A_\sigma)^{(h,h')}
        \oplus3(\sqrt2A_\sigma)^{(h,h')}\oplus0_{16}.
        \label{eq:asymex_ising_local_blocks}
    \end{align}
\end{theorem}

\begin{proof}
    \leanok
    \uses{thm:asymex_ising_sector_vanishing, thm:asymex_ising_gauge_unitary}
    If $\rho(h)\neq\rho(h')$, every summand in
    (\ref{eq:asymex_ising_product_letters}) vanishes by
    Theorem~\ref{thm:asymex_ising_sector_vanishing}.
    The right-hand side of~(\ref{eq:asymex_ising_local_blocks}) also
    vanishes. Otherwise, the sum defining $B^{(h,h')}$ restricts to
    $\rho(g)=\rho(h)$. Formula~(\ref{eq:asymex_ising_gauge_conj})
    reduces the claim to the entries in
    Definitions~\ref{def:asymex_ising_tensors}
    and~\ref{def:asymex_ising_gauge}, checked separately in the three
    middle-label sectors.
\end{proof}

\subsection{Compression onto the \texorpdfstring{$\sigma$}{sigma} tensor}%
\label{sec:asymex_ising_compression}

\begin{theorem}[Compression of the Ising product tensor]%
    \label{thm:asymex_ising_compression}
    % Principal declaration: IsingTwist.isingCompression.
    \lean{IsingTwist.isingCompression,
        IsingTwist.isingBondObject_dim_eq, IsingTwist.isingBondObject_z_eq}
    \leanok
    \uses{thm:asymex_ising_letter_identity, thm:asym_multiblock}
    The tensor $B$ of Definition~\ref{def:asymex_ising_bond_object}
    has the block triangular form of Theorem~\ref{thm:asym_multiblock}
    with two weighted copies of $\sqrt2A_\sigma$, of weights $5$ and
    $3$, and a sixteen-dimensional zero block, $24=4+4+16$.
    Let $V_s$ include block $s$ in the original auxiliary coordinates,
    and let $W_s$ project onto it, for $s=0,1$.
    These compression pairs are mutually biorthogonal:
    $W_sV_t=\delta_{st}\Id_4$. For every pair of physical labels,
    \begin{align}
        W_sB^{(h,h')}V_s & =(\mu_s\sqrt2A_\sigma)^{(h,h')},
    \end{align}
    where $\mu_0=5$ and $\mu_1=3$.
\end{theorem}

\begin{proof}
    \leanok
    \uses{thm:asymex_explicit_gauge, thm:asymex_zsqrt2_ring}
    Theorem~\ref{thm:asymex_ising_letter_identity} gives the block-diagonal
    form of every conjugated local matrix over $\Z[\sqrt2]$.
    Theorem~\ref{thm:asymex_zsqrt2_ring} preserves this identity over
    $\C$. If $\iota_s$ and $p_s$ are the coordinate inclusion and
    projection for block $s$, set $V_s=G^{-1}\iota_s$ and $W_s=p_sG$.
    Then $W_sV_t=p_s\iota_t$, and
    (\ref{eq:asymex_ising_local_blocks}) gives the compression identity.
    Block diagonality gives triangularity as in
    Theorem~\ref{thm:asymex_explicit_gauge}, and clause~(vii) of
    Theorem~\ref{thm:asym_multiblock} gives the zero-block dimension
    $24-(4+4)=16$.
\end{proof}

\begin{theorem}[The extension splits]%
    \label{thm:asymex_ising_split}
    % Principal declaration: IsingTwist.isingBondObject_remainder.
    \lean{IsingTwist.isingBondObject_remainder}
    \leanok
    \uses{thm:asymex_ising_compression}
    The remainder of the compression in
    Theorem~\ref{thm:asymex_ising_compression} vanishes identically.
    Thus the extension splits, and each nonzero block has local
    intertwiners in both directions. For every pair $(h,h')$,
    \begin{align}
        B^{(h,h')}V_s & =V_s(\mu_s\sqrt2A_\sigma)^{(h,h')}, \\
        W_sB^{(h,h')} & =(\mu_s\sqrt2A_\sigma)^{(h,h')}W_s.
    \end{align}
\end{theorem}

\begin{proof}
    \leanok
    \uses{thm:asymex_ising_letter_identity, prop:asym_split_intertwiners}
    Every conjugated local matrix is block diagonal by
    Theorem~\ref{thm:asymex_ising_letter_identity}. Its off-diagonal
    blocks, which form the conjugated remainder, therefore vanish.
    Proposition~\ref{prop:asym_split_intertwiners} then gives the two
    intertwining identities.
\end{proof}

In the following diagram, $C_s=\mu_s\sqrt2A_\sigma$, and $h$ denotes
a single letter of the pair alphabet.

\begin{center}
    % Formula: B V_s = V_s C_s and W_s B = C_s W_s of
    %   Theorem~\ref{thm:asymex_ising_split}, with C_s = mu_s sqrt2 A_sigma.
    % Ink-to-index map: the ring node V_s (resp. W_s) is the compression map
    %   into (resp. out of) the bond space; the labelled node B carries the
    %   physical letter h and the open virtual legs are the bond spaces of B
    %   on the left and of the slot on the right.
    % Contraction set: the single virtual wire joining B and the ring node on
    %   each side; no physical contraction.
    % Boundary signature (virtual-west, virtual-east, physical-up,
    %   physical-down): each equality is (1,1,1,0).
    \tenkzkernel
    \begin{tenkz}[rows={wire}, cols=2, west=open, east=open]
        \tn[label pos=s, ports={90:physical:$h$}]{B} & \tn[skin=ring]{V_s}
    \end{tenkz}
    \quad $=$ \quad
    \begin{tenkz}[rows={wire}, cols=2, west=open, east=open]
        \tn[skin=ring]{V_s} & \tn[label pos=s, ports={90:physical:$h$}]{C_s}
    \end{tenkz}\\
    \begin{tenkz}[rows={wire}, cols=2, west=open, east=open]
        \tn[skin=ring]{W_s} & \tn[label pos=s, ports={90:physical:$h$}]{B}
    \end{tenkz}
    \quad $=$ \quad
    \begin{tenkz}[rows={wire}, cols=2, west=open, east=open]
        \tn[label pos=s, ports={90:physical:$h$}]{C_s} & \tn[skin=ring]{W_s}
    \end{tenkz}
\end{center}

\subsection{The length-dependent product coefficient}%
\label{sec:asymex_ising_trace}

For a tensor $T$ and a word $w=(i_1,\ldots,i_{|w|})$, write
$T^w=T^{i_1}\cdots T^{i_{|w|}}$.

\begin{theorem}[The trace identity for the Ising product tensor]%
    \label{thm:asymex_ising_trace}
    \lean{IsingTwist.isingBondObject_trace_evalWord, IsingTwist.isingTwist_mpo}
    \leanok
    \uses{thm:asymex_ising_compression, prop:asym_compression_trace,
        def:mpo_family}
    For every nonempty word $w$ over the alphabet of pair labels,
    \begin{align}
        \tr(B^w) & =(5^{|w|}+3^{|w|})\tr((\sqrt2A_\sigma)^w).
        \label{eq:asymex_ising_trace}
    \end{align}
    Equivalently, at every positive system size $N$,
    \begin{align}
        O_N(\Theta_2)O_N(\sqrt2A_\sigma)
         & =(5^N+3^N)O_N(\sqrt2A_\sigma).
    \end{align}
\end{theorem}

\begin{proof}
    \leanok
    \uses{thm:asymex_ising_compression, prop:asym_compression_trace}
    Proposition~\ref{prop:asym_compression_trace} applied to the
    compression in Theorem~\ref{thm:asymex_ising_compression} gives
    $\tr(B^w)$ as the sum of the two weighted block traces.
    Pulling each weight out of the matrix product contributes
    $5^{|w|}$ and $3^{|w|}$. The zero block contributes nothing because
    $w$ is nonempty. Reading~(\ref{eq:asymex_ising_trace}) at each pair
    of physical configurations gives the operator identity.
\end{proof}

\begin{remark}[Two fusion channels give the same output tensor]%
    \label{rem:asymex_ising_hidden_channel}
    The fusions $1\otimes\sigma=\sigma$ and
    $\psi\otimes\sigma=\sigma$ give two copies of the same output tensor
    with different weights. The coefficient in
    (\ref{eq:asymex_ising_trace}) is the power sum
    $5^N+3^N=\tr(\diag(5,3)^N)$.
    This is a length-dependent coefficient for the stated periodic
    product identity. Its power-sum form alone does not identify these
    tensors with a basis of normal tensors (BNT) in the positive vertical
    decomposition of an MPDO renormalization fixed point, nor prove that
    the length dependence cannot be removed under other choices of tensors.
\end{remark}

The two-copy space can be written as $\C^2\otimes P_\sigma$.
Let $\Pi_1$ and $\Pi_\psi$ be the coordinate projections on $\C^2$.
In the following diagram, the ring denotes the multiplicity operator
$\Theta=5\Pi_1+3\Pi_\psi$, not the direct-sum MPO tensor previously
denoted by $\Theta$. Each $\sqrt2A_\sigma$ node acts on $P_\sigma$.
On the combined bond space the ring acts as
$\Theta\otimes\Id_{P_\sigma}$, and a tensor letter acts as
$\Id_2\otimes(\sqrt2A_\sigma)^h$.
One weight insertion per site therefore gives
\begin{align}
    \tr\!\left(\prod_{v=0}^{N-1}
    (\Theta\otimes(\sqrt2A_\sigma)^{h_v})\right)
     & =\tr(\Theta^N)
    \tr\!\left(\prod_{v=0}^{N-1}(\sqrt2A_\sigma)^{h_v}\right).
\end{align}

\begin{center}
    % Formula: the physics reading of Remark~\ref{rem:asymex_ising_hidden_channel},
    %   the strand sqrt2 A_sigma with the weighted bond object
    %   Theta = 5 Pi_1 + 3 Pi_psi inserted at one bond.
    % Ink-to-index map: each labelled node is sqrt2 A_sigma, physical legs
    %   omitted; the ring node on the inner bond is the weighted sum of the
    %   two rank-one idempotents cut out by the two hidden bond objects.
    % Contraction set: the horizontal virtual wires joining consecutive
    %   sqrt2 A_sigma nodes, closed periodically through the ring node.
    % Boundary signature: (virtual-west=0, virtual-east=0, physical-up=3,
    %   physical-down=0) for this three-site representative of the periodic
    %   chain; the ellipsis carries the remaining sites.
    % Each bond carries C^2 tensor P_sigma; the ring acts on C^2 and
    % the sigma node on P_sigma. There is one weight insertion per site.
    \tenkzkernel
    \begin{tenkz}[cols=6, boundary=periodic, physical=up]
        \tn[label pos=135, ports={90:physical:$h_1$}]{\sqrt2A_\sigma} &
        \tn[skin=ring]{\Theta} &
        \tn[label pos=135, ports={90:physical:$h_2$}]{\sqrt2A_\sigma} &
        \tn[skin=ring]{\Theta} &
        \tn[skin=dots]{} & \tn[skin=ring]{\Theta}
    \end{tenkz}
    \qquad $\Theta=5\,\Pi_1+3\,\Pi_\psi$
\end{center}

\subsection{Coefficients from the Ising fusion rules}

The power sum $5^N+3^N$ is one value of a family of structure
coefficients built from the Ising fusion multiplicities and one
inserted weighted label, which we now define in general. The
definition and the theorems of this subsection write $L$ for the
chain length that the surrounding prose calls $N$.

\begin{definition}[Structure coefficients from the Ising fusion rules]
    \label{def:asymex_ising_coefficients}
    \lean{TNLean.Algebra.IsingCrossing.Label,
        TNLean.Algebra.IsingCrossing.fusionMultiplicity,
        TNLean.Algebra.IsingCrossing.coefficient}
    \leanok
    Let the labels be $1,\psi,\sigma$, and let $N_{ab}^{c}\in\N$
    be the multiplicity of $c$ in the Ising fusion rules. For each
    label $a$, these rules are
    \begin{align}
        1\times a          & =a\times 1=a,             \\
        \psi\times\psi     & =1,                       \\
        \psi\times\sigma   & =\sigma\times\psi=\sigma, \\
        \sigma\times\sigma & =1+\psi.
    \end{align}
    For arbitrary weights $w_a\in\C$ and $L\in\N$, define
    \begin{align}
        c_{a,b,c}^{(L)}
         & =\sum_r\sum_t w_r^L N_{ar}^{t}N_{tb}^{c}.
        \label{eq:asymex_ising_coeff}
    \end{align}
    Both sums run over $\{1,\psi,\sigma\}$.
    The definition includes $L=0$ and imposes no positivity or
    nonvanishing condition on the weights.
\end{definition}

\begin{theorem}[The three outputs for two sigma labels]
    \label{thm:asymex_ising_sigma}
    \lean{TNLean.Algebra.IsingCrossing.coefficient_sigma_sigma_one,
        TNLean.Algebra.IsingCrossing.coefficient_sigma_sigma_psi,
        TNLean.Algebra.IsingCrossing.coefficient_sigma_sigma_sigma,
        TNLean.Algebra.IsingCrossing.fusionMultiplicity_sigma_sigma_sigma}
    \leanok
    \uses{def:asymex_ising_coefficients}
    For every choice of complex weights and every $L\in\N$,
    \begin{align}
        c_{\sigma,\sigma,1}^{(L)}      & =w_1^L+w_\psi^L, \\
        c_{\sigma,\sigma,\psi}^{(L)}   & =w_1^L+w_\psi^L, \\
        c_{\sigma,\sigma,\sigma}^{(L)} & =2w_\sigma^L.
    \end{align}
    In contrast, the base fusion multiplicity is
    $N_{\sigma\sigma}^{\sigma}=0$.
\end{theorem}
\begin{proof}
    \leanok
    \uses{def:asymex_ising_coefficients}
    Expand both three-element sums in~(\ref{eq:asymex_ising_coeff})
    and substitute the fusion table. For output $1$ or $\psi$, the
    surviving inserted labels are $1$ and $\psi$. For output $\sigma$,
    the inserted label $\sigma$ contributes through both intermediate
    labels $1$ and $\psi$.
\end{proof}

\begin{theorem}[Associativity of the structure coefficients]
    \label{thm:asymex_ising_assoc}
    \lean{TNLean.Algebra.IsingCrossing.coefficient_assoc}
    \leanok
    \uses{def:asymex_ising_coefficients}
    For arbitrary complex weights, every $L\in\N$, and all labels
    $a,b,c,d$,
    \begin{align}
        \sum_e c_{a,b,e}^{(L)}c_{e,c,d}^{(L)}
         & =\sum_f c_{b,c,f}^{(L)}c_{a,f,d}^{(L)}.
    \end{align}
\end{theorem}
\begin{proof}
    \leanok
    \uses{def:asymex_ising_coefficients}
    For each of the $3^4$ choices of $a,b,c,d$, expanding the finite sums
    and substituting the fusion multiplicities reduces both sides to
    identical polynomials in $w_1^L,w_\psi^L,w_\sigma^L$.
\end{proof}

\subsection{Action on a fusion-tree product state}%
\label{sec:asymex_ising_action}

We next apply the full weighted direct sum
$\Theta_3=5A_1\oplus3A_\psi\oplus2A_\sigma$ to the constant closed path
with label $6=(\sigma,1,\sigma)$. This is an operator acting on a
state, rather than the operator product of
Section~\ref{sec:asymex_ising_bond_object}.
The elementary basis action determines the resulting vector;
Corollary~\ref{cor:asym_action_tensors} describes its local compression.

\begin{definition}[The full weighted direct sum and the fusion-tree product states]%
    \label{def:asymex_ising_action}
    \lean{IsingTwist.thetaThreeZ, IsingTwist.thetaThree, IsingTwist.productStateZ,
        IsingTwist.productState, IsingTwist.sectorState, IsingTwist.sectorAction}
    \leanok
    \uses{def:asymex_ising_tensors, def:mps_tensor, cor:asym_action_tensors}
    Set $\Theta_3=5A_1\oplus3A_\psi\oplus2A_\sigma$, with auxiliary space
    $\C^{10}=P_1\oplus P_\psi\oplus P_\sigma$.
    Its sigma summand is $2A_\sigma=\sqrt2\widehat A_\sigma$.
    For a fusion-tree label $h_0$, define the bond-dimension-one tensor
    $\psi_{h_0}^h=\delta_{h,h_0}$. Its periodic vector is
    $\ket{V^{(N)}(\psi_{h_0})}=\ket{h_0}^{\otimes N}$.
    In particular, $\psi_{\sigma1\sigma}=\psi_6$ generates
    $\ket{\sigma1\sigma}^{\otimes N}$, the fusion tree of the trivial
    anyon chain. We also write $\psi_{111}=\psi_0$ and
    $\psi_{\psi1\psi}=\psi_3$.
    The action tensor $\Theta_3\cdot\psi_{\sigma1\sigma}$ is the tensor
    of Corollary~\ref{cor:asym_action_tensors}, with one-site matrices
    $(\Theta_3\cdot\psi_{\sigma1\sigma})^h
        =\sum_{h'}\Theta_3^{hh'}\otimes\psi_{\sigma1\sigma}^{h'}$.
\end{definition}

Both $Q_N$ and $S_N$ fix $\ket{6}^{\otimes N}$: label $6$ is closed,
fixed by $f$, and has positive sign. The sigma contribution follows
from~(\ref{eq:asymex_ising_sigma_product}):
\begin{align}
    O_N(\sqrt2\widehat A_\sigma)\ket{6}^{\otimes N}
     & =2^N\left(\ket{0}^{\otimes N}+\ket{3}^{\otimes N}\right).
\end{align}
The factor $2^N$ comes from two factors of $\sqrt2$ at each site,
not from a length-independent normalization.

\begin{theorem}[One-site injectivity of the fusion-tree product-state tensors]%
    \label{thm:asymex_ising_action_normal}
    \lean{IsingTwist.productState_isNormal}
    \leanok
    \uses{def:asymex_ising_action, def:normal}
    Every fusion-tree product-state tensor $T$ has auxiliary dimension
    $D=1$ and satisfies $\spn_{\C}\{T(i) : i\} = \MN{1}$.
\end{theorem}

\begin{proof}
    \leanok
    The auxiliary space is one-dimensional, so its matrix algebra is
    $\C$. The local matrix indexed by $h_0$ is the scalar $1$.
\end{proof}

\begin{theorem}[The action tensor is diagonal]%
    \label{thm:asymex_ising_action_diagonal}
    \lean{IsingTwist.sectorActionZ, IsingTwist.sectorAction_eq}
    \leanok
    \uses{def:asymex_ising_action, cor:asym_action_tensors}
    Let $C=\Theta_3\cdot\psi_{\sigma1\sigma}$ be the action tensor.
    Its only nonzero one-site matrices are
    \begin{align}
        C^0 & =2E_{88},         \\
        C^3 & =2E_{99},         \\
        C^6 & =5E_{22}+3E_{55}.
        \label{eq:asymex_ising_action_six}
    \end{align}
    Here $E_{ij}$ denotes the matrix unit in $\MN{10}$, with
    indices starting at zero.
\end{theorem}

\begin{proof}
    \leanok
    The product-state letter is $\psi_6^h=\delta_{h,6}$, so
    $C^k=\Theta_3^{k,6}$. The tables in
    Definition~\ref{def:asymex_ising_tensors} give $E_{22}$ in each
    of the $A_1$ and $A_\psi$ summands at output $6$.
    Their coordinates in $P_1\oplus P_\psi\oplus P_\sigma$ are $2$ and $5$.
    In the sigma summand the nonzero entries occur at outputs $0,3$
    and auxiliary coordinates $8,9$. Multiplying
    $\sqrt2\widehat A_\sigma$ gives coefficient $2$ at both coordinates.
    All other entries vanish.
\end{proof}

\begin{theorem}[Compression of the weighted Verlinde action]%
    \label{thm:asymex_ising_action_compression}
    % Principal declaration: IsingTwist.sectorCompression.
    \lean{IsingTwist.sectorSlots, IsingTwist.sectorLabel, IsingTwist.sectorWeights,
        IsingTwist.sectorTargets, IsingTwist.sectorGauge, IsingTwist.sectorCompression,
        IsingTwist.sectorAction_z_eq, IsingTwist.sectorAction_dim_eq}
    \leanok
    \uses{thm:asymex_ising_action_diagonal, thm:asym_multiblock}
    The action tensor of Theorem~\ref{thm:asymex_ising_action_diagonal}
    has the block triangular form of Theorem~\ref{thm:asym_multiblock}
    with four weighted bond-dimension-one product-state tensors.
    Two generate $\ket{\sigma1\sigma}^{\otimes N}$, with weights
    $5$ and $3$, from the summands indexed by $1$ and $\psi$.
    The other two generate $\ket{111}^{\otimes N}$ and
    $\ket{\psi1\psi}^{\otimes N}$, each with weight $2$, from the
    summand indexed by $\sigma$.
    The zero block has dimension six, and $10=1+1+1+1+6$.
\end{theorem}

\begin{proof}
    \leanok
    \uses{thm:asymex_explicit_gauge}
    Let $H$ reorder the auxiliary coordinates as
    $(2,5,8,9,0,1,3,4,6,7)$. For every physical letter $h$,
    \begin{align}
        HC^hH^{-1}
         & =5\delta_{h,6}\oplus3\delta_{h,6}
        \oplus2\delta_{h,0}\oplus2\delta_{h,3}\oplus0_6.
        \label{eq:asymex_ising_action_blocks}
    \end{align}
    This follows by permuting the diagonal matrices of
    Theorem~\ref{thm:asymex_ising_action_diagonal}.
    Block diagonality gives triangularity and identifies the four
    target tensors as in Theorem~\ref{thm:asymex_explicit_gauge},
    coordinate projections and inclusions form the reduction pairs, and
    clause~(vii) of Theorem~\ref{thm:asym_multiblock} gives the
    zero-block dimension $10-4=6$.
\end{proof}

\begin{theorem}[The extension of the action tensor splits]%
    \label{thm:asymex_ising_action_split}
    \lean{IsingTwist.sectorAction_remainder}
    \leanok
    \uses{thm:asymex_ising_action_compression}
    The remainder $R$ of the compression of
    Theorem~\ref{thm:asymex_ising_action_compression} vanishes identically:
    \begin{align}
        R = 0.
        \notag
    \end{align}
\end{theorem}

\begin{proof}
    \leanok
    Formula~(\ref{eq:asymex_ising_action_blocks}) is block diagonal.
    Subtracting its diagonal blocks therefore gives zero, so the
    conjugated remainder and hence the remainder vanish.
\end{proof}

\begin{theorem}[The weighted Verlinde action]%
    \label{thm:asymex_ising_action_trace}
    \lean{IsingTwist.sectorAction_trace_evalWord, IsingTwist.thetaThree_mpo_sectorState}
    \leanok
    \uses{thm:asymex_ising_action_compression, prop:asym_compression_trace}
    For every nonempty word $w$ over the ten fusion-tree labels,
    \begin{align}
        \tr((\Theta_3\cdot\psi_{\sigma1\sigma})^w)
         & =(5^{|w|}+3^{|w|})\tr((\psi_{\sigma1\sigma})^w) \notag \\
         & \quad+2^{|w|}\left(\tr((\psi_{111})^w)
        +\tr((\psi_{\psi1\psi})^w)\right).
    \end{align}
    Equivalently, at every positive system size $N$,
    \begin{align}
        O_N(\Theta_3)\ket{\sigma1\sigma}^{\otimes N}
         & =(5^N+3^N)\ket{\sigma1\sigma}^{\otimes N} \notag \\
         & \quad+2^N\left(\ket{111}^{\otimes N}
        +\ket{\psi1\psi}^{\otimes N}\right).
        \label{eq:asymex_ising_verlinde}
    \end{align}
    The two contributions to $\ket{\sigma1\sigma}^{\otimes N}$ come
    from the summands indexed by $1$ and $\psi$, with different weights.
\end{theorem}

\begin{proof}
    \leanok
    \uses{thm:asymex_ising_action_compression, prop:asym_compression_trace}
    Proposition~\ref{prop:asym_compression_trace} applied to the
    compression in Theorem~\ref{thm:asymex_ising_action_compression}
    gives the sum of the four weighted block traces.
    For each target tensor $\mu\psi_{h_0}$,
    $\tr((\mu\psi_{h_0})^w)=\mu^{|w|}\tr(\psi_{h_0}^w)$.
    The weights are $5,3,2,2$ and the labels are $6,6,0,3$,
    which gives the trace formula. The periodic action identity follows
    by contracting the action tensor and summing its four target
    vectors, with the same factors $\mu^N$.
    This proves~(\ref{eq:asymex_ising_verlinde}) and agrees with the
    closed-path basis calculation.
\end{proof}

At length zero the action tensor has trace $10$, whereas its four
one-dimensional target blocks have total empty trace $4$. The difference
is the empty trace of the six-dimensional zero block. At positive length
that block contributes zero, which explains the nonempty-word
and positive-size hypotheses.

\subsection{Contraction of a matrix-unit factor}

The following identities contract a finite matrix-unit factor before any
fusion decomposition of the remaining matrices. They apply to every
physical letter; no selection of matrix entries is required.

\begin{theorem}[Reconstruction from matrix blocks]
    \label{thm:asymex_matrix_unit_register_reconstruction}
    \lean{Matrix.sum_single_kronecker_submatrix}
    \leanok
    Let $R$ be a commutative semiring, let $I$ be a finite index set,
    and let $K$ be an index set. For a matrix $Z$ indexed by
    $I\times K$, write $Z_{ab}$ for its block with entries
    ${(Z_{ab})}_{ij}=Z_{(a,i),(b,j)}$. If $E_{ab}$ denotes the matrix
    unit indexed by $I$, then
    \begin{align}
        \sum_{a,b\in I}E_{ab}\otimes Z_{ab} & =Z.
        \label{eq:asymex_matrix_unit_register_reconstruction}
    \end{align}
\end{theorem}

\begin{proof}\leanok
    At the entry indexed by $((c,i),(d,j))$, the matrix unit
    $E_{ab}$ vanishes unless $a=c$ and $b=d$. The sum therefore
    reduces to ${(Z_{cd})}_{ij}=Z_{(c,i),(d,j)}$.
\end{proof}

\begin{definition}[Reassociation of the two matrix-unit factors]
    \label{def:asymex_matrix_unit_register_reassociation}
    \lean{MPOTensor.registerFusionEquiv}
    \leanok
    For index sets $V,I,K$, let
    \begin{align}
        e:(V\times I)\times(K\times I)
                       & \longrightarrow (V\times(I\times K))\times I, \notag \\
        e((v,a),(i,b)) & =((v,(a,i)),b).
        \label{eq:asymex_matrix_unit_register_reassociation}
    \end{align}
    This is a bijection. For a matrix on the first index set, denote
    the matrix obtained by reindexing its rows and columns along $e$ by $\mathcal R$.
\end{definition}

\begin{theorem}[Contraction in a stacked product]
    \label{thm:asymex_matrix_unit_register_contraction}
    \lean{MPOTensor.registerFusion_contraction}
    \leanok
    \uses{def:asymex_matrix_unit_register_reassociation}
    Let $R$ be a commutative semiring, let $H,I$ be finite index
    sets, and let $V,K$ be index sets. For each $h\in H$, let $X_h$
    be a matrix indexed by $V$ and let $Z_h$ be a matrix indexed by
    $I\times K$. For every $c,d\in I$,
    \begin{align}
        \mathcal R\left(
        \sum_{h\in H}\sum_{a,b\in I}
            (X_h\otimes E_{ab})\otimes({(Z_h)}_{ab}\otimes E_{cd})
        \right)
         & =\left(\sum_{h\in H}X_h\otimes Z_h\right)\otimes E_{cd}.
        \label{eq:asymex_matrix_unit_register_contraction}
    \end{align}
\end{theorem}

\begin{proof}\leanok
    \uses{def:asymex_matrix_unit_register_reassociation}
    At the entry $(((v,(a,i)),c'),((w,(b,j)),d'))$ of the left-hand side,
    the intermediate matrix unit leaves precisely
    \begin{align}
        \sum_{h\in H}X_h(v,w)Z_h((a,i),(b,j))E_{cd}(c',d'),
        \notag
    \end{align}
    which is the corresponding entry of the right-hand side.
\end{proof}

For $|V|=|K|=16$ and $|I|=4$, the stacked bond has dimension $4096$,
whereas after this contraction the remaining matrix has dimension $1024$
and is tensored with a $4$-dimensional matrix unit; its fusion
decomposition is a separate step.

%
\subsection{The three-object weighted direct sum}

Retaining all three hidden bond objects gives four fusion channels,
since $\sigma\otimes\sigma=1\oplus\psi$. The following result concerns
the local product of the weighted direct sum with the $\sigma$ strand.
The full boundary tensors require the additional register
contraction described in~\cite{gap:tnlean_ising_three_object_twist_scope}.

\begin{definition}[The three-object product tensor]%
    \label{def:asymex_ising_three_product}
    \lean{IsingTwist.isingBondObjectThree, IsingTwist.thetaThree_eq_directSum,
        IsingTwist.isingThreeBlockDim, IsingTwist.isingThreeTargets}
    \leanok
    \uses{def:asymex_ising_action, def:asymex_ising_tensors,
        def:mpdo_operator_product}
    Put $\widehat A_\sigma=\sqrt2A_\sigma$,
    $\Theta_3=5A_1\oplus3A_\psi\oplus2A_\sigma$, and
    $B_3=\Theta_3\star\widehat A_\sigma$, of auxiliary dimension forty.
    Its four weighted targets are
    $5\widehat A_\sigma$, $3\widehat A_\sigma$,
    $2\sqrt2A_1$, and $2\sqrt2A_\psi$, of dimensions $4,4,3,3$.
\end{definition}

\begin{lemma}[Four-channel Ising product]%
    \label{thm:asymex_ising_three_trace}
    \lean{IsingTwist.thetaThree_mpo_mul_isingSigma,
        IsingTwist.isingBondObjectThree_trace_evalWord}
    \leanok
    \uses{def:asymex_ising_three_product, thm:asymex_ising_sigma_fusion,
        thm:asymex_ising_sigma_square}
    At every positive length $L$,
    \begin{align}
        O_L(\Theta_3)O_L(\widehat A_\sigma)
         & = (5^L+3^L)O_L(\widehat A_\sigma)
        +(2\sqrt2)^L\bigl(O_L(A_1)+O_L(A_\psi)\bigr).
        \label{eq:asymex_ising_three_product}
    \end{align}
    For every nonempty word $w$, the corresponding trace identity is
    \begin{align}
        \operatorname{tr}(B_3^w)
         & = (5^{|w|}+3^{|w|})\operatorname{tr}(\widehat A_\sigma^w)
        +(2\sqrt2)^{|w|}
        \bigl(\operatorname{tr}(A_1^w)+\operatorname{tr}(A_\psi^w)\bigr).
        \label{eq:asymex_ising_three_trace}
    \end{align}
\end{lemma}

\begin{proof}
    \leanok
    \uses{thm:asymex_ising_sigma_fusion, thm:asymex_ising_sigma_square}
    The periodic operator of a direct sum is the sum of the periodic
    operators of its summands. Applying the fusion rules
    $1\otimes\sigma=\sigma$, $\psi\otimes\sigma=\sigma$, and
    $\sigma\otimes\sigma=1\oplus\psi$ to those summands yields
    (\ref{eq:asymex_ising_three_product}). Every pair-alphabet word is
    the word of a pair of physical configurations, so taking that matrix
    element gives (\ref{eq:asymex_ising_three_trace}).
\end{proof}

\begin{lemma}[Compression of the three-object product]%
    \label{thm:asymex_ising_three_compression}
    \lean{IsingTwist.isingBondObjectThree_compression_exists,
        IsingTwist.isingThreeCompression, IsingTwist.isingBondObjectThree_dim_eq,
        IsingTwist.isingBondObjectThree_z_eq}
    \leanok
    \uses{def:asymex_ising_three_product, thm:asymex_ising_three_trace,
        thm:asymex_ising_normal, thm:asym_multiblock}
    The tensor $B_3$ admits a multi-block asymmetric compression onto the
    four weighted targets in Definition~\ref{def:asymex_ising_three_product}.
    The auxiliary dimensions satisfy $40=4+4+3+3+26$: in an
    upper-triangular presentation, twenty-six additional one-dimensional
    diagonal blocks vanish.
\end{lemma}

\begin{proof}
    \leanok
    \uses{thm:asymex_ising_normal, thm:asymex_ising_three_trace,
        thm:asym_multiblock}
    Each target is a nonzero scalar multiple of a normal tensor of
    positive auxiliary dimension; applying Theorem~\ref{thm:asym_multiblock}
    to (\ref{eq:asymex_ising_three_trace}) and subtracting the sum of the
    four target dimensions from forty leaves $z=26$ zero diagonal blocks.
\end{proof}

\begin{lemma}[A split four-channel Ising compression]%
    \label{thm:asymex_ising_three_split}
    \lean{IsingTwist.isingThreeSplitCompression,
        IsingTwist.isingBondObjectThree_split_remainder,
        IsingTwist.isingBondObjectThree_mul_right_eq_right_mul,
        IsingTwist.isingBondObjectThree_left_mul_eq_mul_left}
    \leanok
    \uses{def:asymex_ising_three_product}
    There is an explicit four-target compression of $B_3$ with zero remainder.
    Its left and right maps $W_s,V_s$ satisfy, for every letter $a$ and each
    of the four weighted targets $C_s$,
    \begin{align*}
        B_3^aV_s & = V_sC_s^a, \\
        W_sB_3^a & = C_s^aW_s.
    \end{align*}
    In its auxiliary coordinates the letters are the direct sum of
    $5\widehat A_\sigma$, $3\widehat A_\sigma$,
    $2\sqrt2 A_1$, $2\sqrt2 A_\psi$, and twenty-six zero matrices.
\end{lemma}

\begin{proof}
    \leanok
    \uses{thm:asymex_ising_compression, thm:asymex_ising_sigma_square,
        prop:asym_split_intertwiners}
    Grouping the first two summands of $\Theta_3$ writes the stacked
    product as the direct sum of the twenty-four-dimensional two-object
    product and $\sqrt2$ times the sixteen-dimensional $\sigma$ square.
    In the direct sum of their gauges, the first inverse is its transpose
    and the second inverse is half its transpose; their letter identities
    give the four stated blocks together with sixteen and ten zero
    coordinates. Rearranging the auxiliary coordinates yields the
    block-diagonal form, so the remainder vanishes and multiplication by
    the block inclusion and projection gives the two intertwining
    identities.
\end{proof}
%
